
\documentclass[11pt,a4paper]{report}

% =========================
% Encodage & langue
% =========================
\usepackage[T1]{fontenc}
\usepackage[utf8]{inputenc}
\usepackage[french]{babel}
% Désactive l'espacement automatique sur la ponctuation, sinon les labels avec ':' peuvent casser avec cleveref.
\frenchbsetup{AutoSpacePunctuation=false}
% Et on désactive explicitement les shorthands sur la ponctuation concernée.
\AtBeginDocument{\frenchbsetup{AutoSpacePunctuation=false}\shorthandoff{:;!?}}
\usepackage{lmodern}

% =========================
% Packages requis (liste imposée)
% =========================
\usepackage{geometry}
\usepackage{microtype}
\usepackage{amsmath, amssymb, amsthm}
\usepackage{mathtools}
\usepackage{bm}
\usepackage{bbm}
\usepackage{graphicx}
\usepackage{booktabs}
\usepackage{longtable}
\usepackage{enumitem}
\usepackage{listings}
\usepackage{xcolor}
\usepackage{tikz}
\usepackage{pgfplots}
\usepackage{hyperref}
\usepackage{cleveref}

\pgfplotsset{compat=1.18}
\geometry{margin=2.5cm}

% =========================
% Hyperliens
% =========================
\hypersetup{
  colorlinks=true,
  linkcolor=blue,
  citecolor=blue,
  urlcolor=blue,
  pdftitle={Cours de finance quantitative (M1) --- version étendue},
  pdfauthor={},
  pdfsubject={Pricing, volatilité, calibration, couverture, taux, allocation},
  pdfkeywords={Black-Scholes, Ito, Heston, SABR, volatilite, hedging, RL, yield curve, Markowitz}
}

% =========================
% Macros mathématiques propres
% =========================
\newcommand{\dd}{\mathrm{d}}
\newcommand{\ee}{\mathrm{e}}
\newcommand{\ii}{\mathrm{i}}
\newcommand{\E}{\mathbb{E}}
\newcommand{\Var}{\mathrm{Var}}
\newcommand{\Cov}{\mathrm{Cov}}
\newcommand{\Prob}{\mathbb{P}}
% \P est défini en LaTeX (symbole \textparagraph). On le libère pour l'utiliser comme mesure de probabilité.
\let\P\relax
\newcommand{\P}{\mathbb{P}}
\newcommand{\Q}{\mathbb{Q}}
\newcommand{\R}{\mathbb{R}}
\newcommand{\N}{\mathbb{N}}
\newcommand{\1}{\mathbbm{1}}
\newcommand{\Normal}{\mathcal{N}}
\newcommand{\F}{\mathcal{F}}
\newcommand{\G}{\mathcal{G}}

\newcommand{\abs}[1]{\left|#1\right|}
\newcommand{\norm}[1]{\left\|#1\right\|}
\newcommand{\paren}[1]{\left(#1\right)}
\newcommand{\bracket}[1]{\left[#1\right]}
\newcommand{\set}[1]{\left\{#1\right\}}
\newcommand{\ip}[2]{\left\langle #1,#2\right\rangle}

\newcommand{\pd}[2]{\frac{\partial #1}{\partial #2}}
\newcommand{\pdd}[2]{\frac{\partial^2 #1}{\partial #2^2}}

\newcommand{\BS}{\mathrm{BS}}
\newcommand{\vega}{\mathrm{Vega}}
\newcommand{\DeltaH}{\Delta}
\newcommand{\PnL}{\mathrm{PnL}}
\newcommand{\CVaR}{\mathrm{CVaR}}
\newcommand{\VaR}{\mathrm{VaR}}

% =========================
% Environnements (définitions, théorèmes, etc.)
% =========================
\theoremstyle{definition}
\newtheorem{definition}{Définition}[chapter]
\newtheorem{example}{Exemple}[chapter]

\theoremstyle{plain}
\newtheorem{proposition}{Proposition}[chapter]
\newtheorem{theorem}{Théorème}[chapter]
\newtheorem{lemma}{Lemme}[chapter]
\newtheorem{corollary}{Corollaire}[chapter]

\theoremstyle{remark}
\newtheorem{remark}{Remarque}[chapter]

% =========================
% Encadré "Ce que le jury doit retenir"
% =========================
\newenvironment{jurybox}{%
  \begin{center}
  \setlength{\tabcolsep}{6pt}%
  \renewcommand{\arraystretch}{1.1}%
  \begin{tabular}{|p{0.94\textwidth}|}
  \hline
  \small\textbf{Ce que le jury doit retenir.}\par\medskip
}{%
  \\ \hline
  \end{tabular}
  \end{center}
}

% =========================
% Listings (pour pseudo-code/algorithmes, pas du code applicatif)
% =========================
\lstdefinestyle{algo}{
  basicstyle=\ttfamily\small,
  columns=fullflexible,
  frame=single,
  framerule=0.2pt,
  rulecolor=\color{black},
  tabsize=2,
  numbers=left,
  numberstyle=\tiny\color{gray},
  xleftmargin=2em,
  framexleftmargin=1.5em,
  breaklines=true,
  showstringspaces=false
}
\lstset{style=algo}

% =========================
% TikZ
% =========================
\usetikzlibrary{arrows.meta, positioning, shapes.geometric, calc}

\begin{document}

% =========================
% Page de titre
% =========================
\begin{titlepage}
  \centering
  \vspace*{2cm}
  {\LARGE\bfseries Cours de finance quantitative (M1)\par}
  \vspace{0.5cm}
  {\Large Pricing, volatilité, calibration, couverture, taux et allocation\par}
  \vspace{0.25cm}
  {\large Version étendue (objectif : \textasciitilde 120 pages)\par}
  \vspace{0.2cm}
  {\large Document pédagogique ``from ground up''\par}
  \vfill
  {\large \today\par}
\end{titlepage}

% =========================
% Résumé
% =========================
\pagenumbering{roman}
\chapter*{Résumé}
Ce document est un cours complet (niveau Master 1 finance) reliant, de manière progressive,
les briques théoriques nécessaires à la modélisation des marchés (processus stochastiques),
au pricing et à la couverture d'options (Black--Scholes, martingales, mesure risque-neutre),
à la construction et la calibration d'une surface de volatilité implicite (Heston, sauts, SABR, volatilité ``rough''),
ainsi qu'à la prise de décision en gestion (courbe des taux, optimisation de portefeuille de type Markowitz, risk parity,
réduction dimensionnelle par ACP).

Le fil conducteur est systématiquement le même :
\begin{enumerate}[label=\arabic*)]
  \item \textbf{Pourquoi cette notion existe en finance} (problème économique, instrument, risque).
  \item \textbf{Modélisation mathématique} (définitions, hypothèses, limites).
  \item \textbf{Dérivations détaillées} (pas de ``magie'') pour relier le modèle à une formule exploitable.
  \item \textbf{Numérique et pratique} (stabilité, calibration, erreurs, arbitrage).
\end{enumerate}
Chaque chapitre contient des définitions formelles, des explications intuitives, des exemples chiffrés,
et un encadré ``Ce que le jury doit retenir''.

\tableofcontents
\listoffigures
\clearpage
\pagenumbering{arabic}

\chapter{Architecture du cours \& notations}
\label{chap:archi_notations}

\section{Le fil conducteur : de la donnée de marché à la décision}
\label{sec:fil_conducteur}

En finance de marché, on peut schématiser une grande partie des problèmes quantitatifs en une chaîne logique :
\begin{enumerate}[label=\alph*)]
  \item \textbf{Observation} : on collecte des prix (actions, taux, options) et on les convertit en quantités plus stables (rendements, volatilité implicite, facteurs de courbe).
  \item \textbf{Modélisation} : on postule une dynamique aléatoire pour les sous-jacents (par exemple une diffusion, ou une diffusion avec sauts), et on choisit un cadre de valorisation (absence d'arbitrage).
  \item \textbf{Pricing \& couverture} : on calcule des prix théoriques d'options, et des sensibilités (\emph{greeks}) pour gérer le risque.
  \item \textbf{Calibration} : on ajuste les paramètres du modèle pour qu'il reproduise les prix observés, typiquement via la surface de volatilité implicite.
  \item \textbf{Décision} : on transforme ces informations (prix, risques, courbe des taux) en décisions de trading, de couverture ou d'allocation.
\end{enumerate}

\begin{figure}[h!]
\centering
\begin{tikzpicture}[
  node distance=10mm,
  box/.style={draw, rounded corners, align=center, minimum width=3.7cm, minimum height=1.0cm},
  arrow/.style={-Latex, thick}
]
\node[box] (md) {Données\\de marché};
\node[box, right=of md] (pricing) {Modèles\\stochastiques\\\& pricing};
\node[box, right=of pricing] (calib) {Calibration\\(surface IV)};
\node[box, below=of pricing] (hedge) {Couverture\\\& gestion\\du risque};
\node[box, right=of calib] (decision) {Décision\\(trading /\\allocation)};

\draw[arrow] (md) -- (pricing);
\draw[arrow] (pricing) -- (calib);
\draw[arrow] (calib) -- (decision);
\draw[arrow] (pricing) -- (hedge);
\draw[arrow] (hedge) -| (decision);
\end{tikzpicture}
\caption{Architecture conceptuelle : comment les briques théoriques s'enchaînent.}
\label{fig:pipeline_conceptuel}
\end{figure}

Ce cours suit exactement ce flux : \cref{fig:pipeline_conceptuel} n'est pas un ``diagramme joli'' ; c'est un
\emph{plan de bataille} pour comprendre pourquoi on introduit Brownien, Itô, la mesure risque-neutre, les modèles de volatilité,
puis des méthodes de couverture et enfin l'allocation.

\section{Notations standard (taux, spot, forward, discount)}
\label{sec:notations}

On travaille sur un horizon $[0,T]$. Les notations ci-dessous seront utilisées de façon cohérente tout au long du document.

\begin{longtable}{@{}llp{8.5cm}@{}}
\toprule
Symbole & Nom & Commentaire \\
\midrule
$S_t$ & \emph{spot} & Prix du sous-jacent (action/indice) au temps $t$. \\
$K$ & strike & Prix d'exercice d'une option. \\
$T$ & maturité & Date d'échéance d'un produit dérivé. \\
$r_t$ & taux court & Taux instantané (modélisation possible, ou approximé constant). \\
$r$ & taux constant & Hypothèse simplificatrice fréquente pour dériver Black--Scholes. \\
$q$ & taux de dividende & Dividend yield continu (ou taux de repo) pour un indice. \\
$B_t$ & compte sans risque & $B_t=\exp(\int_0^t r_u\,\dd u)$ si taux continu. \\
$P(0,T)$ & discount factor & Valeur aujourd'hui d'1 euro certain à $T$. \\
$F_{0,T}$ & forward & Prix forward à $0$ pour livraison à $T$ (sous hypothèses). \\
$\sigma$ & volatilité & Écart-type instantané (Black--Scholes) ou paramètre de diffusion. \\
$\sigma_{\mathrm{imp}}(K,T)$ & vol implicite & Volatilité Black--Scholes qui reproduit un prix de marché. \\
\bottomrule
\end{longtable}

\section{Probabilités : espace, filtration, martingales (sans measure theory lourde)}
\label{sec:proba_filtration}

\begin{definition}[Cadre probabiliste]
On note $(\Omega,\F,(\F_t)_{t\in[0,T]},\Prob)$ un espace probabilisé filtré :
\begin{itemize}
  \item $\Omega$ est l'ensemble des scénarios possibles ;
  \item $\F$ est la collection des événements observables ;
  \item $(\F_t)$ représente l'information disponible au temps $t$ (on dit \emph{filtration}) ;
  \item $\Prob$ est la probabilité ``physique'' (ou ``historique'') décrivant la réalité statistique.
\end{itemize}
\end{definition}

Intuition : $\F_t$ contient tout ce qu'un trader \emph{peut savoir} au temps $t$ (prix passés, news arrivées avant $t$, etc.).
Un processus $(X_t)$ est \emph{adapté} si $X_t$ est connu à l'instant $t$ (donc $X_t$ est $\F_t$-mesurable).

\begin{definition}[Martingale]
Un processus $(M_t)$ est une \emph{martingale} (sous $\Prob$) si, pour $s\le t$,
\[
\E\bracket{\,M_t\mid \F_s\,}=M_s.
\]
\end{definition}

Intuition : une martingale est un ``jeu équitable'' : le meilleur prédicteur du futur (conditionnellement à l'information) est le présent.
Dans un marché sans arbitrage, une idée clé est que les \emph{prix actualisés} deviennent des martingales sous une probabilité particulière (la mesure risque-neutre), ce qui est l'outil central du pricing.

\begin{example}[Prix forward sous taux et dividendes constants]
Supposons qu'un actif ne verse pas de dividende ($q=0$) et que l'on puisse emprunter/prêter au taux $r$ constant.
Un contrat forward paye $S_T - F_{0,T}$ à maturité.
En l'absence d'arbitrage, on peut répliquer le payoff en :
\begin{itemize}
\item achetant 1 unité de l'actif aujourd'hui, coût $S_0$ ;
\item finançant cet achat par emprunt au taux $r$ (donc on doit $S_0\ee^{rT}$ à $T$).
\end{itemize}
À maturité, la position vaut $S_T - S_0\ee^{rT}$.
Pour qu'un forward ait prix initial nul, il faut $F_{0,T}=S_0\ee^{rT}$.
Si l'actif verse un dividende continu $q$, la réplication donne $F_{0,T}=S_0\ee^{(r-q)T}$.
\end{example}

\begin{jurybox}
\begin{itemize}
\item Le cours est organisé comme une chaîne : données $\rightarrow$ modèle $\rightarrow$ pricing/greeks $\rightarrow$ calibration $\rightarrow$ couverture et décision.
\item Les notations $S_t$, $K$, $T$, $r$, $q$, $P(0,T)$ et $\sigma_{\mathrm{imp}}(K,T)$ doivent être sues \emph{sans hésiter}.
\item Une martingale formalise l'idée d'un ``prix juste'' : pas de dérive prédictible une fois l'information prise en compte.
\end{itemize}
\end{jurybox}

\part{Calcul stochastique \& pricing (socle indispensable)}

\chapter{Brownien et intégrale stochastique}
\label{chap:brownien}

\section{Pourquoi un Brownien en finance ?}
Le premier problème est de modéliser l'incertitude sur un prix $S_t$.
En pratique, les rendements sur des pas de temps courts ont :
\begin{itemize}
  \item une part imprévisible (news, flux, surprises) ;
  \item une agrégation de très nombreux micro-chocs ;
  \item des propriétés proches d'une somme de petits incréments.
\end{itemize}
Le \emph{mouvement brownien} est un modèle canonique pour représenter un bruit continu,
analogue à une limite de marche aléatoire.

\section{Définition formelle et propriétés clés}
\begin{definition}[Mouvement brownien standard]
Un processus $(W_t)_{t\ge 0}$ est un \emph{mouvement brownien standard} (adapté à $(\F_t)$) si :
\begin{enumerate}[label=\roman*)]
  \item $W_0=0$ ;
  \item les incréments sont indépendants : pour $0\le t_0 < t_1 < \dots < t_n$, les variables $W_{t_{k+1}}-W_{t_k}$ sont indépendantes ;
  \item les incréments sont stationnaires et gaussiens :
  \[
  W_t - W_s \sim \Normal(0, t-s)\quad \text{pour } 0\le s<t ;
  \]
  \item les trajectoires sont continues (presque sûrement).
\end{enumerate}
\end{definition}

\subsection{Intuition : marche aléatoire et limite}
Considérons une marche aléatoire discrète :
\[
X_n=\sum_{k=1}^n \varepsilon_k,
\qquad \varepsilon_k \in \{-1,+1\}\ \text{i.i.d. avec } \Prob(\varepsilon_k=1)=\Prob(\varepsilon_k=-1)=\tfrac12.
\]
On a $\E[X_n]=0$ et $\Var(X_n)=n$.
En temps continu, on redéfinit $t=n\Delta t$ et on rescale :
\[
W_t^{(\Delta t)}=\sqrt{\Delta t}\,X_{t/\Delta t}.
\]
Quand $\Delta t\to 0$, on peut montrer (idée du théorème central limite fonctionnel) que $W^{(\Delta t)}$ converge en loi vers un Brownien $W$.
Ce qui compte ici : \emph{la somme de petits chocs indépendants finit par ressembler à un bruit gaussien continu}.

\subsection{Moments simples}
Du fait de la loi $\Normal(0,t)$ :
\[
\E[W_t]=0,\qquad \Var(W_t)=t,\qquad \Cov(W_s,W_t)=\min(s,t).
\]
La covariance $\min(s,t)$ encode une mémoire simple : le chemin jusqu'à $s$ est commun aux deux, puis les incréments après $s$ sont indépendants.

\section{Intégrale stochastique d'Itô}
\label{sec:ito_integrale}

En finance, on manipule souvent des quantités du type ``\emph{exposition} $\times$ \emph{bruit}''.
Exemple : $dS_t = \sigma S_t\,\dd W_t$ signifie que l'incrément de prix est proportionnel au niveau $S_t$ (effet multiplicatif) et au bruit $\dd W_t$.
Pour donner un sens mathématique à $\int_0^T H_t\,\dd W_t$, il faut définir une intégrale qui respecte la nature très irrégulière du Brownien.

\subsection{Construction (idée) par approximations simples}
On commence par les processus en escalier :
\[
H_t = \sum_{k=0}^{n-1} H_{t_k}\,\1_{(t_k,t_{k+1}]}(t),
\qquad 0=t_0<\dots<t_n=T,
\]
avec $H_{t_k}$ mesurable par rapport à $\F_{t_k}$ (donc ``connu'' au début de l'intervalle).
On définit alors
\[
\int_0^T H_t\,\dd W_t \;\triangleq\; \sum_{k=0}^{n-1} H_{t_k}\bigl(W_{t_{k+1}}-W_{t_k}\bigr).
\]
L'idée : sur chaque intervalle, on gèle $H$ au début (convention d'Itô), puis on multiplie par l'incrément de Brownien.

Ensuite, pour des processus plus généraux (par exemple continus, ou carrés-intégrables), on prend une limite en $L^2$.

\subsection{Isométrie d'Itô : le ``Pythagore'' du bruit}
\begin{theorem}[Isométrie d'Itô (forme simple)]
Si $H$ est adapté et vérifie $\E\left[\int_0^T H_t^2\,\dd t\right]<\infty$, alors
\[
\E\left[\left(\int_0^T H_t\,\dd W_t\right)^2\right] = \E\left[\int_0^T H_t^2\,\dd t\right].
\]
\end{theorem}

\textbf{À quoi ça sert ?}
Cette formule est fondamentale pour :
\begin{itemize}
  \item contrôler la variance d'une stratégie de couverture (car les termes en $\dd W$ créent l'aléa) ;
  \item prouver que certains processus sont des martingales ;
  \item analyser numériquement des schémas de simulation (erreur quadratique).
\end{itemize}

\subsection{Martingales associées}
Une conséquence importante : si $H$ est adapté et intégrable,
\[
M_t=\int_0^t H_s\,\dd W_s
\]
est une martingale (sous $\Prob$). Intuition : l'intégrale accumule un bruit centré ; conditionnellement au passé, l'espérance des futurs incréments est nulle.

\begin{example}[Intégrale d'une constante]
Pour $H_t\equiv \sigma$, on obtient
\[
\int_0^T \sigma\,\dd W_t = \sigma (W_T - W_0)=\sigma W_T \sim \Normal(0,\sigma^2 T).
\]
La variance est $\sigma^2 T$, ce qui est cohérent avec l'isométrie d'Itô :
$\E[(\int_0^T \sigma \dd W)^2]=\int_0^T \sigma^2 \dd t=\sigma^2 T$.
\end{example}

\section{Limites et ce que le Brownien ne capture pas}
Le Brownien est \emph{continu} et \emph{gaussien}. Or les marchés réels présentent :
\begin{itemize}
  \item des queues plus épaisses que la normale (kurtosis) ;
  \item des sauts (annonces macro, gaps) ;
  \item de la volatilité stochastique et clustering ;
  \item parfois des corrélations de volatilité à court terme non triviales (``roughness'').
\end{itemize}
Ces phénomènes motivent les modèles enrichis étudiés plus tard (sauts, Heston, SABR, rough volatility).

\begin{jurybox}
\begin{itemize}
\item Un Brownien $W_t$ est la limite naturelle d'une marche aléatoire rescalée ; c'est le ``bruit de base'' des modèles continus.
\item L'intégrale d'Itô $\int H_t\,\dd W_t$ est définie en gelant $H$ au début de chaque intervalle.
\item L'isométrie d'Itô donne directement la variance : c'est l'outil de contrôle du risque des termes stochastiques.
\item Le Brownien ne capture ni sauts ni volatilité variable : d'où les modèles de volatilité (Partie II).
\end{itemize}
\end{jurybox}

\chapter{Variation quadratique et lemme d'Itô}
\label{chap:ito}

\section{Variation quadratique : la quantité qui remplace la dérivée}
\label{sec:qv}

Un point choc (mais essentiel) : un Brownien est continu mais \emph{n'est pas dérivable}.
Pour faire du calcul différentiel, on a besoin d'un concept qui mesure la ``rugosité'' d'un chemin.
C'est le rôle de la \emph{variation quadratique}.

\begin{definition}[Variation quadratique]
Soit $(X_t)$ un processus continu.
Pour une partition $\pi=\{0=t_0<\dots<t_n=T\}$, posez
\[
V^{(2)}(X;\pi)=\sum_{k=0}^{n-1} (X_{t_{k+1}}-X_{t_k})^2.
\]
Si, quand le pas de la partition $|\pi|=\max_k(t_{k+1}-t_k)\to 0$, la quantité $V^{(2)}(X;\pi)$ converge en probabilité vers une limite, on appelle cette limite la \emph{variation quadratique} et on la note
\[
[X]_T.
\]
\end{definition}

\subsection{Exemples fondamentaux}
\begin{proposition}[Variation quadratique du Brownien]
Pour un Brownien standard $(W_t)$,
\[
[W]_T = T.
\]
\end{proposition}

\textbf{Intuition :} sur un petit intervalle $\Delta t$, on a $\Delta W \sim \Normal(0,\Delta t)$ donc typiquement $(\Delta W)^2$ est de l'ordre $\Delta t$.
En sommant $n$ termes de taille $\Delta t$ sur $[0,T]$, on obtient une quantité de taille $T$.

\begin{proposition}[Processus de variation finie]
Si $A_t$ est un processus de variation finie (par exemple $A_t=\int_0^t a_s\,\dd s$ avec $a$ intégrable),
alors
\[
[A]_T=0.
\]
\end{proposition}

\textbf{À quoi ça sert ?}
Cela sépare deux ``types'' d'incréments :
\begin{itemize}
\item les incréments de type $\dd t$ (variation finie) ;
\item les incréments de type $\dd W$ (rugosité), dont le carré donne aussi du $\dd t$ via la variation quadratique.
\end{itemize}

\section{La règle \texorpdfstring{$(\dd W)^2=\dd t$}{(dW)^2=dt} et ses conséquences}
\label{sec:dw2dt}

Dans les calculs heuristiques d'Itô, on utilise les ordres de grandeur :
\[
\Delta W = O(\sqrt{\Delta t}),\qquad (\Delta W)^2 = O(\Delta t),\qquad (\Delta t)^2 = o(\Delta t),\qquad \Delta t\,\Delta W=o(\Delta t).
\]
En langage de différentiels :
\[
(\dd W_t)^2 = \dd t,\qquad \dd t\,\dd W_t = 0,\qquad (\dd t)^2=0.
\]
Ces règles ne sont pas de la magie : elles viennent exactement de la variation quadratique.

\section{Dérivation détaillée du lemme d'Itô (1D)}
\label{sec:derive_ito}

\begin{theorem}[Lemme d'Itô (1D)]
Soit $(X_t)$ solution de l'équation différentielle stochastique
\[
\dd X_t = a(t,X_t)\,\dd t + b(t,X_t)\,\dd W_t,
\]
et soit $f(t,x)$ une fonction $C^{1,2}$ (une dérivée en $t$, deux en $x$).
Alors le processus $Y_t=f(t,X_t)$ vérifie
\[
\dd Y_t = \pd{f}{t}(t,X_t)\,\dd t + \pd{f}{x}(t,X_t)\,\dd X_t + \frac12 \pdd{f}{x}(t,X_t)\,(\dd X_t)^2,
\]
avec la règle $(\dd X_t)^2=b(t,X_t)^2\,\dd t$.
En développant :
\[
\dd f(t,X_t)
=
\left(\pd{f}{t} + a\,\pd{f}{x} + \frac12 b^2\,\pdd{f}{x}\right)\dd t
+ b\,\pd{f}{x}\,\dd W_t.
\]
\end{theorem}

\subsection{Preuve (pas à pas, niveau M1)}
On considère une partition $0=t_0<\dots<t_n=T$ et on note $\Delta X_k=X_{t_{k+1}}-X_{t_k}$, $\Delta t_k=t_{k+1}-t_k$.
Sur chaque intervalle, on fait un Taylor à l'ordre 2 en $x$ et à l'ordre 1 en $t$ :
\[
f(t_{k+1},X_{t_{k+1}})-f(t_k,X_{t_k})
=
f_t(t_k,X_{t_k})\,\Delta t_k
+ f_x(t_k,X_{t_k})\,\Delta X_k
+ \frac12 f_{xx}(t_k,X_{t_k})\,(\Delta X_k)^2
+ \text{reste}.
\]
Le reste contient des termes d'ordre $(\Delta t_k)^2$, $\Delta t_k \Delta X_k$, et $(\Delta X_k)^3$ (plus précisément, il est $o(\Delta t_k)+o((\Delta X_k)^2)$).
Or $\Delta X_k \approx a\,\Delta t_k + b\,\Delta W_k$.
Donc :
\[
(\Delta X_k)^2
= a^2(\Delta t_k)^2 + 2ab\,\Delta t_k\,\Delta W_k + b^2(\Delta W_k)^2.
\]
Quand on somme sur $k$ et que $|\pi|\to 0$ :
\begin{itemize}
  \item $\sum a^2(\Delta t_k)^2 \to 0$ (car $(\Delta t)^2$ est négligeable) ;
  \item $\sum 2ab\,\Delta t_k\,\Delta W_k \to 0$ en probabilité (ordre $\Delta t^{3/2}$) ;
  \item $\sum b^2 (\Delta W_k)^2 \to \int_0^T b(t,X_t)^2\,\dd t$ grâce à $[W]_T=T$.
\end{itemize}
Ainsi, le seul terme d'ordre 1 qui survit dans $(\Delta X_k)^2$ est $b^2(\Delta W)^2 \sim b^2 \Delta t$,
d'où le fameux $\tfrac12 f_{xx} b^2\,\dd t$.
En passant à la limite, on obtient l'identité différentielle.

\section{Version multidimensionnelle (utilité pratique)}
\label{sec:ito_multi}

\begin{theorem}[Itô multidimensionnel (énoncé)]
Soit $X_t\in\R^d$ solution de
\[
\dd X_t = a(t,X_t)\,\dd t + \Sigma(t,X_t)\,\dd W_t,
\]
où $W_t\in\R^m$ est un Brownien $m$-dimensionnel et $\Sigma$ est une matrice $d\times m$.
Pour $f(t,x)$ de classe $C^{1,2}$,
\[
\dd f(t,X_t)
=
\pd{f}{t}\dd t
+ \nabla f^\top \dd X_t
+ \frac12 \mathrm{Tr}\!\left(\Sigma \Sigma^\top \nabla^2 f\right)\dd t.
\]
\end{theorem}

\textbf{À quoi ça sert ?}
Dès qu'on modélise plusieurs facteurs (spot + variance, taux + volatilité, etc.), c'est cette version qui s'applique.

\section{Exemples classiques en finance}
\label{sec:examples_ito}

\begin{example}[Logarithme d'une diffusion multiplicative]
Supposons $S_t$ tel que
\[
\dd S_t = \mu S_t\,\dd t + \sigma S_t\,\dd W_t.
\]
On pose $X_t=\log S_t$ et $f(s)=\log s$.
On a $f'(s)=1/s$ et $f''(s)=-1/s^2$.
Par Itô :
\[
\dd \log S_t
= \frac{1}{S_t}\dd S_t + \frac12\left(-\frac{1}{S_t^2}\right)(\dd S_t)^2.
\]
Or $(\dd S_t)^2=\sigma^2 S_t^2\,\dd t$.
Donc
\[
\dd \log S_t
= \mu\,\dd t + \sigma\,\dd W_t - \frac12\sigma^2\,\dd t
= \left(\mu-\frac12\sigma^2\right)\dd t + \sigma\,\dd W_t.
\]
Conclusion : $\log S_t$ est une gaussienne (donc $S_t$ est lognormale), base de Black--Scholes.
\end{example}

\begin{example}[Exponentielle martingale]
Pour un Brownien $W_t$ et une constante $\theta$, définissons
\[
M_t = \exp\!\left(\theta W_t - \frac12\theta^2 t\right).
\]
Avec $f(t,w)=\exp(\theta w-\frac12\theta^2 t)$, Itô donne
\[
\dd M_t = \theta M_t\,\dd W_t,
\]
donc $M_t$ est une martingale (car son drift est nul).
Ce type de construction est au cœur des changements de mesure (Girsanov) utilisés en pricing risque-neutre.
\end{example}

\section{Limites et hypothèses}
Le lemme d'Itô s'applique aux \emph{semi-martingales} (dont les diffusions classiques).
Il ne s'applique pas directement à certains modèles ``rough'' construits avec un Brownien fractionnaire ($H\neq 1/2$) : on y revient en Partie II.

\begin{jurybox}
\begin{itemize}
\item La variation quadratique est la clé qui remplace la dérivée : $[W]_T=T$.
\item Les règles $\dd t\,\dd W=0$ et $(\dd W)^2=\dd t$ sont la traduction calculatoire de la variation quadratique.
\item Le lemme d'Itô ajoute un terme $\tfrac12 f_{xx}b^2\,\dd t$ : c'est \emph{le} point qui distingue calcul stochastique et calcul classique.
\item En finance, il explique pourquoi une diffusion multiplicative mène à une loi lognormale (log de $S_t$ gaussien).
\end{itemize}
\end{jurybox}

\chapter{Pricing risque-neutre, réplication et PDE de Black--Scholes}
\label{chap:risk_neutral}

\section{Le problème financier}
Une option est un contrat qui dépend d'un sous-jacent $S_T$ à maturité.
Le problème est de déterminer un prix ``juste'' aujourd'hui $V_0$ :
\begin{itemize}
  \item \textbf{Objectif} : éviter l'arbitrage (pas de profit sans risque et sans capital).
  \item \textbf{Entrées} : dynamique du sous-jacent, taux, dividendes, payoff.
  \item \textbf{Sorties} : prix $V_t$ et sensibilités (greeks) pour couvrir.
\end{itemize}

Deux philosophies coexistent mais sont en fait équivalentes (en marchés complets) :
\begin{enumerate}[label=\alph*)]
  \item \emph{Réplication} : construire un portefeuille auto-financé qui reproduit le payoff.
  \item \emph{Risque-neutre} : calculer une espérance sous une probabilité $\Q$ où les actifs actualisés sont des martingales.
\end{enumerate}

\section{Modèle Black--Scholes (hypothèses)}
On suppose :
\begin{itemize}
  \item un marché \textbf{frictionless} (pas de coûts, pas de contraintes) ;
  \item possibilité de trader en continu ;
  \item taux sans risque constant $r$ et dividend yield constant $q$ ;
  \item sous-jacent diffusif :
  \[
  \dd S_t = (r-q)S_t\,\dd t + \sigma S_t\,\dd W_t^{\Q}.
  \]
\end{itemize}
La présence de $r-q$ dans le drift sous $\Q$ sera justifiée ci-dessous.
Sous la probabilité historique $\Prob$, on écrirait plutôt $\dd S_t=\mu S_t\dd t + \sigma S_t\dd W_t^{\Prob}$ avec $\mu$ inconnu et variable ; la magie du risque-neutre est de l'éliminer du pricing.

\section{Dérivation PDE par portefeuille de couverture}
\label{sec:derive_pde}

Soit $V(t,S)$ le prix d'une option européenne (fonction suffisamment régulière).
Appliquons Itô à $V(t,S_t)$ :
\[
\dd V
=
\pd{V}{t}\dd t + \pd{V}{S}\dd S + \frac12 \pdd{V}{S}\,(\dd S)^2.
\]
Or $(\dd S)^2 = \sigma^2 S^2 \dd t$ (car $(\dd W)^2=\dd t$), donc
\[
\dd V
=
\left(\pd{V}{t} + (r-q)S\pd{V}{S} + \frac12\sigma^2 S^2 \pdd{V}{S}\right)\dd t
+ \sigma S \pd{V}{S}\,\dd W^{\Q}.
\]

\subsection{Portefeuille auto-financé}
Considérons un portefeuille
\[
\Pi_t = V(t,S_t) - \Delta_t S_t,
\]
c'est-à-dire : \emph{long} une option, \emph{short} $\Delta_t$ actions.
Son différentiel est
\[
\dd \Pi = \dd V - \Delta\,\dd S \quad \text{(auto-financement : on ajuste $\Delta$ en réinvestissant dans le portefeuille).}
\]
Substitution :
\[
\dd \Pi
=
\left(\pd{V}{t} + (r-q)S\pd{V}{S} + \frac12\sigma^2 S^2 \pdd{V}{S}\right)\dd t
+ \sigma S \pd{V}{S}\,\dd W
- \Delta\bigl((r-q)S\dd t + \sigma S\,\dd W\bigr).
\]
Regroupons :
\[
\dd \Pi
=
\left(\pd{V}{t} + (r-q)S\pd{V}{S} + \frac12\sigma^2 S^2 \pdd{V}{S} - \Delta (r-q)S\right)\dd t
+ \sigma S\bigl(\pd{V}{S} - \Delta\bigr)\dd W.
\]

\subsection{Choix du hedge : annuler le risque instantané}
Choisissons
\[
\Delta = \pd{V}{S}.
\]
Alors le terme en $\dd W$ disparaît :
\[
\dd \Pi
=
\left(\pd{V}{t} + \frac12\sigma^2 S^2 \pdd{V}{S}\right)\dd t.
\]
Le portefeuille est (instantanément) \emph{sans risque}.
En absence d'arbitrage, un actif sans risque doit rapporter le taux sans risque :
\[
\dd \Pi = r \Pi\,\dd t = r\bigl(V - S\pd{V}{S}\bigr)\dd t.
\]
On obtient la PDE de Black--Scholes :
\[
\boxed{
\pd{V}{t} + \frac12\sigma^2 S^2 \pdd{V}{S} + (r-q)S\pd{V}{S} - rV = 0
}
\]
avec condition terminale $V(T,S)=\text{payoff}(S)$.

\textbf{À quoi ça sert ?}
La PDE donne :
\begin{itemize}
  \item un cadre de valorisation (prix cohérents sans arbitrage) ;
  \item une façon de calculer $V$ (analyse/PDE ou méthode numérique) ;
  \item la stratégie de couverture $\Delta=\partial_S V$.
\end{itemize}

\section{Mesure risque-neutre et prix comme espérance}
\label{sec:rn_expectation}

La PDE est un objet analytique. En finance, on préfère souvent une formule d'espérance :
\[
V_0 = \E^{\Q}\left[\ee^{-\int_0^T r_u\,\dd u}\,\text{payoff}(S_T)\right].
\]
Avec $r$ constant : $V_0=\ee^{-rT}\E^{\Q}[\text{payoff}(S_T)]$.

\subsection{Idée de Girsanov (sans preuves lourdes)}
Sous la probabilité historique $\Prob$ :
\[
\frac{\dd S_t}{S_t} = \mu\,\dd t + \sigma\,\dd W_t^{\Prob}.
\]
On définit le \emph{prix du risque} (market price of risk)
\[
\lambda = \frac{\mu-(r-q)}{\sigma}.
\]
Il existe une nouvelle probabilité $\Q$ telle que
\[
\dd W_t^{\Q} = \dd W_t^{\Prob} + \lambda\,\dd t
\]
est un Brownien sous $\Q$.
En remplaçant, on obtient
\[
\frac{\dd S_t}{S_t} = (r-q)\,\dd t + \sigma\,\dd W_t^{\Q}.
\]
\textbf{Conclusion :} sous $\Q$, le drift ``physique'' $\mu$ disparaît au profit de $r-q$.

\subsection{Feynman--Kac : PDE $\leftrightarrow$ espérance}
\begin{theorem}[Feynman--Kac (version utile ici)]
Considérons la PDE
\[
\pd{V}{t} + \mathcal{L}V - rV = 0,\qquad V(T,\cdot)=g(\cdot),
\]
où $\mathcal{L}$ est le générateur de la diffusion $S_t$ sous $\Q$.
Alors la solution s'écrit
\[
V(t,S_t) = \E^{\Q}\left[\exp\!\left(-\int_t^T r\,\dd u\right) g(S_T)\,\Big|\,\F_t\right].
\]
\end{theorem}

Intuition : la PDE impose exactement que le processus actualisé $\ee^{-\int_0^t r\,\dd u}V(t,S_t)$ soit une martingale,
ce qui force sa valeur actuelle à être l'espérance conditionnelle de sa valeur terminale actualisée.

\begin{example}[Prix d'un zéro-coupon en taux constant]
Un zéro-coupon paye 1 à $T$.
Sous taux constant $r$, sa valeur est
\[
P(0,T)=\ee^{-rT}.
\]
Ce résultat se retrouve :
\begin{itemize}
\item par réplication (placer 1 euro à $r$) ;
\item par Feynman--Kac : $V_0=\E^\Q[\ee^{-rT}\cdot 1]=\ee^{-rT}$.
\end{itemize}
\end{example}

\section{Limites du cadre Black--Scholes}
Les hypothèses qui ``cassent'' en pratique :
\begin{itemize}
  \item volatilité constante $\sigma$ (les marchés ont une surface $\sigma_{\mathrm{imp}}(K,T)$) ;
  \item trading continu sans coûts (en réalité hedging discret + coûts) ;
  \item distribution lognormale (sous-estime les queues et l'asymétrie) ;
  \item marché complet (dès qu'on ajoute des facteurs non tradables, la réplication parfaite peut échouer).
\end{itemize}
Ces limites motivent la calibration de surfaces et des modèles plus riches (Partie II) et le hedging discret/risk-aware (Partie III).

\begin{jurybox}
\begin{itemize}
\item La PDE Black--Scholes vient d'un portefeuille \emph{option} $-$ $\Delta$ \emph{actions} avec $\Delta=\partial_S V$.
\item En absence d'arbitrage, un portefeuille sans risque doit croître au taux $r$ : cela force la PDE.
\item La mesure risque-neutre $\Q$ remplace le drift $\mu$ par $r-q$ et permet une formule d'espérance.
\item PDE et espérance sont deux faces de la même pièce (Feynman--Kac).
\end{itemize}
\end{jurybox}

\chapter{Changement de mesure \& numéraires : pourquoi $\Q$ n'est pas une ``astuce''}
\label{chap:measure_change}

Ce chapitre approfondit la mécanique derrière la mesure risque-neutre.
En M1, on voit souvent apparaître $\Q$ comme une ``mesure magique'' qui permet d'écrire
\(V_0=\E^{\Q}[\text{payoff actualisé}]\). L'objectif ici est de comprendre précisément :
\begin{enumerate}[label=\alph*)]
\item \emph{ce qui change} quand on passe de $\P$ à $\Q$ ;
\item \emph{pourquoi} ce changement est nécessaire (absence d'arbitrage) ;
\item \emph{comment} choisir un numéraire adapté (zéro-coupon, money-market, forward measure) ;
\item \emph{quelles limites} : marchés incomplets, risques non couvrables, choix non-unique de mesure.
\end{enumerate}

\section{Mesure physique $\P$ vs mesure de pricing}

\subsection{Deux questions différentes}
\begin{itemize}
\item \textbf{Prévision (statistique)} : ``Quelle est la distribution future de $S_T$ ?''
\quad $\Rightarrow$ on travaille naturellement sous $\P$ (mesure ``réelle'').
\item \textbf{Valorisation (absence d'arbitrage)} : ``Quel est le prix aujourd'hui d'un payoff $g(S_T)$ ?''
\quad $\Rightarrow$ on veut un prix cohérent avec la possibilité de couverture/ réplication.
\end{itemize}

Le point clé : \emph{le prix d'un dérivé n'est pas l'espérance sous $\P$},
car un prix doit incorporer l'actualisation et la rémunération du risque.
Dans un marché complet (comme Black--Scholes), l'absence d'arbitrage force un prix unique,
et ce prix s'écrit comme une espérance sous une mesure de martingale.

\subsection{Le rôle de la prime de risque}
Sous $\P$, la dynamique d'un sous-jacent risqué (avec dividendes continus $q$) s'écrit typiquement
\[\frac{\dd S_t}{S_t} = \mu\,\dd t + \sigma\,\dd W_t^{\P}.\]
Le paramètre $\mu$ \emph{n'est pas} observable de manière stable à partir des prix instantanés :
il mélange anticipation et prime de risque.
En revanche, les options liquides révèlent principalement \emph{la distribution risque-neutre}
et donc la volatilité implicite, beaucoup plus stable à court terme.

\section{Numéraire et principe de martingale}

\subsection{Définition}
\begin{definition}[Numéraire]
Un \emph{numéraire} $N_t$ est un actif strictement positif utilisé comme unité de compte.
On exprime les prix relatifs $\widetilde{X}_t = X_t/N_t$.
\end{definition}

Exemples :
\begin{itemize}
\item \textbf{Money-market} $B_t = \exp(\int_0^t r_u\,\dd u)$ (compte bancaire).
\item \textbf{Zéro-coupon} $P(t,T)$ comme numéraire (utile pour produits de taux).
\item \textbf{Actif risqué} $S_t$ comme numéraire (utile pour certains changements de mesure en FX).
\end{itemize}

\subsection{Théorème (version utilisable en M1)}
\begin{theorem}[Principe de martingale associé à un numéraire]
Supposons absence d'arbitrage.
Alors, pour un numéraire $N_t$, il existe une mesure $\Q^{N}$ telle que
\[\frac{X_t}{N_t}\quad\text{est une martingale sous }\Q^{N}\]
pour tout actif admissible $X$ (ou au moins pour la classe de payoffs considérés).
En particulier, le prix à $t$ d'un payoff $X_T$ s'écrit
\[
X_t = N_t\,\E^{\Q^{N}}\left[\frac{X_T}{N_T}\,\Big|\,\F_t\right].
\]
\end{theorem}

\textbf{À quoi ça sert ?}
\begin{itemize}
\item Choisir le bon numéraire simplifie souvent les calculs (par exemple, sous la \emph{forward measure}, un forward devient une martingale).
\item Cela clarifie que ``la mesure de pricing'' dépend d'un choix d'unité de compte.
\end{itemize}

\section{Girsanov (intuition + formule)}

\subsection{Intuition : absorber un drift}
Dans un modèle de diffusion, changer de mesure revient à \emph{modifier le drift} du Brownien.
On garde la même volatilité, mais on ``décale'' le mouvement aléatoire pour que des prix actualisés deviennent martingales.

\subsection{Énoncé (cas 1D)}
Soit $W_t^{\P}$ un Brownien sous $\P$ et soit un processus adapté $(\theta_t)$.
Définissons la densité de Radon--Nikodym
\[
\frac{\dd \Q}{\dd \P}\Big|_{\F_T} = \exp\Bigl(-\int_0^T \theta_u\,\dd W_u^{\P} - \tfrac12\int_0^T \theta_u^2\,\dd u\Bigr).
\]
Alors, sous $\Q$, le processus
\[
W_t^{\Q} := W_t^{\P} + \int_0^t \theta_u\,\dd u
\]
est un Brownien.

\textbf{Lecture financière :}
\(\theta\) est le \emph{prix du risque} (market price of risk) : il relie la dérive réelle $\mu$ au drift risque-neutre $r-q$.

\subsection{Application à Black--Scholes}
Sous $\P$ : \(\dd S_t = \mu S_t\,\dd t + \sigma S_t\,\dd W_t^{\P}\).
On veut sous $\Q$ : \(\dd S_t = (r-q) S_t\,\dd t + \sigma S_t\,\dd W_t^{\Q}\).
En identifiant :
\[
\dd W_t^{\Q} = \dd W_t^{\P} + \theta\,\dd t,
\qquad \text{avec }\theta = \frac{\mu-(r-q)}{\sigma}.
\]
On voit explicitement : passer à $\Q$ \emph{retire} la prime de risque $\mu-(r-q)$.

\begin{example}[Mini-exemple chiffré]
Prenons $\mu=10\%$, $r=3\%$, $q=1\%$, $\sigma=20\%$.
Alors $\theta = \frac{0.10-(0.03-0.01)}{0.20}=\frac{0.08}{0.20}=0.4$.
Sous $\Q$, le Brownien ``avance'' en moyenne de $-0.4\,\dd t$ par rapport à celui sous $\P$.
Ce n'est pas une correction ``petite'' : cela illustre pourquoi un prix ne peut pas être une espérance sous $\P$.
\end{example}

\section{Forward measure : quand un forward devient une martingale}
\label{sec:forward_measure}

\subsection{Forward price et numéraire}
Le \emph{forward} de maturité $T$ (sur action avec dividendes continus) vaut
\[F(t,T) = S_t\,\exp((r-q)(T-t)).\]
Sous la mesure $\Q$ associée au money-market, $F(t,T)$ n'est pas forcément une martingale car il contient explicitement la croissance à $r-q$.

Si l'on prend comme numéraire le zéro-coupon $P(t,T)$, on travaille sous la \emph{$T$-forward measure} $\Q^T$.
Dans ce cadre, des quantités exprimées en unités de $P(t,T)$ deviennent martingales.
Pour un produit dont le payoff est payé à $T$, cela simplifie :
\[
V_t = P(t,T)\E^{\Q^T}[\text{payoff} \mid \F_t].
\]

\subsection{À quoi ça sert concrètement ?}
\begin{itemize}
\item En \textbf{taux}, beaucoup d'instruments payent à une date précise (swaplets, caplets). La forward measure évite des termes d'actualisation complexes.
\item En \textbf{FX}, choisir comme numéraire l'actif domestique vs étranger clarifie le drift du taux de change.
\end{itemize}

\section{Limites : marchés incomplets et non-unicité de $\Q$}
\label{sec:incomplete_markets}

Dès qu'un modèle contient des risques non tradables (volatilité stochastique, sauts, illiquidité),
la réplication parfaite échoue : le marché devient \emph{incomplet}.
Alors :
\begin{itemize}
\item il peut exister une \emph{infinité} de mesures de martingale ;
\item le prix n'est plus unique sans un critère additionnel (préférence, risque, minimal martingale measure, entropie minimale, etc.).
\end{itemize}
Dans la pratique, la calibration à la surface d'options sélectionne implicitement une mesure de pricing ``compatible'' avec les prix.

\begin{figure}[h!]
\centering
\begin{tikzpicture}[node distance=10mm, font=\small]
  \node[draw, rectangle, align=center, minimum width=0.32\textwidth] (P) {$\P$\\\footnotesize dynamique ``réelle''\\\footnotesize drift $\mu$};
  \node[draw, rectangle, right=of P, align=center, minimum width=0.32\textwidth] (Q) {$\Q$\\\footnotesize pricing (numéraire $B_t$)\\\footnotesize drift $r-q$};
  \node[draw, rectangle, below=of Q, align=center, minimum width=0.32\textwidth] (QT) {$\Q^T$\\\footnotesize pricing (numéraire $P(t,T)$)\\\footnotesize ``forward measure''};
  \draw[-{Latex[length=2mm]}] (P) -- node[above]{Girsanov ($\theta$)} (Q);
  \draw[-{Latex[length=2mm]}] (Q) -- node[right]{changement de numéraire} (QT);
\end{tikzpicture}
\caption{Lecture opérationnelle : $\P$ pour prévoir, une mesure de martingale ($\Q$, $\Q^T$, etc.) pour pricer, en fonction du numéraire.}
\end{figure}

\begin{jurybox}
\begin{itemize}
\item $\P$ répond à ``que va-t-il se passer ?'' ; $\Q$ répond à ``combien ça vaut sans arbitrage ?''.
\item Un numéraire $N_t$ définit une mesure $\Q^N$ telle que $X_t/N_t$ est une martingale.
\item Girsanov montre explicitement que changer de mesure modifie le drift : $\theta=(\mu-(r-q))/\sigma$ en BS.
\item En marché incomplet, $\Q$ n'est plus unique : la calibration et/ou un critère de risque complète la théorie.
\end{itemize}
\end{jurybox}

\chapter{Options vanilles : formule de Black--Scholes et greeks}
\label{chap:bs_greeks}

\section{Payoffs européens et parité put--call}
\begin{definition}[Call/Put européens]
À maturité $T$, avec strike $K$ :
\[
\text{Call : } (S_T-K)^+ = \max(S_T-K,0),\qquad
\text{Put : } (K-S_T)^+ = \max(K-S_T,0).
\]
\end{definition}

\begin{proposition}[Parité put--call (avec dividendes continus)]
Pour des options européennes de même $K,T$,
\[
C_0 - P_0 = S_0 \ee^{-qT} - K\ee^{-rT}.
\]
\end{proposition}

\textbf{Intuition :} un call et un put diffèrent exactement d'un forward.
En effet, $\text{Call}-\text{Put} = S_T-K$ (vérifier cas par cas), donc en actualisant :
$C_0-P_0 = \text{PV}(S_T) - \text{PV}(K) = S_0\ee^{-qT}-K\ee^{-rT}$.

\section{Distribution lognormale sous $\Q$}
Sous Black--Scholes et sous la mesure risque-neutre :
\[
\frac{\dd S_t}{S_t} = (r-q)\dd t + \sigma \dd W_t^{\Q}.
\]
En appliquant Itô à $\log S_t$ (voir \cref{sec:examples_ito}) :
\[
\dd\log S_t = \left((r-q)-\frac12\sigma^2\right)\dd t + \sigma \dd W_t^{\Q}.
\]
En intégrant de $0$ à $T$ :
\[
\log S_T = \log S_0 + \left((r-q)-\frac12\sigma^2\right)T + \sigma W_T^{\Q}.
\]
Comme $W_T^{\Q}\sim \Normal(0,T)$, on obtient
\[
\log S_T \sim \Normal\!\left(\log S_0 + \left((r-q)-\frac12\sigma^2\right)T,\ \sigma^2 T\right).
\]

\section{Dérivation détaillée de la formule Black--Scholes (par espérance)}
\label{sec:derive_bs_expectation}

Le prix d'un call européen est
\[
C_0 = \ee^{-rT}\E^{\Q}\bracket{(S_T-K)^+}.
\]
On utilise l'identité $(S_T-K)^+ = S_T\1_{S_T>K}-K\1_{S_T>K}$ :
\[
C_0 = \ee^{-rT}\Bigl(\E^{\Q}[S_T\1_{S_T>K}] - K\,\Q(S_T>K)\Bigr).
\]
On va calculer séparément les deux termes.

\subsection{Étape 1 : calcul de $\Q(S_T>K)$}
Définissons
\[
Y = \log S_T.
\]
On sait que $Y\sim \Normal(m, s^2)$ avec
\[
m=\log S_0 + \left((r-q)-\frac12\sigma^2\right)T,\qquad s=\sigma\sqrt{T}.
\]
Alors
\[
\Q(S_T>K) = \Q(\log S_T > \log K) = \Q\left(\frac{Y-m}{s} > \frac{\log K - m}{s}\right).
\]
Si $Z=\frac{Y-m}{s}\sim \Normal(0,1)$, on obtient
\[
\Q(S_T>K) = 1-\Phi\left(\frac{\log K - m}{s}\right)=\Phi\left(\frac{m-\log K}{s}\right).
\]
En développant :
\[
\frac{m-\log K}{s}
=
\frac{\log(S_0/K) + \left((r-q)-\frac12\sigma^2\right)T}{\sigma\sqrt{T}}
\;\triangleq\; d_2.
\]
Donc
\[
\Q(S_T>K)=\Phi(d_2).
\]

\subsection{Étape 2 : calcul de $\E^{\Q}[S_T\1_{S_T>K}]$}
On écrit $S_T=\ee^{Y}$.
On veut $\E[\ee^Y\1_{Y>\log K}]$ où $Y$ est gaussien.
Une astuce standard consiste à \emph{compléter le carré}.

La densité de $Y$ est
\[
f_Y(y)=\frac{1}{s\sqrt{2\pi}}\exp\left(-\frac{(y-m)^2}{2s^2}\right).
\]
Alors
\[
\E[\ee^Y\1_{Y>\log K}]
=\int_{\log K}^{\infty}\ee^{y} f_Y(y)\,\dd y
=\int_{\log K}^{\infty}\frac{1}{s\sqrt{2\pi}}
\exp\left(y-\frac{(y-m)^2}{2s^2}\right)\dd y.
\]
Regardons l'exposant :
\[
y-\frac{(y-m)^2}{2s^2}
=
-\frac{1}{2s^2}\Bigl((y-m)^2-2s^2 y\Bigr).
\]
On réécrit le terme en $y$ :
\[
(y-m)^2-2s^2 y
=
y^2 -2my + m^2 -2s^2 y
=
y^2 -2(m+s^2)y + m^2.
\]
Complétons le carré :
\[
y^2 -2(m+s^2)y + m^2
=
(y-(m+s^2))^2 - (m+s^2)^2 + m^2.
\]
Donc
\[
y-\frac{(y-m)^2}{2s^2}
=
-\frac{(y-(m+s^2))^2}{2s^2} + \frac{(m+s^2)^2-m^2}{2s^2}.
\]
Or
\[
\frac{(m+s^2)^2-m^2}{2s^2} = \frac{2ms^2 + s^4}{2s^2}= m+\frac12 s^2.
\]
Ainsi
\[
\exp\left(y-\frac{(y-m)^2}{2s^2}\right)
=
\exp\left(m+\frac12 s^2\right)\cdot
\exp\left(-\frac{(y-(m+s^2))^2}{2s^2}\right).
\]
On obtient
\[
\E[\ee^Y\1_{Y>\log K}]
=
\exp\left(m+\frac12 s^2\right)\int_{\log K}^{\infty}\frac{1}{s\sqrt{2\pi}}
\exp\left(-\frac{(y-(m+s^2))^2}{2s^2}\right)\dd y.
\]
Mais l'intégrale est une probabilité gaussienne :
si $Y'\sim \Normal(m+s^2,s^2)$, alors
\[
\int_{\log K}^{\infty} f_{Y'}(y)\dd y = \Q(Y'>\log K)=\Phi\left(\frac{(m+s^2)-\log K}{s}\right).
\]
On définit
\[
d_1 = \frac{(m+s^2)-\log K}{s} = \frac{\log(S_0/K)+\left((r-q)+\frac12\sigma^2\right)T}{\sigma\sqrt{T}}.
\]
Donc
\[
\E^{\Q}[S_T\1_{S_T>K}] = \exp\left(m+\frac12 s^2\right)\Phi(d_1).
\]
Enfin, $\exp(m+\tfrac12 s^2)=\E[\ee^Y]=\E[S_T]=S_0\ee^{(r-q)T}$ sous $\Q$.
Ainsi
\[
\E^{\Q}[S_T\1_{S_T>K}] = S_0\ee^{(r-q)T}\Phi(d_1).
\]

\subsection{Formule finale}
On injecte dans $C_0$ :
\[
C_0 = \ee^{-rT}\Bigl(S_0\ee^{(r-q)T}\Phi(d_1) - K\Phi(d_2)\Bigr)
=
S_0\ee^{-qT}\Phi(d_1) - K\ee^{-rT}\Phi(d_2).
\]
\[
\boxed{
C_0 = S_0\ee^{-qT}\Phi(d_1) - K\ee^{-rT}\Phi(d_2),
\qquad
d_{1,2}=\frac{\log(S_0/K)+(r-q\pm \frac12\sigma^2)T}{\sigma\sqrt{T}}.
}
\]
Et par parité put--call :
\[
P_0 = K\ee^{-rT}\Phi(-d_2) - S_0\ee^{-qT}\Phi(-d_1).
\]

\section{Greeks : définitions, calculs et interprétation}
\label{sec:greeks}

\begin{definition}[Greeks principaux]
Pour une option de prix $V(S,t;\theta)$ (où $\theta$ regroupe les paramètres),
on définit :
\[
\Delta=\pd{V}{S},\quad
\Gamma=\pdd{V}{S},\quad
\vega=\pd{V}{\sigma},\quad
\Theta=\pd{V}{t},\quad
\rho=\pd{V}{r}.
\]
\end{definition}

\subsection{Delta et gamma (les plus importants pour le hedging)}
Pour un call BS :
\[
\Delta_{\text{call}} = \ee^{-qT}\Phi(d_1),
\qquad
\Delta_{\text{put}} = \ee^{-qT}\bigl(\Phi(d_1)-1\bigr).
\]
\textbf{Interprétation :} $\Delta$ est la quantité d'actions à short/long pour neutraliser le risque instantané $\dd W$.
C'est le hedge ``de premier ordre''.

La densité normale standard est $\varphi(x)=\frac{1}{\sqrt{2\pi}}\exp(-x^2/2)$.
On obtient
\[
\Gamma = \pd{\Delta}{S} = \frac{\ee^{-qT}\varphi(d_1)}{S\sigma\sqrt{T}}.
\]
\textbf{Interprétation :} $\Gamma$ mesure la variation de $\Delta$ quand $S$ bouge.
Un portefeuille ``gamma'' important nécessite de rebalancer souvent.

\subsection{Vega (sensibilité à la volatilité)}
Pour call ou put :
\[
\vega = \pd{V}{\sigma} = S_0\ee^{-qT}\varphi(d_1)\sqrt{T}.
\]
\textbf{Interprétation :} vega est maximale à la monnaie et augmente avec $\sqrt{T}$.
Elle est cruciale pour l'inversion de volatilité implicite (Newton) et la calibration.

\subsection{Theta et rho}
Les expressions explicites existent, mais ce qui compte surtout en pratique :
\begin{itemize}
\item $\Theta$ capture l'érosion temporelle : à paramètres fixes, un call perd de la valeur quand $T$ diminue (sauf effets particuliers).
\item $\rho$ mesure l'effet du taux sur la valeur actualisée du strike.
\end{itemize}

\begin{example}[Exemple chiffré simple]
Prenons $S_0=100$, $K=100$, $T=1$, $r=3\%$, $q=0$, $\sigma=20\%$.
Alors
\[
d_1=\frac{\log(1)+(0.03+0.5\cdot 0.2^2)}{0.2}= \frac{0.03+0.02}{0.2}=0.25,
\qquad
d_2=d_1-0.2=0.05.
\]
Avec $\Phi(0.25)\approx 0.5987$ et $\Phi(0.05)\approx 0.5199$ :
\[
C_0 \approx 100\cdot 0.5987 - 100\ee^{-0.03}\cdot 0.5199
\approx 59.87 - 50.44 \approx 9.43.
\]
Delta call $\approx 0.5987$.
Gamma $\approx \frac{\varphi(0.25)}{100\cdot 0.2} \approx \frac{0.3867}{20} \approx 0.0193$.
Vega $\approx 100\varphi(0.25)\sqrt{1} \approx 38.67$ (par unité de volatilité, i.e. pour $\sigma$ en absolu).
\end{example}

\section{Limites pratiques des greeks BS}
Les greeks Black--Scholes sont \emph{modèle-dépendants}.
Deux sources d'erreur dominantes :
\begin{itemize}
  \item la volatilité n'est pas constante : la vraie sensibilité est à la surface $\sigma_{\mathrm{imp}}(K,T)$ ;
  \item hedging discret et coûts : même un $\Delta$ parfait instantanément ne garantit pas un $\PnL$ neutre à horizon fini.
\end{itemize}

\begin{jurybox}
\begin{itemize}
\item Sous Black--Scholes, $\log S_T$ est normal : cela permet une dérivation par espérance sans résoudre la PDE.
\item Formule call : $C_0 = S_0\ee^{-qT}\Phi(d_1)-K\ee^{-rT}\Phi(d_2)$, avec $d_1-d_2=\sigma\sqrt{T}$.
\item $\Delta=\partial_S V$ est la quantité d'actions pour neutraliser le risque instantané ; $\Gamma$ mesure la courbure et le besoin de rebalancing.
\item $\vega$ est centrale pour la volatilité implicite et la calibration.
\end{itemize}
\end{jurybox}

\chapter{Volatilité implicite : définition, inversion numérique et stabilité}
\label{chap:implied_vol}

\section{Pourquoi introduire la volatilité implicite ?}
Les desks options ne quotent pas directement des prix en euros sur chaque strike/maturité,
car :
\begin{itemize}
  \item les niveaux de spot et de taux changent, ce qui déforme les prix ;
  \item les options ITM/OTM ont des prix très différents, peu comparables ;
  \item la volatilité est un paramètre plus ``stable'' et plus parlant (niveau de risque).
\end{itemize}
On convertit donc un prix de marché en une \emph{volatilité implicite} : la volatilité BS qui reproduit ce prix.

\section{Définition formelle}
\begin{definition}[Volatilité implicite]
Soit $C_{\mathrm{mkt}}(K,T)$ un prix de call observé (après ajustements éventuels).
La \emph{volatilité implicite} $\sigma_{\mathrm{imp}}(K,T)$ est la solution (si elle existe) de
\[
C_{\mathrm{mkt}}(K,T) = \BS_{\mathrm{call}}(S_0,K,r,q,T;\sigma_{\mathrm{imp}}).
\]
De même pour un put.
\end{definition}

\subsection{Existence et unicité : monotonicité en $\sigma$}
En Black--Scholes, pour un call ou put européen,
\[
\pd{V}{\sigma} = \vega > 0
\]
(sauf limites dégénérées).
Donc $V(\sigma)$ est strictement croissant en $\sigma$ : si un prix mkt est compatible avec l'intervalle des prix BS, la solution est unique.

\begin{remark}[Bornes de non-arbitrage (rappel utile)]
Pour un call européen (avec dividendes continus), les bornes simples sont
\[
\max(0, S_0\ee^{-qT}-K\ee^{-rT}) \le C_0 \le S_0\ee^{-qT}.
\]
Si un prix observé viole ces bornes, l'``implied vol'' n'existe pas (ou signale un problème de données).
\end{remark}

\section{Inversion numérique : bisection, Newton, et pièges}
On veut résoudre l'équation $f(\sigma)=\BS(\sigma)-C_{\mathrm{mkt}}=0$.

\subsection{Méthode de dichotomie (bisection)}
\textbf{Principe :} si $f$ est continue et monotone, on encadre la racine dans $[\sigma_{\min},\sigma_{\max}]$.
On itère :
\[
\sigma_{\mathrm{mid}}=\frac{\sigma_{\min}+\sigma_{\max}}{2},
\]
et on remplace l'intervalle par la moitié qui contient le changement de signe.
\textbf{Avantage :} robuste (converge toujours).
\textbf{Inconvénient :} convergence linéaire (plus lente).

\subsection{Méthode de Newton}
\textbf{Principe :} linéariser $f$ autour de $\sigma_n$ :
\[
f(\sigma)\approx f(\sigma_n)+f'(\sigma_n)(\sigma-\sigma_n).
\]
On impose $f(\sigma_{n+1})\approx 0$, donc
\[
\boxed{
\sigma_{n+1} = \sigma_n - \frac{f(\sigma_n)}{f'(\sigma_n)}
= \sigma_n - \frac{\BS(\sigma_n)-C_{\mathrm{mkt}}}{\vega(\sigma_n)}.
}
\]
\textbf{Avantage :} convergence quadratique près de la solution.
\textbf{Inconvénient :} peut diverger si l'initialisation est mauvaise ou si $\vega$ est très petite.

\subsection{Heuristiques d'initialisation}
Un choix typique :
\begin{itemize}
  \item commencer à $\sigma_0=20\%$ ou $\sigma_0=\sqrt{\frac{2\abs{\log(S_0/K)}}{T}}$ (ordre de grandeur) ;
  \item borner $\sigma$ dans $[\sigma_{\min},\sigma_{\max}]$ (par exemple $[10^{-4},5]$ en absolu) ;
  \item combiner Newton avec bisection (``safeguarded Newton'') : si Newton sort de l'intervalle, on fait une étape de bisection.
\end{itemize}

\subsection{Pseudo-code d'un Newton sécurisé}
\begin{lstlisting}[caption={Newton sécurisé pour volatilité implicite (pseudo-code).},label={lst:newton_iv}]
Input: price C_mkt, initial sigma0, bounds [a,b]
for n = 1..Nmax:
    C = BS_price(sigma)
    v = BS_vega(sigma)
    if abs(C - C_mkt) < tol: break
    sigma_new = sigma - (C - C_mkt)/max(v, v_floor)
    if sigma_new < a or sigma_new > b:
        sigma_new = 0.5*(a+b)      # fallback: bisection
    # update bracket using monotonicity
    if BS_price(sigma_new) > C_mkt: b = sigma_new else a = sigma_new
    sigma = sigma_new
return sigma
\end{lstlisting}

\section{Stabilité : quand l'implied vol devient ``mauvais'' numériquement}
La difficulté majeure vient de $\vega$ :
\[
\vega = S_0\ee^{-qT}\varphi(d_1)\sqrt{T}.
\]
Donc $\vega$ est petite :
\begin{itemize}
  \item pour des maturités très courtes ($\sqrt{T}$) ;
  \item pour des options très ITM/OTM (car $\varphi(d_1)$ devient minuscule).
\end{itemize}
Dans ces zones, un écart de prix minuscule correspond à une variation énorme de $\sigma$ : l'inversion est \emph{mal conditionnée}.
Il faut donc :
\begin{itemize}
  \item contrôler la qualité des prix (bid/ask, erreurs) ;
  \item préférer des métriques robustes (ex. travailler en prix plutôt qu'en IV pour la calibration, ou pondérer par vega) ;
  \item utiliser des bornes et des méthodes robustes.
\end{itemize}

\begin{example}[Illustration de mal-conditionnement]
Considérons une option très OTM : $S_0=100$, $K=200$, $T=0.05$ (environ 12 jours), $r=q=0$.
Le prix est très proche de 0 ; deux feeds peuvent donner $C_{\mathrm{mkt}}=0.0005$ ou $0.0010$.
Ces deux prix correspondent à des volatilités implicites très différentes, car la courbe $C(\sigma)$ est quasi plate près de 0.
Conclusion : l'implied vol ``explose'' en bruit.
\end{example}

\begin{jurybox}
\begin{itemize}
\item La volatilité implicite est définie par inversion de Black--Scholes ; elle n'est pas une volatilité ``réelle''.
\item Unicité : $\BS(\sigma)$ est croissante en $\sigma$ car $\vega>0$.
\item Newton : $\sigma_{n+1}=\sigma_n-(\BS(\sigma_n)-C_{\mathrm{mkt}})/\vega(\sigma_n)$, rapide mais fragile.
\item Zones dangereuses : très court terme et extrêmes strikes (vega faible) $\Rightarrow$ inversion mal conditionnée.
\end{itemize}
\end{jurybox}

\part{Surface de volatilité \& calibration (cœur ``quant'')}

\chapter{Surface d'implied volatility : construction, smile/skew et arbitrage}
\label{chap:surface_iv}

\section{Pourquoi une surface existe-t-elle ? (smile, skew)}
Dans Black--Scholes, une seule volatilité $\sigma$ suffit pour tous les strikes et maturités.
En réalité, si l'on calcule $\sigma_{\mathrm{imp}}(K,T)$ pour chaque option cotée, on observe :
\begin{itemize}
  \item un \textbf{smile} (vol plus élevée loin de la monnaie) ;
  \item un \textbf{skew} (asymétrie : vol plus forte pour puts OTM sur indices actions) ;
  \item une dépendance à la maturité (\emph{term structure}).
\end{itemize}
Ces phénomènes reflètent des caractéristiques réelles des rendements :
asymétrie (crashes plus violents), queues épaisses, volatilité stochastique, sauts,
et effets d'offre/demande (couverture des investisseurs).

\section{Coordonnées de la surface : strike, moneyness, delta}
\label{sec:moneyness}
Plusieurs axes sont utilisés en pratique :
\begin{itemize}
  \item \textbf{Strike} $K$ : naturel mais peu comparable quand $S_0$ change.
  \item \textbf{Log-moneyness} $k=\log(K/F_{0,T})$ : centré sur le forward $F_{0,T}$, plus stable.
  \item \textbf{Delta} : on indexe la smile par la delta BS (ex. 25-delta put/call), pratique pour FX.
\end{itemize}

\begin{definition}[Forward et log-moneyness]
Sous taux/dividendes constants, le forward est $F_{0,T}=S_0\ee^{(r-q)T}$.
On définit souvent
\[
k = \log\left(\frac{K}{F_{0,T}}\right).
\]
La monnaie correspond à $k=0$ (i.e. $K=F_{0,T}$).
\end{definition}

\section{No-arbitrage : contraintes simples et intuition}
\label{sec:noarb_surface}

Une surface n'est pas qu'un interpolateur : elle doit éviter de créer des opportunités d'arbitrage.
Deux familles d'arbitrages sont classiques : \textbf{butterfly} (en strike) et \textbf{calendar} (en maturité).

\subsection{Arbitrage en strike : convexité des prix}
Fixons $T$.
Le prix d'un call $C(K,T)$ comme fonction de $K$ doit vérifier :
\begin{itemize}
  \item $C$ décroît en $K$ (un strike plus haut vaut moins cher) ;
  \item $C$ est convexe en $K$ (car c'est une espérance d'une fonction convexe du strike).
\end{itemize}
Formellement, si $C$ est deux fois dérivable en $K$ :
\[
\pd{C}{K}\le 0,\qquad \pdd{C}{K}\ge 0.
\]
De plus, la densité risque-neutre $f_{S_T}(K)$ est donnée par :
\[
\pdd{C}{K}(K,T)=\ee^{-rT} f_{S_T}(K).
\]
\textbf{À quoi ça sert ?}
Cette identité relie directement la smile aux probabilités risque-neutres : un skew ``put'' reflète une masse de probabilité plus importante dans les scénarios de baisse.

\subsection{Arbitrage en maturité : monotonie de la valeur temps}
À strike fixé, un call européen ne doit pas décroître avec la maturité (hors effets de taux négatifs et dividendes particuliers).
Sur la surface en volatilité, une condition pratique (pas toujours suffisante) est :
\[
w(K,T) \triangleq \sigma_{\mathrm{imp}}(K,T)^2\,T \quad \text{(variance totale)} \quad \text{doit être croissante en } T.
\]
Intuition : plus de temps $\Rightarrow$ plus d'incertitude cumulée (variance totale).

\section{Interpolation et lissage : le minimum viable}
\label{sec:interp}

On a des IV sur une grille (quelques strikes et maturités). Pour pricer des options non cotées, il faut interpoler.
Le piège : une interpolation naïve en $\sigma$ peut créer de l'arbitrage.
Une pratique robuste consiste à travailler en variance totale $w=\sigma^2 T$ :
\begin{itemize}
  \item interpolation (souvent linéaire) en $T$ sur $w$ ;
  \item interpolation en strike (cubic spline, SVI, etc.) avec contraintes de convexité si possible.
\end{itemize}

\begin{example}[Interpolation linéaire en maturité sur la variance totale]
Supposons à un strike donné $K$ :
\[
\sigma_{\mathrm{imp}}(K,T_1)=20\%,\quad T_1=0.5,\qquad
\sigma_{\mathrm{imp}}(K,T_2)=18\%,\quad T_2=1.
\]
Alors $w_1=0.2^2\cdot 0.5=0.02$ et $w_2=0.18^2\cdot 1=0.0324$.
À $T=0.75$, interpolation linéaire :
\[
w(0.75)=w_1 + \frac{0.75-0.5}{1-0.5}(w_2-w_1)=0.02+0.5(0.0124)=0.0262.
\]
Donc $\sigma_{\mathrm{imp}}(K,0.75)=\sqrt{w(0.75)/0.75}\approx \sqrt{0.03493}\approx 18.7\%$.
\end{example}

\begin{jurybox}
\begin{itemize}
\item La surface $\sigma_{\mathrm{imp}}(K,T)$ résume smile/skew/term structure observés.
\item Les quotes peuvent se faire en $K$, log-moneyness ou delta : il faut savoir passer de l'un à l'autre.
\item No-arbitrage (intuition) : en strike, les calls sont décroissants et convexes ; en maturité, la variance totale $w=\sigma^2 T$ doit typiquement croître.
\item Interpoler en variance totale est souvent plus stable qu'interpoler en volatilité.
\end{itemize}
\end{jurybox}

\chapter{Calibration : principe général, métriques d'erreur et choix numériques}
\label{chap:calibration}

\section{Le problème : faire coller un modèle à la surface}
Un modèle (Black--Scholes, Heston, SABR, etc.) dépend de paramètres $\theta$.
La calibration consiste à choisir $\theta$ pour minimiser un écart entre :
\begin{itemize}
  \item les prix/IV du modèle,
  \item et les prix/IV observés sur le marché.
\end{itemize}

\begin{definition}[Problème de calibration (forme générique)]
Soit une grille d'options $(K_i,T_i)_{i=1}^N$ avec prix de marché $P_i^{\mathrm{mkt}}$.
On cherche
\[
\theta^\star \in \arg\min_{\theta\in\Theta} \sum_{i=1}^N w_i\,\ell\!\left(P_i^{\mathrm{model}}(\theta)-P_i^{\mathrm{mkt}}\right),
\]
où $\ell$ est une fonction de perte (souvent $\ell(x)=x^2$) et $w_i$ des pondérations.
\end{definition}

\textbf{À quoi ça sert ?}
Sans calibration, un modèle est un exercice académique.
Avec calibration, il devient un outil opérationnel :
\begin{itemize}
  \item pricer des options illiquides (interpolation/extrapolation cohérente) ;
  \item extraire des paramètres interprétables (skew via $\rho$, vol-of-vol via $\xi$, etc.) ;
  \item générer des scénarios et mesurer des risques (VaR/CVaR) ;
  \item construire des stratégies (hedging, relative value).
\end{itemize}

\section{Prix vs volatilité implicite : que minimiser ?}
Deux choix usuels :

\subsection{Calibration en prix}
\[
\min_\theta \sum_i w_i\bigl(P_i^{\mathrm{model}}(\theta)-P_i^{\mathrm{mkt}}\bigr)^2.
\]
\textbf{Avantage :} cohérent économiquement (euros).
\textbf{Inconvénient :} les options ATM valent beaucoup plus cher que les OTM $\Rightarrow$ elles dominent la loss.

\subsection{Calibration en volatilité implicite}
On convertit prix en IV et on minimise
\[
\min_\theta \sum_i w_i\bigl(\sigma_i^{\mathrm{model}}(\theta)-\sigma_i^{\mathrm{mkt}}\bigr)^2.
\]
\textbf{Avantage :} met toutes les options sur une échelle comparable.
\textbf{Inconvénient :} nécessite une inversion (voir \cref{chap:implied_vol}) et devient instable quand vega est faible.

\subsection{Pondération par vega : une justification simple}
Un argument classique : un petit écart de volatilité correspond à un écart de prix
\[
\Delta P \approx \vega \cdot \Delta \sigma.
\]
Donc minimiser des erreurs de volatilité revient à minimiser des erreurs de prix \emph{pondérées} :
\[
(\Delta \sigma)^2 \approx \frac{(\Delta P)^2}{\vega^2}.
\]
Une pondération naturelle est donc $w_i \propto \vega_i^2$ (ou l'inverse selon convention) pour équilibrer.

\section{Contraintes, bornes et identifiabilité}
\label{sec:constraints}

Les paramètres ont souvent des contraintes structurelles :
\begin{itemize}
  \item volatilités et variances positives ;
  \item intensité de sauts $\lambda\ge 0$ ;
  \item corrélation $\rho\in[-1,1]$ ;
  \item conditions de positivité (ex. Feller en Heston).
\end{itemize}

\textbf{Identifiabilité :} certains paramètres sont difficiles à estimer sur certaines grilles.
Exemple : à très court terme, la vol-of-vol est peu informée ; à long terme, l'intensité de jumps peut être confondue avec la vol stochastique.
C'est un point crucial : une calibration ``parfaite'' en loss peut produire des paramètres non robustes.

\section{Optimisation : méthodes et compromis}
\subsection{Gradient vs derivative-free}
\begin{itemize}
\item \textbf{Gradient-based} (Newton, quasi-Newton L-BFGS) : rapide si la loss est lisse et si on peut calculer des dérivées (ou approx par différences finies).
\item \textbf{Derivative-free} (Nelder--Mead, CMA-ES, differential evolution) : robuste sur surfaces irrégulières mais plus coûteux.
\end{itemize}

\subsection{Multi-start et minima locaux}
Les calibrations non linéaires ont souvent plusieurs minima locaux.
Une pratique ``sérieuse'' :
\begin{itemize}
  \item lancer l'optimisation depuis plusieurs points initiaux (multi-start) ;
  \item imposer des bornes réalistes ;
  \item choisir le meilleur minimum au sens économique (et pas seulement la loss).
\end{itemize}

\subsection{Complexité}
Chaque évaluation de $P_i^{\mathrm{model}}(\theta)$ peut coûter cher (intégrales, FFT, Monte Carlo).
La calibration est donc souvent dominée par \emph{le coût du pricing}.
C'est exactement la motivation des surrogates (chapitre sur réseaux de neurones).

\begin{example}[Calibration impossible : une seule $\sigma$ pour deux smiles]
Supposons deux prix observés à la même maturité :
\[
\sigma_{\mathrm{imp}}(K_1,T)=15\%,\qquad \sigma_{\mathrm{imp}}(K_2,T)=25\%.
\]
Le modèle Black--Scholes avec une $\sigma$ constante ne peut pas matcher simultanément les deux : il produira une unique $\sigma$.
La calibration donnera un compromis (``best fit'') qui reflète la limite du modèle.
\end{example}

\begin{jurybox}
\begin{itemize}
\item La calibration est une optimisation : choisir $\theta$ pour que le modèle reproduise les prix/IV de marché.
\item Choix crucial : minimiser en prix (échelle euros) ou en IV (échelle comparable mais inversion instable).
\item Les pondérations (souvent via vega) ne sont pas du détail : elles déterminent \emph{ce que} la calibration cherche à bien reproduire.
\item Les contraintes et l'identifiabilité conditionnent la robustesse : une calibration ``qui fit'' peut être économiquement absurde si non contrainte.
\end{itemize}
\end{jurybox}

\chapter{Optimisation numérique pour la calibration : stabilité, régularisation, contraintes}
\label{chap:opt_calib}

La calibration est, en pratique, un problème d'optimisation non linéaire sous contraintes.
Sur le papier, "minimiser une erreur" semble simple ; en réalité, les difficultés viennent de :
\begin{itemize}
  \item \textbf{l'ill-conditioning} (Vega faible, paramètres peu identifiables) ;
  \item \textbf{les contraintes} (positivité, bornes, stabilité du modèle) ;
  \item \textbf{le bruit numérique} (intégrales oscillantes, FFT, discrétisation) ;
  \item \textbf{les minima locaux}.
\end{itemize}
Ce chapitre explicite les outils standards pour rendre une calibration reproductible.

\section{Du problème général au cas "moindres carrés"}

\subsection{Formulation}
Écrivons une erreur vecteur $r(\theta)\in\R^m$ (les \emph{résidus}) telle que
\[
r_i(\theta)=\begin{cases}
P_i^{\mathrm{model}}(\theta)-P_i^{\mathrm{mkt}} & \text{si l'on calibre en prix},\\
\sigma_i^{\mathrm{model}}(\theta)-\sigma_i^{\mathrm{mkt}} & \text{si l'on calibre en IV}.
\end{cases}
\]
Les pondérations se mettent naturellement dans une matrice diagonale $W=\mathrm{diag}(w_1,\dots,w_m)$.
Le problème devient un \textbf{moindres carrés pondéré} :
\[
\min_{\theta\in\Theta}\ \Phi(\theta)\ \triangleq\ \tfrac12\,\norm{W\,r(\theta)}^2
=\tfrac12\sum_{i=1}^m w_i^2\,r_i(\theta)^2.
\]

\subsection{Gradient et matrice Jacobienne}
Notons $J(\theta)\in\R^{m\times p}$ la Jacobienne des résidus :
\[(J(\theta))_{i,j}=\pd{r_i}{\theta_j}(\theta).
\]
Par la règle de dérivation d'une norme quadratique :
\[
\nabla\Phi(\theta)=J(\theta)^\top W^\top W\, r(\theta).
\]
Le Hessien exact contient des dérivées secondes. L'approximation "Gauss--Newton" consiste à négliger ces termes :
\[
\nabla^2\Phi(\theta)\approx J(\theta)^\top W^\top W\,J(\theta).
\]
\textbf{Intuition :} quand les résidus sont petits (bonne calibration), les dérivées secondes pèsent peu.

\section{Levenberg--Marquardt : un compromis robuste}

\subsection{Le pas de Gauss--Newton}
Autour de $\theta$, on linéarise $r(\theta+\Delta)\approx r(\theta)+J\Delta$.
Alors
\[\Phi(\theta+\Delta)\approx \tfrac12\norm{W(r+J\Delta)}^2.
\]
Minimiser ce quadratique en $\Delta$ donne le système normal :
\[
(J^\top W^\top WJ)\Delta = -J^\top W^\top W r.
\]
Si $J^\top W^\top WJ$ est mal conditionnée (fréquent !), le pas devient instable.

\subsection{Régularisation de Levenberg}
Levenberg--Marquardt ajoute un terme $\lambda I$ :
\[
\bigl(J^\top W^\top WJ + \lambda I\bigr)\Delta = -J^\top W^\top Wr.
\]
\begin{itemize}
  \item si $\lambda\to 0$, on retrouve Gauss--Newton (rapide près du minimum) ;
  \item si $\lambda\to +\infty$, on se rapproche d'une descente de gradient (plus stable, mais lente).
\end{itemize}

\subsection{Choix adaptatif de $\lambda$}
L'heuristique standard :
\begin{enumerate}[label=\alph*)]
  \item proposer $\Delta$ avec un certain $\lambda$ ;
  \item si $\Phi(\theta+\Delta)$ diminue suffisamment, accepter et \emph{réduire} $\lambda$ ;
  \item sinon, refuser et \emph{augmenter} $\lambda$.
\end{enumerate}
C'est ce mécanisme qui rend LM très robuste en calibration.

\begin{lstlisting}[caption={Levenberg--Marquardt (schéma pédagogique)},label={lst:lm}]
input: theta0, lambda0>0
theta <- theta0 ; lambda <- lambda0
repeat:
    r <- residuals(theta)
    J <- jacobian(theta)
    solve: (J' W' W J + lambda I) Delta = - J' W' W r
    if Phi(theta+Delta) < Phi(theta):
        theta <- theta + Delta
        lambda <- lambda / 2
    else:
        lambda <- 2 * lambda
until convergence
output: theta
\end{lstlisting}

\section{Contraintes : rendre l'optimisation "légale"}

\subsection{Bornes et transformations}
Beaucoup de paramètres doivent rester dans un domaine : $\kappa>0$, $\theta>0$, $\xi>0$, $\rho\in(-1,1)$.
Deux stratégies :
\begin{itemize}
  \item \textbf{optimisation bornée} (L-BFGS-B, trust-region) ;
  \item \textbf{reparamétrisation} sans contrainte.
\end{itemize}
Exemples de reparamétrisation :
\[
\kappa = \ee^{a},\quad \theta=\ee^{b},\quad \xi=\ee^{c},\quad \rho=\tanh(d).
\]
Alors on optimise sur $(a,b,c,d)\in\R^4$.

\subsection{Exemple : condition de Feller (Heston)}
Dans Heston, la variance suit un processus de type CIR. Pour éviter que $v_t$ touche trop souvent 0,
on impose souvent (pas toujours strictement nécessaire numériquement) la \textbf{condition de Feller} :
\[
2\kappa\theta \ge \xi^2.
\]
Cette contrainte couple les paramètres et rend l'optimisation plus difficile.
Une manière simple de l'imposer via reparamétrisation est, par exemple,
\[2\kappa\theta = \xi^2 + \ee^{u}\quad (u\in\R),
\]
ce qui garantit automatiquement l'inégalité.

\section{Identifiabilité et échelles : pourquoi "tout fitter" est impossible}

\subsection{Un paramètre "faible" est invisible}
Si une direction de paramètres change très peu les prix, le problème est quasi-plat :
\(J^\top J\) a une petite valeur propre, donc le pas d'optimisation explose.
Typiquement :
\begin{itemize}
  \item la corrélation $\rho$ est surtout identifiée par le skew court terme ;
  \item $\kappa$ et $\theta$ peuvent se compenser sur des maturités limitées ;
  \item $\xi$ est mal identifié si la surface est presque plate.
\end{itemize}
\textbf{Conséquence pratique :} calibrer sur trop peu de maturités rend certains paramètres non identifiables.

\subsection{Choix d'échelle : prix vs IV}
En calibrant en prix, les options ITM (plus chères) dominent la loss.
En calibrant en IV, on compare des quantités homogènes, mais on introduit le problème d'inversion (vega faible).
Une alternative utile est de calibrer en \emph{prix normalisés} :
\[
r_i(\theta)=\frac{P_i^{\mathrm{model}}(\theta)-P_i^{\mathrm{mkt}}}{\max(\text{tick},\,\vega_i\,\sigma_i)}.
\]
Ce type de normalisation limite l'impact disproportionné d'une option peu informative.

\section{Diagnostics : savoir quand une calibration est mauvaise}

\subsection{Erreurs structurées}
Une calibration "mauvaise" n'est pas forcément une grosse MSE. Ce qui compte est la \emph{structure} des erreurs :
\begin{itemize}
  \item si toutes les maturités sont sous-pricées en même temps : erreur de taux/dividendes ;
  \item si le skew est mal reproduit : corrélation/vol-of-vol mal identifiées ;
  \item si seules les ailes sont mauvaises : modèle trop léger en queues (sauts/roughness manquants).
\end{itemize}
Un diagnostic utile consiste à tracer $r_i$ en fonction de $k=\log(K/F)$ et $T$.

\subsection{Sensibilités de paramètres}
Le Jacobien $J$ est aussi une information économique :
\(\pd{\sigma_{\mathrm{imp}}}{\rho}<0\) sur indices (skew put) est attendu.
Si la calibration produit des sensibilités incohérentes, c'est souvent un signe d'instabilité numérique.

\begin{jurybox}
\begin{itemize}
\item La calibration est un moindres carrés non linéaire : résidus $r(\theta)$, Jacobienne $J(\theta)$.
\item Gauss--Newton est rapide mais instable si $J^\top J$ est mal conditionnée.
\item Levenberg--Marquardt stabilise via $(J^\top J+\lambda I)$ avec $\lambda$ adaptatif.
\item Les contraintes (positivité, bornes, Feller, $\rho\in(-1,1)$) se gèrent par optimisation bornée ou reparamétrisation.
\item Une ``bonne" calibration n'est pas juste une petite loss : les erreurs doivent être économiquement plausibles.
\end{itemize}
\end{jurybox}

\chapter{Modèle de Heston : volatilité stochastique, pricing par fonction caractéristique, FFT}
\label{chap:heston}

\section{Pourquoi Heston ?}
La surface de volatilité implicite des indices actions présente souvent un skew marqué : les puts OTM sont chers (IV élevée),
ce qui traduit (i) une probabilité risque-neutre de baisse plus importante, (ii) une aversion au risque de crash,
et (iii) une dynamique où la volatilité augmente quand le spot baisse (\emph{leverage effect}).
Le modèle de Heston est l'un des plus utilisés car :
\begin{itemize}
  \item il introduit une volatilité \emph{stochastique} (variance aléatoire) ;
  \item il permet de générer un skew via une corrélation négative spot/variance ;
  \item il reste suffisamment tractable : pricing européen via fonction caractéristique (semi-fermé).
\end{itemize}

\section{Définition du modèle (sous mesure risque-neutre)}
\label{sec:heston_sde}
\begin{definition}[Modèle de Heston]
Sous la mesure risque-neutre $\Q$, on modélise le spot et la variance instantanée par :
\begin{align*}
\dd S_t &= (r-q)S_t\,\dd t + \sqrt{v_t}\,S_t\,\dd W_t^{(1)},\\
\dd v_t &= \kappa(\theta - v_t)\,\dd t + \xi \sqrt{v_t}\,\dd W_t^{(2)},\\
\dd\langle W^{(1)},W^{(2)}\rangle_t &= \rho\,\dd t,
\end{align*}
où $\kappa>0$ (vitesse de retour à la moyenne), $\theta>0$ (variance de long terme),
$\xi>0$ (vol-of-vol) et $\rho\in[-1,1]$ (corrélation).
\end{definition}

\subsection{Interprétation financière des paramètres}
\begin{itemize}
  \item $\theta$ : niveau ``normal'' de variance à long terme (donc vol de long terme $\sqrt{\theta}$).
  \item $\kappa$ : vitesse de mean-reversion ; plus $\kappa$ est grand, plus $v_t$ revient vite vers $\theta$.
  \item $\xi$ : intensité des chocs sur la variance (plus $\xi$ est grand, plus la variance bouge vite).
  \item $\rho$ : corrélation spot/variance.
  Sur indices, on observe souvent $\rho<0$ : quand $S$ baisse, $v$ monte, ce qui génère un skew put (IV plus élevée pour strikes bas).
\end{itemize}

\subsection{Condition de Feller (positivité de la variance)}
\begin{proposition}[Condition de Feller]
Pour le processus de Cox--Ingersoll--Ross (CIR) ci-dessus, une condition suffisante pour éviter que $v_t$ touche zéro est
\[
2\kappa\theta \ge \xi^2.
\]
\end{proposition}

\textbf{Intuition :}
$\kappa\theta$ contrôle la force de rappel vers un niveau positif,
tandis que $\xi$ contrôle la force des chocs qui peuvent pousser $v_t$ vers zéro.
Quand $\xi$ est trop grand, le bruit domine et la variance peut heurter 0.
En pratique, on peut calibrer même si Feller n'est pas strictement vérifiée, mais cela peut rendre certaines intégrales plus instables.

\section{Passage à la log-dynamique}
On pose $X_t=\log S_t$.
Par Itô (comme dans \cref{sec:examples_ito}) :
\[
\dd X_t = \left((r-q)-\frac12 v_t\right)\dd t + \sqrt{v_t}\,\dd W_t^{(1)}.
\]
Le couple $(X_t,v_t)$ est un processus \emph{affine} : c'est cette structure qui rend la fonction caractéristique calculable.

\section{Fonction caractéristique : dérivation via l'ansatz affine}
\label{sec:heston_cf}

On veut pricer des options européennes. Une voie efficace est de calculer la fonction caractéristique de $X_T$ :
\[
\phi(u) = \E^\Q\left[\ee^{\ii u X_T}\mid X_0=\log S_0,\ v_0\right].
\]
Dans un modèle affine, on cherche une solution de la forme
\[
\phi(u) = \exp\Bigl(A(T;u) + B(T;u)\,v_0 + \ii u \log S_0\Bigr),
\]
où $A$ et $B$ sont à déterminer.

\subsection{Étape 1 : équation de Kolmogorov backward}
Considérons la fonction
\[
f(t,x,v)=\E^\Q\left[\ee^{\ii u X_T}\mid X_t=x,\ v_t=v\right].
\]
Elle vérifie une PDE backward (Feynman--Kac) :
\[
\pd{f}{t} + \mathcal{L} f = 0,\qquad f(T,x,v)=\ee^{\ii u x},
\]
où $\mathcal{L}$ est le générateur de $(X_t,v_t)$.

\subsection{Étape 2 : générateur du couple $(X,v)$}
À partir des SDE :
\[
\dd X = \left((r-q)-\frac12 v\right)\dd t + \sqrt{v}\,\dd W^{(1)},\qquad
\dd v = \kappa(\theta-v)\dd t + \xi\sqrt{v}\,\dd W^{(2)},
\]
et $\dd\langle W^{(1)},W^{(2)}\rangle=\rho \dd t$.
Le générateur est (résultat standard de calcul d'Itô en 2D) :
\[
\mathcal{L}f
=
\left((r-q)-\frac12 v\right)\pd{f}{x}
+ \kappa(\theta-v)\pd{f}{v}
+ \frac12 v\,\pdd{f}{x}
+ \rho \xi v\,\pd{^2 f}{x\,\partial v}
+ \frac12 \xi^2 v\,\pdd{f}{v}.
\]
Les trois termes de second ordre viennent des variations quadratiques :
$(\dd X)^2=v\,\dd t$, $(\dd v)^2=\xi^2 v\,\dd t$, et $\dd X\,\dd v = \rho\xi v\,\dd t$.

\subsection{Étape 3 : insérer l'ansatz et obtenir des ODE}
On pose
\[
f(t,x,v)=\exp\Bigl(A(\tau)+B(\tau)v+\ii u x\Bigr),\qquad \tau=T-t.
\]
On calcule les dérivées :
\[
\pd{f}{x}=\ii u f,\quad \pdd{f}{x}=-(u^2)f,\quad
\pd{f}{v}=B f,\quad \pdd{f}{v}=B^2 f,\quad
\pd{^2 f}{x\,\partial v}=\ii u B f.
\]
On calcule aussi $\pd{f}{t} = -A'(\tau)f - B'(\tau) v f$ (car $\tau=T-t$).

On remplace dans la PDE $\pd{f}{t} + \mathcal{L} f=0$, puis on divise par $f$.
On obtient une identité de la forme
\[
0 = \text{termes constants en } v \ +\ \text{termes proportionnels à } v.
\]
Comme cela doit être vrai pour tout $v$, on identifie séparément :
\begin{align*}
A'(\tau) &= \kappa\theta\,B(\tau) + (r-q)\ii u,\\
B'(\tau) &= \frac12(u^2+\ii u) + \bigl(\rho\xi\ii u-\kappa\bigr)B(\tau) + \frac12\xi^2 B(\tau)^2.
\end{align*}
La seconde équation est une équation de Riccati.
Avec conditions terminales $A(0)=0$ et $B(0)=0$ (puisque $f(T,x,v)=\ee^{\ii u x}$),
on peut résoudre explicitement et obtenir la formule fermée de $\phi(u)$.

\subsection{Forme fermée (résultat standard)}
On définit
\[
d = \sqrt{(\rho\xi\ii u-\kappa)^2 + \xi^2(u^2+\ii u)},
\qquad
g = \frac{\kappa-\rho\xi\ii u-d}{\kappa-\rho\xi\ii u+d}.
\]
Alors une écriture usuelle est
\begin{align*}
A(T;u) &= (r-q)\ii u T + \frac{\kappa\theta}{\xi^2}\left((\kappa-\rho\xi\ii u-d)T - 2\log\left(\frac{1-g\ee^{-dT}}{1-g}\right)\right),\\
B(T;u) &= \frac{\kappa-\rho\xi\ii u-d}{\xi^2}\cdot \frac{1-\ee^{-dT}}{1-g\ee^{-dT}}.
\end{align*}
Donc
\[
\phi(u)=\exp\Bigl(A(T;u)+B(T;u)v_0+\ii u\log S_0\Bigr).
\]

\begin{remark}[Stabilité numérique]
Il existe plusieurs écritures équivalentes de $\phi(u)$.
Certaines peuvent souffrir de problèmes de branche de la racine carrée complexe ou de logarithme (``Heston trap'').
En pratique numérique, on choisit une implémentation stable (``Little Heston trap'' ou variantes).
Pour un cours M1, l'essentiel est de comprendre : \emph{la structure affine mène à une solution exponentielle avec une Riccati}.
\end{remark}

\section{Pricing européen via inversion de Fourier}
\label{sec:heston_pricing}

À partir de $\phi(u)$, on peut pricer un call européen de strike $K$.
Une représentation classique est :
\[
C_0 = S_0\ee^{-qT}P_1 - K\ee^{-rT}P_2,
\]
où $P_1$ et $P_2$ sont des ``probabilités'' calculées par intégrales de Fourier.
L'idée est proche de la récupération d'une CDF à partir d'une fonction caractéristique :
pour une variable $X$, on peut reconstruire $\Prob(X\le x)$ via une intégrale impliquant $\E[\ee^{\ii u X}]$.

Sans entrer dans tous les détails techniques, une forme standard est :
\[
P_2 = \frac12 + \frac{1}{\pi}\int_0^\infty \Re\left(\frac{\ee^{-\ii u\log K}\phi(u)}{\ii u}\right)\dd u,
\]
et $P_1$ est obtenu avec un décalage complexe (lié au terme $S_T$ dans le payoff).

Le point important : le pricing devient une \emph{intégrale 1D}.

\section{Carr--Madan : FFT pour pricer beaucoup de strikes à la fois}
\label{sec:carr_madan}

Quand on calibre, on doit pricer des centaines d'options.
Faire une intégrale par strike peut être trop lent.
Carr--Madan propose une astuce :
\begin{itemize}
\item on ``amortit'' le call pour qu'il soit intégrable en strike ;
\item on prend la transformée de Fourier en log-strike ;
\item on utilise une FFT pour obtenir une grille de prix en $O(N\log N)$.
\end{itemize}

\subsection{Idée principale}
On note $k=\log K$ et $C(k)=C(K)$.
Le call n'est pas intégrable en $k$ (il croît comme $S$ pour $k\to -\infty$).
On choisit $\alpha>0$ et on définit le call amorti
\[
c_\alpha(k)=\ee^{\alpha k}C(k).
\]
Sous des conditions raisonnables, $c_\alpha$ devient intégrable.
On considère sa transformée de Fourier
\[
\hat{c}_\alpha(u)=\int_{-\infty}^{\infty} \ee^{\ii u k}c_\alpha(k)\,\dd k.
\]
Carr--Madan montre qu'on peut exprimer $\hat{c}_\alpha(u)$ directement à partir de la fonction caractéristique $\phi$.
Puis, par inversion (et FFT), on récupère $C(k)$ sur une grille de $k$.

\section{Illustration numérique : skew typique généré par Heston}
Pour fixer les idées, voici une smile obtenue sous un jeu de paramètres raisonnable
($S_0=100$, $r=2\%$, $q=0$, $T=1$ an,
$v_0=\theta=0.04$,
$\kappa=1.5$, $\xi=0.3$, $\rho=-0.7$).
On calcule d'abord les prix Heston par intégrale de Fourier, puis on les reconvertit en volatilités implicites BS.

\begin{figure}[h!]
\centering
\begin{tikzpicture}
\begin{axis}[
  width=0.85\textwidth,
  height=0.52\textwidth,
  grid=both,
  xlabel={Strike $K$},
  ylabel={Volatilité implicite $\sigma_{\mathrm{imp}}$},
  ymin=0.12, ymax=0.30,
  xmin=58, xmax=142
]
\addplot+ coordinates { (60,0.267275) (65,0.257064) (70,0.247170) (75,0.237541) (80,0.228138) (85,0.218938) (90,0.209937) (95,0.201154) (100,0.192637) (105,0.184468) (110,0.176768) (115,0.169687) (120,0.163392) (125,0.158027) (130,0.153673) (135,0.150323) (140,0.147893) };
\end{axis}
\end{tikzpicture}
\caption{Smile de volatilité implicite générée par un modèle de Heston (exemple numérique).}
\label{fig:heston_smile_example}
\end{figure}

On observe un skew : les strikes bas (puts OTM) ont une IV plus élevée.
Dans Heston, ce skew est principalement piloté par $\rho$ (corrélation spot/variance) et $\xi$ (vol-of-vol).

\section{Calibration Heston : objectifs et pièges}
En pratique, on calibre $\theta=(\kappa,\theta,\xi,\rho,v_0)$ (parfois on fixe $v_0=\theta$ ou on ajoute un paramètre de risque).
Les pièges classiques :
\begin{itemize}
  \item \textbf{Minima locaux} : multi-start quasi obligatoire.
  \item \textbf{Trade-offs} : $\kappa$ et $\theta$ peuvent compenser $\xi$ sur certaines grilles.
  \item \textbf{Données} : bid/ask et erreurs d'IV (surtout extrêmes strikes) peuvent dominer.
  \item \textbf{Instabilités} : intégrales oscillantes, choix du cutoff, paramètre d'amortissement en FFT.
\end{itemize}

\begin{jurybox}
\begin{itemize}
\item Heston ajoute une variance stochastique $v_t$ de type CIR : mean-reversion ($\kappa,\theta$), vol-of-vol ($\xi$), corrélation ($\rho$).
\item La structure affine permet une fonction caractéristique explicite via un ansatz exponentiel et une Riccati.
\item Les prix européens se calculent par intégrales de Fourier ; Carr--Madan + FFT accélère le pricing multi-strikes.
\item Le skew actions vient surtout de $\rho<0$ ; la calibration est une optimisation non convexe avec contraintes.
\end{itemize}
\end{jurybox}

\chapter{Diffusions à sauts : Merton et Bates}
\label{chap:jumps}

\section{Pourquoi ajouter des sauts ?}
Même avec volatilité stochastique, certains phénomènes sont difficiles à reproduire :
\begin{itemize}
  \item \textbf{gaps} de marché (annonces, overnight) ;
  \item \textbf{smile court terme} très prononcé : à très courte maturité, la diffusion seule donne une densité trop ``fine'' ;
  \item \textbf{queues épaisses} : plus d'événements extrêmes que la gaussienne.
\end{itemize}
Les modèles à sauts ajoutent une composante rare mais significative.

\section{Processus de Poisson composé (rappel)}
\begin{definition}[Poisson et sauts multiplicatifs]
Soit $(N_t)$ un processus de Poisson d'intensité $\lambda>0$ : $\Prob(N_t=n)=\ee^{-\lambda t}(\lambda t)^n/n!$.
On considère des tailles de saut i.i.d. $(J_k)$, et le processus de saut multiplicatif
\[
\prod_{k=1}^{N_t} J_k.
\]
Dans le modèle de Merton, on pose souvent $J_k=\ee^{Y_k}$ avec $Y_k\sim \Normal(\mu_J,\delta_J^2)$.
\end{definition}

\section{Modèle de Merton (jump diffusion) sous $\Q$}
\label{sec:merton_sde}
\begin{definition}[Merton jump diffusion]
Sous $\Q$ :
\[
\frac{\dd S_t}{S_{t^-}} = (r-q-\lambda \kappa_J)\dd t + \sigma \dd W_t + (J-1)\dd N_t,
\]
où $J=\ee^Y$ est le multiplicateur de saut et
\[
\kappa_J \triangleq \E[J-1] = \E[\ee^Y-1] = \ee^{\mu_J+\frac12\delta_J^2}-1.
\]
\end{definition}

\textbf{Pourquoi le terme $-\lambda\kappa_J$ ?}
Parce qu'en moyenne, les sauts ajoutent une dérive :
en un petit $\dd t$, on a $\Prob(\dd N_t=1)\approx \lambda \dd t$ et sinon 0,
donc l'espérance instantanée du terme de saut est $\lambda \E[J-1]\dd t = \lambda \kappa_J \dd t$.
On le soustrait pour conserver le drift risque-neutre $(r-q)$ pour $S_t$.

\section{Pricing : mélange de Black--Scholes (intuition et formule)}
Conditionnellement au nombre de sauts $N_T=n$, la somme des logs de sauts est gaussienne (somme de gaussiennes),
donc $\log S_T$ reste gaussien conditionnellement à $n$.
On obtient donc un mélange de lois lognormales, et un mélange de prix Black--Scholes.

\begin{proposition}[Prix de call Merton comme somme de BS]
Soit $C_{\mathrm{BS}}(S_0,K,r,q,T,\sigma_n)$ un call BS avec une volatilité effective $\sigma_n$.
Alors
\[
C_0^{\mathrm{Merton}}
=
\sum_{n=0}^{\infty} \frac{\ee^{-\lambda T}(\lambda T)^n}{n!}\;
C_{\mathrm{BS}}\Bigl(S_0,K,r,q,T,\ \sigma_n\Bigr),
\]
où $\sigma_n^2 = \sigma^2 + \frac{n\delta_J^2}{T}$ et où le forward est ajusté par la moyenne des sauts (effet de $\mu_J$).
\end{proposition}

\textbf{Intuition :} chaque terme correspond au scénario ``il y a eu $n$ sauts''.
En pratique, la somme est tronquée à un $n_{\max}$ raisonnable (car la probabilité de beaucoup de sauts est faible).

\section{Fonction caractéristique (utile pour FFT)}
Le modèle de Merton est aussi très commode en Fourier.
La fonction caractéristique de $\log S_T$ factorise en :
\[
\phi(u)=\phi_{\text{diffusion}}(u)\cdot \phi_{\text{sauts}}(u),
\]
avec
\[
\phi_{\text{sauts}}(u)=\exp\left(\lambda T\bigl(\E[\ee^{\ii u Y}]-1\bigr)\right).
\]
Pour $Y\sim \Normal(\mu_J,\delta_J^2)$,
\[
\E[\ee^{\ii u Y}] = \exp\left(\ii u\mu_J - \frac12 u^2\delta_J^2\right).
\]
Donc
\[
\phi_{\text{sauts}}(u)=\exp\left(\lambda T\left(\exp\left(\ii u\mu_J - \frac12 u^2\delta_J^2\right)-1\right)\right).
\]
Cela rend Carr--Madan (FFT) très efficace.

\section{Bates : Heston + sauts}
\begin{definition}[Modèle de Bates (idée)]
Le modèle de Bates combine :
\begin{itemize}
  \item une variance stochastique à la Heston ;
  \item des sauts multiplicatifs à la Merton.
\end{itemize}
On garde la dynamique de $v_t$ de Heston et on ajoute un terme $(J-1)\dd N_t$ dans $\dd S_t/S_{t^-}$,
avec correction de drift $-\lambda\kappa_J$.
\end{definition}

\textbf{Pourquoi c'est utile ?}
Heston explique bien le skew ``structurel'' via $\rho$.
Les sauts expliquent souvent une partie du smile court terme et des queues extrêmes.
Bates combine les deux.

\begin{example}[Impact qualitatif sur la smile courte maturité]
À maturité très courte, un modèle diffusif produit une distribution de $\log S_T$ très concentrée (variance $\propto T$).
Ajouter des sauts injecte une probabilité non négligeable de mouvements de grande amplitude même pour $T$ petit,
ce qui augmente la valeur des options très OTM : la smile court terme remonte.
\end{example}

\section{Limites et extensions}
\begin{itemize}
  \item Les sauts de Merton sont gaussiens en log : cela peut être trop simple.
  \item L'intensité $\lambda$ peut être stochastique.
  \item Les sauts peuvent être asymétriques (ex. double exponentielle de Kou).
\end{itemize}

\begin{jurybox}
\begin{itemize}
\item Les sauts capturent gaps et queues épaisses ; ils sont cruciaux pour la smile très court terme.
\item Merton : diffusion + Poisson composé, avec correction de drift $-\lambda\kappa_J$ pour rester risque-neutre.
\item Le prix européen est un mélange de Black--Scholes (conditionnel au nombre de sauts) ou via fonction caractéristique/FFT.
\item Bates = Heston + sauts : skew + queues.
\end{itemize}
\end{jurybox}

\chapter{SABR : dynamique, formule de Hagan et calibration}
\label{chap:sabr}

\section{Pourquoi SABR ?}
En taux et en FX, les marchés aiment :
\begin{itemize}
  \item une paramétrisation \emph{simple} par maturité (smile par ``tenor'') ;
  \item un modèle flexible qui reproduit des smiles très variées ;
  \item une formule d'approximation rapide de la volatilité implicite.
\end{itemize}
SABR (Stochastic Alpha Beta Rho) répond à ces besoins.

\section{Dynamique SABR (sur le forward)}
\begin{definition}[Modèle SABR]
On modélise un forward $F_t$ (par exemple un taux forward ou un taux swap) et sa volatilité instantanée $\alpha_t$ :
\begin{align*}
\dd F_t &= \alpha_t F_t^\beta\,\dd W_t^{(1)},\\
\dd \alpha_t &= \nu \alpha_t\,\dd W_t^{(2)},\\
\dd\langle W^{(1)},W^{(2)}\rangle_t &= \rho\,\dd t,
\end{align*}
avec paramètres $\beta\in[0,1]$, $\nu>0$ et $\rho\in[-1,1]$.
\end{definition}

\subsection{Interprétation des paramètres}
\begin{itemize}
  \item $\beta$ : élasticité de la volatilité au niveau du forward.
  \begin{itemize}
    \item $\beta=1$ : modèle lognormal (type Black) : $\dd F/F = \alpha\,\dd W$.
    \item $\beta=0$ : modèle normal (type Bachelier) : $\dd F = \alpha\,\dd W$.
  \end{itemize}
  \item $\alpha_0$ : niveau initial de vol (``ATM level'').
  \item $\nu$ : vol-of-vol (à quel point $\alpha_t$ bouge).
  \item $\rho$ : corrélation entre mouvements de $F$ et de la vol $\alpha$ (contrôle l'asymétrie/skew).
\end{itemize}

\section{Objectif : une formule rapide d'IV (Hagan)}
SABR n'admet pas de formule fermée simple pour les prix d'options européennes (sauf cas particuliers).
Hagan et al. proposent une approximation asymptotique directe pour la volatilité implicite,
très utilisée en pratique pour calibrer et interpoler des smiles.

\section{Formule de Hagan (Black vol) : expression et lecture}
\label{sec:hagan_formula}

On note $K$ le strike, $T$ la maturité, $F$ le forward initial $F_0$ et $\alpha$ le niveau initial $\alpha_0$.
Définissons :
\[
z = \frac{\nu}{\alpha}\,(F K)^{\frac{1-\beta}{2}}\,\log\left(\frac{F}{K}\right),
\qquad
\chi(z)=\log\left(\frac{\sqrt{1-2\rho z+z^2}+z-\rho}{1-\rho}\right).
\]
Alors l'approximation de Hagan pour la volatilité implicite lognormale (Black) est :
\[
\sigma_{\mathrm{Hagan}}(F,K)
=
\frac{\nu\,\log(F/K)}{\chi(z)}\cdot
\frac{\alpha}{(F K)^{\frac{1-\beta}{2}}}
\cdot
\left[1 + \text{correction}(T)\right],
\]
où la correction en temps (ordre $T$) vaut :
\[
\text{correction}(T)
=
\left(
\frac{(1-\beta)^2}{24}\frac{\alpha^2}{(F K)^{1-\beta}}
+
\frac{\rho\beta\nu\alpha}{4(F K)^{\frac{1-\beta}{2}}}
+
\frac{2-3\rho^2}{24}\nu^2
\right)T.
\]

\textbf{Lecture financière :}
\begin{itemize}
  \item Le facteur $\alpha/(FK)^{(1-\beta)/2}$ donne le niveau de vol en fonction du niveau de taux (effet CEV).
  \item Le ratio $\nu\log(F/K)/\chi(z)$ encode l'asymétrie et la courbure de la smile via $\rho$ et $\nu$.
  \item La correction en $T$ ajoute l'effet du temps sur la smile.
\end{itemize}

\begin{remark}[Pourquoi $\chi(z)$ ?]
Le terme $\chi(z)$ apparaît en résolvant (à l'ordre asymptotique) une équation de diffusion
et en effectuant un changement de variable qui ``linearise'' la dynamique.
Le ratio $\frac{z}{\chi(z)}$ est proche de 1 quand $z$ est petit, et corrige le comportement hors ATM.
\end{remark}

\section{Cas ATM : limite $K\to F$}
Quand $K\to F$, on a $\log(F/K)\to 0$ et donc $z\to 0$.
On utilise l'expansion $\chi(z)\sim z$ pour obtenir la formule ATM :

\[
\sigma_{\mathrm{ATM}}
=
\frac{\alpha}{F^{1-\beta}}
\left[
1 +
\left(
\frac{(1-\beta)^2}{24}\frac{\alpha^2}{F^{2-2\beta}}
+
\frac{\rho\beta\nu\alpha}{4F^{1-\beta}}
+
\frac{2-3\rho^2}{24}\nu^2
\right)T
\right].
\]

\textbf{À quoi ça sert ?}
L'ATM donne une contrainte forte sur $\alpha$ (niveau), tandis que les strikes OTM contraignent surtout $\rho$ et $\nu$ (skew et convexité).

\section{Calibration SABR (pratique)}
\label{sec:sabr_calib}

En pratique, on calibre SABR \emph{maturité par maturité} :
\begin{itemize}
  \item on fixe souvent $\beta$ (ex. 0.5 ou 0.0 selon conventions et régime de taux) ;
  \item on ajuste $(\alpha,\rho,\nu)$ pour fitter 3 quotes (ATM, 25$\Delta$ put, 25$\Delta$ call) ou plus.
\end{itemize}

\subsection{Pourquoi fixer $\beta$ ?}
Parce que $\beta$ est difficile à identifier avec peu de points, et parce qu'il contrôle une dimension ``structurelle''
du marché (relation niveau-vol). Fixer $\beta$ stabilise fortement la calibration.

\begin{example}[Mini-procédure de calibration (idée)]
Supposons qu'on dispose, à une maturité $T$, de trois IV :
\[
\sigma_{\mathrm{ATM}},\quad \sigma_{25\Delta\ \mathrm{put}},\quad \sigma_{25\Delta\ \mathrm{call}}.
\]
On fixe $\beta$, puis on cherche $(\alpha,\rho,\nu)$ qui minimisent l'erreur quadratique entre ces trois valeurs et la formule de Hagan.
On peut initialiser $\alpha$ par la formule ATM, puis ajuster $\rho$ (skew) et $\nu$ (courbure).
\end{example}

\section{Limites et extensions}
\begin{itemize}
  \item Hagan est une approximation : elle peut devenir mauvaise pour très longues maturités ou paramètres extrêmes.
  \item Quand les taux sont proches de zéro ou négatifs, la formulation lognormale (Black) peut poser problème :
  on préfère parfois une volatilité normale (Bachelier) ou des adaptations.
  \item Des extensions : SABR déplacé (shifted), ou modèles multi-facteurs.
\end{itemize}

\begin{jurybox}
\begin{itemize}
\item SABR modélise un forward $F_t$ avec une vol stochastique $\alpha_t$ lognormale ; $\beta$ contrôle la dépendance niveau-vol.
\item La formule de Hagan donne une approximation directe de la volatilité implicite : rapide pour la calibration.
\item $\alpha$ fixe le niveau (ATM), $\rho$ le skew, $\nu$ la courbure (vol-of-vol).
\item Limites : approximation, régime de taux bas/négatif, stabilité hors ATM.
\end{itemize}
\end{jurybox}

\chapter{Volatilité ``rough'' et modèles de type Volterra}
\label{chap:rough}

\section{Motivation empirique : la volatilité est très irrégulière}
Depuis les années 2010, de nombreuses études empiriques suggèrent que la volatilité réalisée
ou certaines proxies (log-vol) ont une régularité plus faible que celle d'un Brownien.
Une manière de quantifier cela est via l'exposant de Hurst $H$ :
\begin{itemize}
  \item $H=1/2$ : Brownien classique (rugosité ``standard'').
  \item $H<1/2$ : trajectoires plus rugueuses, mémoire négative à court terme (\emph{rough}).
\end{itemize}
En pratique, on estime souvent $H\simeq 0.1$ pour la volatilité.

\section{Brownien fractionnaire : définition et intuition}
\begin{definition}[Brownien fractionnaire (fBM)]
Un Brownien fractionnaire $(W_t^H)_{t\ge 0}$ d'exposant $H\in(0,1)$ est un processus gaussien centré
tel que
\[
\E[W_t^H W_s^H] = \frac12\left(t^{2H}+s^{2H}-|t-s|^{2H}\right).
\]
\end{definition}

\textbf{Propriétés importantes :}
\begin{itemize}
  \item $W^H$ a des incréments stationnaires, mais \emph{pas indépendants} si $H\neq 1/2$.
  \item La variance de l'incrément est $\Var(W_t^H-W_s^H)=|t-s|^{2H}$.
  Donc les incréments sont d'ordre $|t-s|^H$.
\end{itemize}

\begin{remark}[Pourquoi Itô ``ne marche pas'' directement]
Pour $H\neq 1/2$, $W^H$ n'est pas une semi-martingale.
Le calcul d'Itô (basé sur la variation quadratique finie et sur les semi-martingales) ne s'applique pas tel quel.
On doit utiliser d'autres outils (intégrales de Young, Malliavin, rough paths), trop avancés pour ce cours.
\end{remark}

\section{Représentation Volterra : ramener du rough à du Brownien}
Une idée clé (très utilisée) est de représenter $W^H$ comme une intégrale contre un Brownien standard :
\[
W_t^H = \int_0^t K_H(t,s)\,\dd W_s,
\]
où $K_H(t,s)$ est un \emph{noyau Volterra} (triangulaire : $s<t$) qui encode la mémoire.
Pour le rough case $H<1/2$, un comportement typique est
\[
K_H(t,s) \approx (t-s)^{H-\frac12}
\quad \text{quand } s\uparrow t,
\]
ce qui amplifie les effets à très court terme et rend les trajectoires plus rugueuses.

\section{Exemple : modèle rBergomi (esprit)}
\begin{definition}[rBergomi (forme simplifiée)]
Un prototype de modèle rough est :
\begin{align*}
\dd S_t &= (r-q)S_t\,\dd t + \sqrt{v_t}\,S_t\,\dd W_t^{(1)},\\
v_t &= \xi_0(t)\exp\left(\eta W_t^H - \frac12 \eta^2 t^{2H}\right),
\end{align*}
où $W_t^H$ est un fBM construit à partir d'un Brownien standard et $\eta$ contrôle l'amplitude.
\end{definition}

\textbf{Intuition :}
\begin{itemize}
  \item $\xi_0(t)$ est la \emph{courbe de variance forward} (niveau moyen de variance par maturité) ;
  \item la partie exponentielle rend $v_t$ lognormal (positif) ;
  \item $H<1/2$ rend les chocs de volatilité très rugueux, ce qui reproduit mieux certaines smiles court terme.
\end{itemize}

\section{Numérique : simulation et approximation markovienne}
\label{sec:rough_numerics}

Le pricing rough est souvent fait par Monte Carlo, car pas de CF simple.
Le défi : simuler $W_t^H$ efficacement.

\subsection{Simulation directe (coût élevé)}
Si l'on simule un vecteur gaussien corrélé de taille $n$ (grille de temps), le coût naïf est $O(n^2)$.
Pour de la calibration, c'est trop lent.

\subsection{Approximation markovienne (somme d'exponentielles)}
Une astuce importante :
approximer le noyau $K_H(t,s)$ par une somme d'exponentielles :
\[
K_H(t,s)\approx \sum_{j=1}^m w_j \ee^{-\lambda_j(t-s)}.
\]
Pourquoi c'est utile ?
Parce que chaque terme exponentiel peut être généré par une équation markovienne auxiliaire (type OU),
ce qui transforme le modèle rough en un modèle multi-facteur markovien (dimension $m$), plus facile à simuler et parfois calibrer.

\begin{example}[Idée d'état auxiliaire]
Définissez des processus
\[
Y_t^{(j)} = \int_0^t \ee^{-\lambda_j(t-s)}\,\dd W_s.
\]
Alors $Y_t^{(j)}$ satisfait
\[
\dd Y_t^{(j)} = -\lambda_j Y_t^{(j)}\,\dd t + \dd W_t,
\]
ce qui est markovien.
On peut reconstruire $W_t^H$ comme combinaison linéaire des $Y^{(j)}$.
\end{example}

\section{Limites et ce qu'il faut retenir}
\begin{itemize}
  \item Les modèles rough améliorent le fit du smile court terme, mais sont plus difficiles à calibrer et simuler.
  \item Pour un praticien M1, l'essentiel est de comprendre la motivation (rugosité empirique) et les conséquences (Itô non applicable directement, simulation plus coûteuse).
\end{itemize}

\begin{jurybox}
\begin{itemize}
\item Rough volatility : la volatilité semble avoir une régularité $H<1/2$ (très rugueuse) à court terme.
\item Le Brownien fractionnaire $W^H$ n'est pas une semi-martingale si $H\neq 1/2$ : Itô classique ne s'applique pas tel quel.
\item Une représentation Volterra $W_t^H=\int_0^t K(t,s)\dd W_s$ permet de travailler avec un Brownien standard.
\item Numériquement, l'approximation markovienne (somme d'exponentielles) est une stratégie clé pour rendre le problème tractable.
\end{itemize}
\end{jurybox}

\chapter{Réseaux de neurones et calibration accélérée : surrogates et inversion}
\label{chap:nn_calibration}

\section{Pourquoi des réseaux en calibration ?}
Dans les modèles riches (Heston, Bates, rough), pricer une option peut nécessiter :
\begin{itemize}
  \item une intégrale oscillante (Fourier) ;
  \item une FFT (grille) ;
  \item ou du Monte Carlo (souvent coûteux).
\end{itemize}
La calibration (chapitre \cref{chap:calibration}) requiert des centaines/milliers d'évaluations.
Même si un pricing unitaire prend 5 ms, la calibration peut prendre plusieurs secondes ou minutes.

Un réseau de neurones peut servir de :
\begin{itemize}
  \item \textbf{surrogate direct} : $\theta \mapsto \sigma_{\mathrm{imp}}(K,T)$ (ou prix) ;
  \item \textbf{surrogate inverse} : surface $\mapsto \theta$ (approximation d'inversion) ;
  \item \textbf{initialiseur} : fournir un bon point de départ pour une calibration classique.
\end{itemize}

\section{Rappel minimal : perceptron multicouche (MLP)}
\begin{definition}[Couche affine + non-linéarité]
Une couche prend une entrée $x\in\R^d$ et produit
\[
h = \varphi(Wx+b),
\]
où $W$ est une matrice, $b$ un biais, et $\varphi$ une non-linéarité (ReLU, tanh, etc.).
Un réseau MLP empile plusieurs couches.
\end{definition}

\subsection{Forward pass}
Si le réseau a $L$ couches, on calcule
\[
h^{(0)}=x,\qquad
h^{(\ell)}=\varphi_\ell\bigl(W_\ell h^{(\ell-1)} + b_\ell\bigr),\quad \ell=1,\dots,L,
\]
et la sortie $\hat{y}=h^{(L)}$.

\section{Apprentissage supervisé : loss et objectif}
\label{sec:supervised}

On dispose d'un dataset $\{(x_i,y_i)\}_{i=1}^n$.
Dans notre contexte :
\begin{itemize}
  \item $x_i$ = paramètres du modèle + caractéristiques du contrat (ex. $K,T,S_0,r,q$),
  \item $y_i$ = prix ou volatilité implicite.
\end{itemize}

\begin{definition}[Loss quadratique (MSE)]
\[
\mathcal{L}(\theta_{\mathrm{NN}})=\frac{1}{n}\sum_{i=1}^n \norm{\hat{y}_i - y_i}^2,
\]
où $\theta_{\mathrm{NN}}$ désigne les poids/biais du réseau.
\end{definition}

\textbf{À quoi ça sert ?}
Minimiser $\mathcal{L}$ revient à approximer la fonction inconnue $x\mapsto y$.
Si le réseau généralise, on obtient un pricing quasi-instantané.

\section{Backprop : la règle de la chaîne (idée)}
Le réseau est une composition de fonctions.
Pour une couche : $h=\varphi(z)$ avec $z=Wx+b$.
La dérivée de la loss par rapport aux paramètres suit la règle de la chaîne :
\[
\pd{\mathcal{L}}{W} = \pd{\mathcal{L}}{z}\,\pd{z}{W}.
\]
La backprop calcule ces gradients efficacement en remontant de la sortie vers l'entrée.

\begin{example}[Mini-exemple (une couche)]
Supposons $\hat{y}=w x$ (scalaire) et $\mathcal{L}=\frac12(\hat{y}-y)^2$.
Alors
\[
\pd{\mathcal{L}}{w} = (\hat{y}-y)\cdot x.
\]
C'est exactement un cas de règle de la chaîne : dérivée de la loss times dérivée de la sortie.
\end{example}

\section{Génération de dataset synthétique (la clé en quant)}
\label{sec:synthetic_data}

Pour pricer, on sait générer des données de façon \emph{contrôlée} :
\begin{enumerate}[label=\arabic*)]
  \item tirer des paramètres $\theta$ dans un domaine réaliste (ex. Heston) ;
  \item tirer des contrats $(K,T)$ sur une grille ;
  \item calculer des prix/IV via le modèle ``exact'' (FFT/intégrale/MC) ;
  \item entraîner le réseau sur ces couples.
\end{enumerate}

\textbf{Attention :} si le domaine de paramètres est trop large, le réseau aura du mal.
Si le domaine est trop étroit, il ne généralisera pas hors de cette zone.
C'est un vrai choix de design.

\section{Surrogate direct vs inverse : avantages et limites}
\subsection{Surrogate direct}
\[
(\theta,K,T)\mapsto \sigma_{\mathrm{imp}}.
\]
\textbf{Avantage :} bien posé (fonction bien définie), facile à entraîner.
\textbf{Usage :} accélérer la calibration classique (on minimise une loss mais avec le réseau au lieu du pricing exact).

\subsection{Surrogate inverse}
\[
\{\sigma_{\mathrm{imp}}(K_i,T_i)\}_{i=1}^N \mapsto \theta.
\]
\textbf{Problème :} l'inversion n'est pas toujours unique.
Deux jeux de paramètres peuvent générer des smiles proches sur une grille limitée.
Cela crée une ambiguïté : le réseau peut ``moyenner'' des solutions.

Solutions possibles :
\begin{itemize}
  \item prédire une distribution (incertitude) plutôt qu'un point ;
  \item imposer des régularisations et contraintes ;
  \item utiliser l'inverse NN seulement comme initialisation, puis raffiner par calibration exacte.
\end{itemize}

\section{Overfitting et validation}
Un réseau peut apprendre ``par cœur'' le dataset synthétique sans généraliser.
Bonnes pratiques :
\begin{itemize}
  \item split train/validation/test ;
  \item early stopping ;
  \item normalisation des inputs ;
  \item contrôle de la complexité (nombre de couches/neuronnes).
\end{itemize}

\begin{example}[Toy surrogate pour Black--Scholes]
On peut entraîner un réseau à approximer $C_{\mathrm{BS}}(S,K,T,r,q,\sigma)$.
Même si la formule est fermée, l'exercice montre :
\begin{itemize}
  \item comment choisir des entrées (features) ;
  \item comment normaliser (ex. log-moneyness) ;
  \item comment évaluer l'erreur relative/absolue.
\end{itemize}
\end{example}

\begin{jurybox}
\begin{itemize}
\item Les réseaux servent surtout à accélérer la calibration : remplacer un pricing coûteux par une approximation rapide.
\item Surrogate direct ($\theta\to$ prix/IV) est mieux posé que l'inverse (surface $\to \theta$), souvent non-unique.
\item Dataset synthétique : c'est un avantage majeur en quant, mais le domaine de paramètres doit être choisi intelligemment.
\item Risque principal : sur-généralisation ou sous-généralisation ; la validation est obligatoire.
\end{itemize}
\end{jurybox}

\part{Couverture discrète \& reinforcement learning}

\chapter{Hedging discret : PnL, coûts de transaction et mesures de risque}
\label{chap:discrete_hedging}

\section{Pourquoi le hedging discret est un vrai sujet ?}
En théorie Black--Scholes, on hedge en continu et sans coûts : la réplication est parfaite.
En pratique :
\begin{itemize}
  \item on hedge à des fréquences finies (minute, heure, jour) ;
  \item chaque trade coûte (bid/ask, commissions, impact) ;
  \item la volatilité et la dynamique sont mal spécifiées.
\end{itemize}
Donc même un ``delta hedging'' produit un \emph{PnL aléatoire}.
La question devient : \textbf{quelle stratégie minimise le risque (ou maximise l'utilité) compte tenu des coûts ?}

\section{Définition du PnL d'une stratégie de couverture}
\label{sec:pnl_definition}

Supposons qu'on \emph{vende} une option européenne de payoff $g(S_T)$ (donc on est short).
On reçoit $V_0$ au départ. On hedge avec le sous-jacent et un compte cash.

On fixe des dates $0=t_0<t_1<\dots<t_N=T$.
Notons :
\begin{itemize}
  \item $S_i=S_{t_i}$ ;
  \item $\Delta_i$ la position en sous-jacent détenue sur $(t_i,t_{i+1}]$ ;
  \item $B_i$ le montant en cash au temps $t_i$ (placé au taux $r$ entre les dates).
\end{itemize}

La valeur du portefeuille de hedge au temps $t_i$ est :
\[
\Pi_i = \Delta_i S_i + B_i.
\]

\subsection{Auto-financement sans coûts}
Sans coûts, la stratégie est auto-financée si, lors du rebalancing,
l'achat/vente de sous-jacent est financé par le cash :
\[
B_{i+1} = B_i \ee^{r\Delta t_i} - (\Delta_{i+1}-\Delta_i)S_{i+1},
\qquad \Delta t_i=t_{i+1}-t_i.
\]
Au temps final, on liquide le sous-jacent et on paie le payoff.
Le PnL (profit net) pour le vendeur d'option est :
\[
\PnL = V_0\ee^{rT} + \sum_{i=0}^{N-1} \Delta_i (S_{i+1}-S_i)\ee^{r(T-t_{i+1})} - g(S_T),
\]
où on a exprimé les gains de trading en valeur à $T$ (actualisation inverse).

\subsection{Avec coûts de transaction proportionnels}
Un modèle simple : chaque variation de position coûte une fraction $c$ du notionnel :
\[
\text{coût}_i = c\,S_{i+1}\,\abs{\Delta_{i+1}-\Delta_i}.
\]
Alors l'équation de cash devient
\[
B_{i+1} = B_i \ee^{r\Delta t_i} - (\Delta_{i+1}-\Delta_i)S_{i+1} - c\,S_{i+1}\,\abs{\Delta_{i+1}-\Delta_i}.
\]
Ces coûts rendent le rebalancing fréquent potentiellement destructeur : on veut un compromis \emph{risque vs coûts}.

\section{Erreur de couverture et intuition}
Même si $\Delta_i$ est le delta BS évalué au temps $t_i$,
le hedge n'est plus parfait car :
\begin{itemize}
  \item la delta est recalculée à des dates discrètes (on manque la convexité entre deux dates) ;
  \item la volatilité implicite utilisée peut bouger ;
  \item des coûts s'accumulent.
\end{itemize}

Une approximation classique (sous BS, sans coûts) est que l'erreur de couverture est liée à $\Gamma$ :
\[
\PnL \approx \frac12 \sum_{i=0}^{N-1} \Gamma_{t_i} (S_{i+1}-S_i)^2.
\]
En moyenne, $(S_{i+1}-S_i)^2 \approx \sigma^2 S_i^2 \Delta t_i$,
ce qui suggère que la variance du PnL décroît avec le pas de temps, mais pas forcément linéairement.

\section{Mesures de risque pour évaluer une stratégie}
\label{sec:risk_measures_hedge}

\subsection{Variance / MSE}
On peut minimiser
\[
\min \Var(\PnL) \quad \text{ou} \quad \min \E[\PnL^2].
\]
C'est naturel si on veut ``stabiliser'' le résultat, mais ça traite symétriquement gains et pertes.

\subsection{Value-at-Risk et CVaR}
\begin{definition}[VaR et CVaR (niveau $\alpha$)]
Pour une perte $L=-\PnL$ :
\[
\VaR_\alpha(L) = \inf\{ \ell : \Prob(L\le \ell)\ge \alpha\}.
\]
La CVaR (aussi appelée Expected Shortfall) est
\[
\CVaR_\alpha(L)=\E[L\mid L\ge \VaR_\alpha(L)].
\]
\end{definition}

\textbf{Intuition :}
\begin{itemize}
  \item VaR : seuil de perte à ne pas dépasser avec probabilité $\alpha$.
  \item CVaR : perte moyenne dans le pire $(1-\alpha)\%$ des cas (plus sensible aux queues).
\end{itemize}

\subsection{Utilité exponentielle (approche économique)}
Une autre approche : maximiser l'espérance d'utilité,
par exemple utilité exponentielle $U(x)=-\exp(-\gamma x)$,
ce qui pénalise fortement les grandes pertes.

\begin{example}[Comparaison qualitative]
Une stratégie qui gagne souvent un peu mais perd rarement beaucoup peut avoir :
\begin{itemize}
  \item une variance faible,
  \item mais une CVaR très mauvaise.
\end{itemize}
En gestion des risques, la CVaR est souvent plus pertinente car elle regarde explicitement la queue de distribution.
\end{example}

\section{Ce qui change quand on introduit une stratégie ``apprenante''}
Le delta hedging est une règle analytique dérivée d'un modèle.
Mais si le modèle est imparfait et si on ajoute des coûts,
on peut vouloir \emph{apprendre} une politique $\Delta_i=\pi(\text{état}_i)$ qui :
\begin{itemize}
  \item réduit les coûts (moins de rebalancing),
  \item contrôle une mesure de risque (variance, CVaR),
  \item et reste raisonnable économiquement.
\end{itemize}
C'est précisément là qu'intervient la formulation MDP et le RL (chapitres suivants).

\begin{jurybox}
\begin{itemize}
\item Le hedging discret transforme un problème ``déterministe'' (réplication) en un problème de \emph{contrôle sous incertitude} : le PnL devient aléatoire.
\item Les coûts de transaction rendent la fréquence de rebalancing centrale : plus souvent = moins de risque diffusif, mais plus de coûts.
\item Les objectifs pertinents ne se limitent pas à la variance : VaR/CVaR ou utilité peuvent mieux refléter la contrainte du risque extrême.
\item Une stratégie apprenante cherche une politique $\Delta=\pi(\text{état})$ optimale pour un critère donné.
\end{itemize}
\end{jurybox}

\chapter{MDP pour la couverture et algorithmes RL tabulaires}
\label{chap:mdp_rl}

\section{Pourquoi formuler le hedging comme un MDP ?}
Une stratégie de hedging discret prend des décisions séquentielles :
à chaque date, on observe un état du marché et on choisit une action (nouvelle delta).
Cela ressemble exactement à un problème de contrôle en temps discret.

\section{Définition d'un MDP (Markov Decision Process)}
\begin{definition}[MDP]
Un MDP est un quintuplet $(\mathcal{S},\mathcal{A},P,R,\gamma)$ où :
\begin{itemize}
  \item $\mathcal{S}$ : ensemble des états ;
  \item $\mathcal{A}$ : ensemble des actions ;
  \item $P(s'|s,a)$ : probabilité de transition ;
  \item $R(s,a,s')$ : récompense (ou coût) ;
  \item $\gamma\in(0,1]$ : facteur d'actualisation (en finance de hedging court, on peut prendre $\gamma=1$).
\end{itemize}
\end{definition}

\subsection{État, action et reward pour le hedging}
Un choix typique :
\begin{itemize}
  \item état $s_t$ : $(t, S_t, v_t, \Delta_t,\ldots)$ ou des features comme moneyness, temps restant, greeks.
  \item action $a_t$ : nouvelle position $\Delta_{t^+}$ (ou ajustement $\Delta_{t^+}-\Delta_t$).
  \item reward : négatif des coûts et du risque.
  Par exemple :
  \[
  r_t = -\text{coûts}_t - \lambda \cdot (\text{risque}_t),
  \]
  où $\lambda$ règle le compromis.
\end{itemize}

\section{Valeurs : $V^\pi$ et $Q^\pi$}
\begin{definition}[Fonctions valeur]
Pour une politique $\pi(a|s)$ :
\[
V^\pi(s)=\E^\pi\left[\sum_{t=0}^{T} \gamma^t r_t \,\middle|\, s_0=s\right],
\qquad
Q^\pi(s,a)=\E^\pi\left[\sum_{t=0}^{T} \gamma^t r_t \,\middle|\, s_0=s, a_0=a\right].
\]
\end{definition}

\textbf{À quoi ça sert ?}
\begin{itemize}
  \item $V^\pi$ mesure la performance future d'une politique.
  \item $Q^\pi$ permet de choisir l'action optimale : $\pi^\star(s)=\arg\max_a Q^\star(s,a)$.
\end{itemize}

\section{Équations de Bellman}
\label{sec:bellman}

\begin{theorem}[Bellman (policy evaluation)]
Pour une politique $\pi$ :
\[
Q^\pi(s,a)=\E\left[r + \gamma \E_{a'\sim\pi(\cdot|s')}[Q^\pi(s',a')] \,\middle|\, s,a\right].
\]
\end{theorem}

\begin{theorem}[Bellman optimalité]
La fonction optimale $Q^\star$ vérifie :
\[
Q^\star(s,a)=\E\left[r + \gamma \max_{a'} Q^\star(s',a') \,\middle|\, s,a\right].
\]
\end{theorem}

\textbf{Intuition :}
le futur optimal, c'est ``récompense immédiate + meilleur futur possible''.

\section{Apprentissage TD : le cœur du RL}
\label{sec:td}

On n'a pas accès à la vraie transition $P$ ni aux vraies espérances.
On observe des trajectoires $(s_t,a_t,r_t,s_{t+1})$ et on met à jour des estimations de $Q$.

\begin{definition}[Erreur TD]
Pour une transition :
\[
\delta_t = r_t + \gamma \widehat{V}(s_{t+1}) - \widehat{V}(s_t).
\]
Pour $Q$ :
\[
\delta_t = r_t + \gamma \widehat{Q}(s_{t+1},a_{t+1}) - \widehat{Q}(s_t,a_t).
\]
\end{definition}

\subsection{SARSA (on-policy)}
SARSA met à jour avec l'action réellement prise au temps suivant :
\[
\boxed{
Q(s_t,a_t)\leftarrow Q(s_t,a_t) + \alpha\Bigl[r_t + \gamma Q(s_{t+1},a_{t+1}) - Q(s_t,a_t)\Bigr].
}
\]
C'est \emph{on-policy} : on apprend la valeur de la politique effectivement suivie (avec exploration).

\subsection{Expected SARSA}
On remplace $Q(s_{t+1},a_{t+1})$ par son espérance sous la politique :
\[
Q(s_t,a_t)\leftarrow Q(s_t,a_t) + \alpha\Bigl[r_t + \gamma \E_{a'\sim\pi}[Q(s_{t+1},a')] - Q(s_t,a_t)\Bigr].
\]
Moins de variance, parfois plus stable.

\subsection{Q-learning (off-policy)}
Q-learning vise directement $Q^\star$ :
\[
\boxed{
Q(s_t,a_t)\leftarrow Q(s_t,a_t) + \alpha\Bigl[r_t + \gamma \max_{a'}Q(s_{t+1},a') - Q(s_t,a_t)\Bigr].
}
\]
C'est \emph{off-policy} : on peut explorer avec une politique $\epsilon$-greedy
tout en évaluant le greedy optimal.

\section{On-policy vs off-policy : pourquoi c'est important ?}
\begin{itemize}
  \item On-policy : plus stable si l'environnement dépend fortement de la politique (rare en hedging, mais possible via impact).
  \item Off-policy : plus efficace en données (on peut réutiliser des trajectoires générées par d'autres politiques).
\end{itemize}

\section{Importance sampling : corriger le biais off-policy}
\label{sec:is}

Supposons qu'on a des données générées sous une politique de comportement $b(a|s)$,
mais qu'on veut évaluer une autre politique $\pi(a|s)$.
On introduit des poids
\[
w_t = \frac{\pi(a_t|s_t)}{b(a_t|s_t)}.
\]
Intuition : si une action est rare sous $b$ mais probable sous $\pi$, il faut augmenter son poids.

\begin{remark}[Variance]
Les ratios peuvent exploser si $b$ met très peu de masse là où $\pi$ en met beaucoup,
ce qui crée une variance énorme.
En RL pratique, c'est une difficulté majeure : on stabilise via clipping, régularisation, ou politiques proches.
\end{remark}

\begin{example}[Mini-exemple de mise à jour Q-learning]
Supposons $\gamma=1$, $\alpha=0.1$.
On observe $(s,a,r,s')$ avec $r=-0.5$ (coût), et $\max_{a'}Q(s',a')=2.0$.
Si $Q(s,a)=1.0$ alors
\[
Q(s,a)\leftarrow 1.0 + 0.1(-0.5 + 2.0 - 1.0)=1.0 + 0.1(0.5)=1.05.
\]
Le Q augmente car l'état futur promet un bon reward, malgré un coût immédiat.
\end{example}

\begin{jurybox}
\begin{itemize}
\item Un MDP formalise le hedging discret : état (marché + position), action (ajuster la hedge), reward (PnL/coûts/risque).
\item Bellman relie $Q$ d'aujourd'hui à ``reward + futur'' ; TD learning approxime cela à partir de trajectoires.
\item SARSA (on-policy) apprend la politique suivie ; Q-learning (off-policy) apprend l'optimal greedy.
\item Importance sampling corrige le biais off-policy mais peut exploser en variance.
\end{itemize}
\end{jurybox}

\chapter{DQN : Q-learning avec réseaux (design, stabilité, extensions)}
\label{chap:dqn}

\section{Pourquoi DQN ?}
Les méthodes tabulaires supposent que l'espace d'états $\mathcal{S}$ est petit et discret.
En hedging réel, l'état contient des variables continues (spot, temps, vol, etc.) :
impossible de stocker une table $Q(s,a)$.

DQN (Deep Q-Network) approxime $Q(s,a)$ par un réseau de neurones $Q_\theta(s,a)$.

\section{Objectif : minimiser une loss TD}
Pour Q-learning, l'équation de Bellman optimalité suggère :
\[
Q^\star(s,a) \approx r + \gamma \max_{a'} Q^\star(s',a').
\]
On définit une \emph{cible} (target) :
\[
y = r + \gamma \max_{a'} Q_{\theta^-}(s',a'),
\]
où $\theta^-$ sont des paramètres \emph{gelés} (target network).
Puis on minimise la loss quadratique
\[
\mathcal{L}(\theta)=\E\left[(Q_\theta(s,a)-y)^2\right].
\]
On fait une descente de gradient stochastique sur des minibatches.

\section{Replay buffer : pourquoi on le fait}
\begin{itemize}
  \item Les transitions $(s_t,a_t,r_t,s_{t+1})$ sont corrélées dans le temps.
  \item La descente de gradient suppose des données i.i.d. (idéalement).
\end{itemize}
On stocke donc des transitions dans une mémoire (buffer), puis on échantillonne aléatoirement des minibatches.
Cela réduit la corrélation et améliore la stabilité.

\section{Target network : pourquoi on le fait}
Sans target network, la cible $y$ dépend de $\theta$ lui-même :
on essaie de poursuivre une cible mouvante, ce qui peut diverger.
On introduit donc $\theta^-$ mis à jour lentement (copie périodique ou moyenne exponentielle).
C'est une astuce de stabilisation essentielle.

\section{Exploration : epsilon-greedy}
On choisit l'action :
\[
a=
\begin{cases}
\text{action aléatoire} & \text{avec probabilité } \epsilon,\\
\arg\max_a Q_\theta(s,a) & \text{avec probabilité } 1-\epsilon.
\end{cases}
\]
En hedging, l'exploration doit être contrôlée (on ne veut pas de positions absurdes) :
on borne les actions (delta min/max) et on peut explorer autour d'un delta ``baseline''.

\section{Pseudo-code DQN (version minimale)}
\begin{lstlisting}[caption={DQN (version minimale, pseudo-code).},label={lst:dqn}]
Initialize Q_theta and target Q_theta_minus = Q_theta
Initialize replay buffer D

for episode = 1..E:
    reset environment -> state s
    for t = 1..T:
        with prob epsilon choose random action a
        else a = argmax_a Q_theta(s,a)
        step -> (r, s_next)
        store (s,a,r,s_next) in D
        sample minibatch from D
        y = r + gamma * max_a' Q_theta_minus(s_next, a')
        minimize (Q_theta(s,a) - y)^2 over minibatch
        every K steps: theta_minus <- theta
        s <- s_next
\end{lstlisting}

\section{Stabilité : la ``deadly triad''}
Il y a trois ingrédients qui, combinés, rendent l'apprentissage instable :
\begin{itemize}
  \item \textbf{fonction approximatrice} (réseau de neurones) ;
  \item \textbf{bootstrapping} (cible dépend d'une estimation) ;
  \item \textbf{off-policy} (données générées par une autre politique).
\end{itemize}
DQN contient les trois : d'où la nécessité de replay buffer + target network + autres techniques.

\subsection{Trucs pratiques (niveau M1)}
\begin{itemize}
  \item \textbf{clipping des rewards} ou normalisation : évite des gradients explosifs.
  \item \textbf{Double DQN} : réduit le biais d'optimisme du $\max$ en séparant sélection et évaluation.
  \item \textbf{Dueling architecture} : sépare valeur d'état et avantage d'action (utile si beaucoup d'actions similaires).
\end{itemize}

\section{Actions continues : limites de DQN}
Dans le hedging, l'action naturelle $\Delta$ est continue.
DQN suppose un ensemble discret d'actions.
Solution : discretiser $\Delta$ (ex. 41 niveaux entre $-1$ et $+1$), mais cela :
\begin{itemize}
  \item augmente la taille de l'espace d'actions ;
  \item impose une granularité artificielle.
\end{itemize}
Extensions possibles :
\begin{itemize}
  \item \textbf{DDPG/TD3} : actor-critic pour actions continues.
  \item \textbf{SAC} : ajoute entropie, souvent très stable.
\end{itemize}

\section{Design du reward en hedging (le point non négociable)}
Un reward mal défini donne une stratégie mauvaise mais ``optimale'' pour la mauvaise métrique.
Exemples :
\begin{itemize}
  \item reward = PnL brut : le modèle peut prendre du risque extrême.
  \item reward = $-\PnL^2$ : minimise variance, mais ignore CVaR.
  \item reward = $-\text{coûts} - \lambda \cdot \text{risque}$ : compromis explicite.
\end{itemize}
En pratique, on doit décider : \emph{qu'est-ce qu'on veut optimiser ?} (stabilité, drawdown, tail risk, etc.).

\begin{example}[Reward simple risk-aware]
À chaque pas :
\[
r_t = - c S_t|\Delta_t-\Delta_{t^-}| - \lambda(\Delta_t S_t)^2,
\]
où le premier terme pénalise le turnover (coûts) et le second pénalise une exposition directionnelle trop grande.
Ce n'est pas universel, mais illustre la logique : reward = coût + pénalité de risque.
\end{example}

\begin{jurybox}
\begin{itemize}
\item DQN approxime $Q(s,a)$ par un réseau et apprend via une loss TD $(Q_\theta - y)^2$.
\item Replay buffer (dé-corrélation) et target network (cible stable) sont essentiels pour éviter la divergence.
\item DQN discretise les actions : pour actions continues, on préfère des méthodes actor-critic (DDPG/TD3/SAC).
\item Le design du reward est \emph{le} point critique : on optimise ce qu'on code, pas ce qu'on souhaite implicitement.
\end{itemize}
\end{jurybox}

\part{Courbe des taux : discounting, forwards, calibration}

\chapter{Discount factors, taux spot et forwards : dérivations claires}
\label{chap:yield_basics}

\section{Pourquoi une courbe des taux ?}
Pour pricer des dérivés (options, swaps, forwards), il faut :
\begin{itemize}
  \item actualiser des cashflows futurs ;
  \item parfois projeter des taux forwards (pour coupons flottants) ;
  \item être cohérent avec les instruments liquides du marché.
\end{itemize}
La courbe des taux est l'objet qui relie toutes ces opérations.

\section{Discount factor, taux zéro, taux forward : définitions}
\begin{definition}[Discount factor]
Le discount factor $P(0,T)$ est la valeur aujourd'hui d'1 euro certain à la date $T$ :
\[
P(0,T)=\text{PV}_0(1 \text{ payé à } T).
\]
\end{definition}

\begin{definition}[Taux zéro (spot) continûment composé]
Le taux spot $y(0,T)$ est défini par
\[
P(0,T)=\ee^{-y(0,T)\,T}
\quad\Longleftrightarrow\quad
y(0,T) = -\frac{1}{T}\log P(0,T).
\]
\end{definition}

\begin{definition}[Forward instantané]
Le forward instantané $f(0,T)$ est
\[
f(0,T) = -\pd{}{T}\log P(0,T).
\]
\end{definition}

\subsection{Relations entre $y$ et $f$}
Partant de $\log P(0,T)=-\int_0^T f(0,u)\,\dd u$, on obtient
\[
y(0,T) = \frac{1}{T}\int_0^T f(0,u)\,\dd u.
\]
Donc :
\begin{itemize}
  \item $f$ est la ``densité'' de taux ;
  \item $y$ est une moyenne de $f$ sur $[0,T]$.
\end{itemize}

\section{Forward simple sur une période}
\begin{definition}[Forward simple entre $T_1$ et $T_2$]
Le taux forward simple $F(0;T_1,T_2)$ (avec année fraction $\tau=T_2-T_1$) est défini par
\[
(1+\tau F(0;T_1,T_2)) = \frac{P(0,T_1)}{P(0,T_2)}.
\]
Donc
\[
F(0;T_1,T_2) = \frac{1}{\tau}\left(\frac{P(0,T_1)}{P(0,T_2)}-1\right).
\]
\end{definition}

\textbf{Dérivation (arbitrage) :}
acheter un zéro-coupon $T_2$ et vendre un $T_1$ crée un cashflow nul aujourd'hui et déterministe sur $[T_1,T_2]$,
ce qui définit un taux forward sans arbitrage.

\section{Bootstrapping (idée) : construire $P(0,T)$ à partir d'instruments liquides}
Le bootstrapping consiste à :
\begin{itemize}
  \item prendre des instruments courts (dépôts, OIS) pour obtenir $P(0,T)$ aux petites maturités ;
  \item utiliser des swaps (ou OIS swaps) pour étendre aux maturités plus longues.
\end{itemize}
L'idée essentielle : chaque instrument donne une équation linéaire en discount factors inconnus.
On les résout séquentiellement en avançant en maturité.

\begin{example}[Forward à partir de deux taux zéro]
Supposons $y(0,1)=3\%$ et $y(0,2)=3.5\%$ (continu).
Alors
\[
P(0,1)=\ee^{-0.03}\approx 0.9704,\qquad P(0,2)=\ee^{-0.07}\approx 0.9324.
\]
Le forward simple sur $[1,2]$ ($\tau=1$) est :
\[
F(0;1,2)=\frac{P(0,1)}{P(0,2)}-1 \approx \frac{0.9704}{0.9324}-1\approx 0.0408 = 4.08\%.
\]
\end{example}

\section{Interpolation : point sensible}
On observe $P(0,T)$ ou $y(0,T)$ sur une grille.
Pour pricer, on doit interpoler.
Bonnes pratiques :
\begin{itemize}
  \item interpoler sur $\log P(0,T)$ (qui est souvent plus lisse) ;
  \item éviter des forwards négatifs non voulus ;
  \item utiliser des interpolations monotones/convexes si on veut préserver l'absence d'arbitrage (plus avancé).
\end{itemize}

\begin{jurybox}
\begin{itemize}
\item $P(0,T)$ est l'objet de base : tout le reste (taux spot, forwards) en découle.
\item $y(0,T)=-(1/T)\log P(0,T)$ et $f(0,T)=-\partial_T\log P(0,T)$ : $y$ est une moyenne de $f$.
\item Le forward simple sur $[T_1,T_2]$ vérifie $(1+\tau F)=P(0,T_1)/P(0,T_2)$.
\item Bootstrapping = résoudre séquentiellement des équations d'instruments liquides pour obtenir les discount factors.
\end{itemize}
\end{jurybox}

\chapter{Modèles paramétriques : Nelson--Siegel et Nelson--Siegel--Svensson}
\label{chap:ns_nss}

\section{Pourquoi des modèles paramétriques ?}
Une courbe bootstrapée donne des points $(T_i,P(0,T_i))$ ou $(T_i,y(0,T_i))$.
Pour :
\begin{itemize}
  \item obtenir une courbe lisse (et des forwards stables),
  \item extrapoler raisonnablement au-delà des points,
  \item résumer la forme de la courbe par quelques facteurs,
\end{itemize}
on utilise des modèles paramétriques.

\section{Nelson--Siegel (NS) : formule et interprétation}
\begin{definition}[Nelson--Siegel (sur taux zéro)]
Le taux zéro continûment composé est modélisé par
\[
y(T)=\beta_0 + \beta_1 \frac{1-\ee^{-T/\tau}}{T/\tau}
+ \beta_2\left(\frac{1-\ee^{-T/\tau}}{T/\tau}-\ee^{-T/\tau}\right),
\]
avec paramètres $(\beta_0,\beta_1,\beta_2,\tau)$ et $T>0$.
\end{definition}

\subsection{Lecture ``level / slope / curvature''}
Définissons les fonctions de base :
\[
\ell(T)=1,\quad
s(T)=\frac{1-\ee^{-T/\tau}}{T/\tau},\quad
c(T)=\frac{1-\ee^{-T/\tau}}{T/\tau}-\ee^{-T/\tau}.
\]
Alors
\[
y(T)=\beta_0 \ell(T)+\beta_1 s(T)+\beta_2 c(T).
\]
\begin{itemize}
  \item $\beta_0$ (level) : comme $\lim_{T\to\infty} s(T)=0$ et $\lim_{T\to\infty} c(T)=0$,
  on a $\lim_{T\to\infty} y(T)=\beta_0$. Donc $\beta_0$ fixe le niveau long terme.
  \item $\beta_1$ (slope) : $s(0^+)=1$, donc $y(0^+)\approx \beta_0+\beta_1$. $\beta_1$ affecte surtout le court terme.
  \item $\beta_2$ (curvature) : $c(0^+)=0$ et $c(T)$ est en cloche (maximum vers $T\approx \tau$).
  Donc $\beta_2$ ajuste la courbure au moyen terme.
\end{itemize}

\subsection{Rôle de $\tau$}
Le paramètre $\tau$ positionne la zone de courbure :
plus $\tau$ est grand, plus le ``hump'' se déplace vers les maturités longues.

\section{Nelson--Siegel--Svensson (NSS)}
\begin{definition}[NSS]
NSS ajoute un second terme de courbure :
\[
y(T)=\beta_0 + \beta_1 s(T;\tau_1) + \beta_2 c(T;\tau_1) + \beta_3 c(T;\tau_2),
\]
avec deux paramètres d'échelle $\tau_1,\tau_2$.
\end{definition}
Cela permet de fitter des courbes avec deux bosses (rare mais utile en pratique).

\section{Calibration (moindres carrés)}
\label{sec:ns_calib}

On observe des taux $y_i^{\mathrm{mkt}}$ aux maturités $T_i$.
On minimise :
\[
\min_{\beta,\tau} \sum_{i=1}^n w_i\left(y(T_i;\beta,\tau)-y_i^{\mathrm{mkt}}\right)^2.
\]
\textbf{Note :} si $\tau$ est fixé, la calibration en $\beta$ est un problème linéaire (régression),
car $y$ est linéaire en $\beta$.
On peut donc :
\begin{itemize}
  \item grid-search sur $\tau$,
  \item et pour chaque $\tau$, résoudre un least squares en $\beta$.
\end{itemize}
C'est une stratégie très robuste (évite une optimisation non linéaire complète).

\begin{example}[Illustration : forme NS avec paramètres simples]
Prenons (en continûment composé) :
\[
\beta_0=3\%,\quad \beta_1=-2\%,\quad \beta_2=3\%,\quad \tau=1.5.
\]
On obtient une courbe qui démarre sous 3\% (car $\beta_1<0$) puis remonte et se stabilise vers 3\%.
\end{example}

\begin{figure}[h!]
\centering
\begin{tikzpicture}
\begin{axis}[
  width=0.85\textwidth,
  height=0.52\textwidth,
  grid=both,
  xlabel={Maturité $T$ (années)},
  ylabel={Taux zéro $y(T)$},
  xmin=0, xmax=10,
  ymin=0.0, ymax=0.06
]
\addplot+[
  domain=0.05:10,
  samples=200
]
{0.03 + (-0.02)*((1-exp(-x/1.5))/(x/1.5)) + (0.03)*(((1-exp(-x/1.5))/(x/1.5))-exp(-x/1.5))};
\end{axis}
\end{tikzpicture}
\caption{Exemple de courbe Nelson--Siegel (paramètres illustratifs).}
\label{fig:ns_example}
\end{figure}

\section{Limites et usage pratique}
\begin{itemize}
  \item NS/NSS sont empiriques : ce ne sont pas des modèles ``d'absence d'arbitrage'' au sens strict.
  \item Les forwards dérivés peuvent être sensibles (car $f=-\partial_T \log P$ amplifie le bruit).
  \item Malgré cela, NS/NSS sont très utilisés comme \emph{résumé factoriel} et pour des interpolations stables.
\end{itemize}

\begin{jurybox}
\begin{itemize}
\item Nelson--Siegel modélise $y(T)$ avec 4 paramètres : $\beta_0$ (niveau long), $\beta_1$ (pente court), $\beta_2$ (courbure moyen), $\tau$ (échelle).
\item NSS ajoute un second terme de courbure pour plus de flexibilité.
\item Calibration robuste : grid sur $\tau$ puis least squares en $\beta$ (car linéaire en $\beta$ à $\tau$ fixé).
\item Limite : pas un modèle strictement arbitrage-free ; prudence sur les forwards dérivés.
\end{itemize}
\end{jurybox}

\part{Allocation de portefeuille \& gestion du risque (décision d'investissement)}

\chapter{Markowitz : moyenne--variance, min-variance, max Sharpe}
\label{chap:markowitz}

\section{Pourquoi Markowitz ?}
Un desk ou un gérant doit transformer des anticipations (rendements, risques, corrélations)
en allocations $w$.
Markowitz propose un cadre simple :
\begin{itemize}
  \item le rendement est résumé par son espérance ;
  \item le risque est résumé par la variance ;
  \item on optimise un compromis rendement/risque.
\end{itemize}
C'est un modèle de base, mais il structure énormément de raisonnements en gestion.

\section{Cadre : vecteur de rendements et matrice de covariance}
\begin{definition}[Rendements et portefeuille]
Soit $R\in\R^n$ le vecteur de rendements (sur un horizon donné).
On note
\[
\mu=\E[R]\in\R^n,\qquad \Sigma=\Cov(R)\in\R^{n\times n}.
\]
Pour des poids $w\in\R^n$ (somme à 1), le rendement du portefeuille est
\[
R_p = w^\top R,
\]
avec
\[
\E[R_p]=w^\top \mu,\qquad \Var(R_p)=w^\top \Sigma w.
\]
\end{definition}

\textbf{Hypothèses :}
\begin{itemize}
  \item $\mu$ et $\Sigma$ existent et sont stables ;
  \item la variance est une mesure adéquate du risque ;
  \item pas de non-linéarités (options), pas de coûts (dans la version de base).
\end{itemize}

\section{Problème moyenne--variance (forme contrainte)}
\label{sec:mv_constrained}
Une formulation canonique :
\[
\min_{w} \quad w^\top \Sigma w
\quad \text{s.c.}\quad
\begin{cases}
w^\top \mu = m,\\
w^\top \1 = 1.
\end{cases}
\]
On cherche la variance minimale pour un rendement cible $m$.

\subsection{Dérivation par multiplicateurs de Lagrange (détaillée)}
On forme le Lagrangien :
\[
\mathcal{L}(w,\lambda,\nu)= w^\top \Sigma w - \lambda(w^\top \mu - m) - \nu(w^\top \1 - 1).
\]
Condition du premier ordre en $w$ :
\[
\nabla_w \mathcal{L} = 2\Sigma w - \lambda \mu - \nu \1 = 0.
\]
Donc
\[
w = \frac12 \Sigma^{-1}(\lambda \mu + \nu \1).
\]
On peut réabsorber le $\tfrac12$ dans $\lambda,\nu$ et écrire
\[
w = \Sigma^{-1}(a\mu + b\1),
\]
pour certains scalaires $a,b$.

Imposons les contraintes.
Définissons trois scalaires utiles :
\[
A=\1^\top \Sigma^{-1}\1,\qquad
B=\1^\top \Sigma^{-1}\mu,\qquad
C=\mu^\top \Sigma^{-1}\mu.
\]
Alors
\[
\1^\top w = aB + bA = 1,\qquad
\mu^\top w = aC + bB = m.
\]
C'est un système linéaire en $(a,b)$ :
\[
\begin{pmatrix}
B & A\\
C & B
\end{pmatrix}
\begin{pmatrix}
a\\ b
\end{pmatrix}
=
\begin{pmatrix}
1\\ m
\end{pmatrix}.
\]
Son déterminant vaut $D=AC-B^2$ (positif si $\Sigma$ définie positive et $\mu$ non colinéaire à $\1$).
On obtient
\[
a = \frac{Bm-A m}{D},\qquad
b = \frac{C- B m}{D}.
\]
Donc la solution explicite est
\[
\boxed{
w(m)=\Sigma^{-1}\left(\frac{Bm-A m}{D}\mu + \frac{C- B m}{D}\1\right).
}
\]

\subsection{Frontière efficiente}
On peut aussi obtenir la variance minimale associée à $m$ :
\[
\sigma_p^2(m)=w(m)^\top \Sigma w(m)=\frac{A m^2 - 2B m + C}{D}.
\]
C'est une parabole en $(\sigma, m)$ : la frontière efficiente.

\section{Portefeuille de variance minimale globale (GMV)}
Si on enlève la contrainte de rendement (on ne fixe que $w^\top \1=1$),
on obtient :
\[
w_{\mathrm{GMV}} = \frac{\Sigma^{-1}\1}{\1^\top \Sigma^{-1}\1}=\frac{\Sigma^{-1}\1}{A}.
\]
C'est le portefeuille le moins risqué possible (sous la mesure variance) parmi ceux investis à 100\%.

\section{Max Sharpe (portefeuille tangentiel)}
Si on a un actif sans risque de taux $r_f$,
on veut maximiser le ratio de Sharpe
\[
\mathrm{SR}(w)=\frac{w^\top \mu - r_f}{\sqrt{w^\top \Sigma w}}
\quad \text{avec } w^\top \1 = 1.
\]
Le portefeuille tangentiel (sans contrainte long-only) est proportionnel à :
\[
w_{\mathrm{tan}} \propto \Sigma^{-1}(\mu - r_f \1).
\]
Puis on normalise pour que la somme fasse 1.

\section{Contraintes, estimation risk, robustesse}
\label{sec:markowitz_limits}
Les deux gros problèmes en pratique :
\begin{itemize}
  \item \textbf{contraintes} : long-only, limites de poids, exposition sectorielle, etc. $\Rightarrow$ optimisation quadratique (QP) sans solution fermée.
  \item \textbf{estimation} : $\mu$ est très bruité, $\Sigma$ aussi (surtout en grande dimension). Les poids Markowitz deviennent instables et extrêmes.
\end{itemize}
Solutions usuelles :
\begin{itemize}
  \item shrinkage de covariance (Ledoit--Wolf), factor models ;
  \item robust optimization (pénaliser l'incertitude sur $\mu$) ;
  \item privilégier des portefeuilles moins sensibles à $\mu$ (min-variance, risk parity).
\end{itemize}

\begin{example}[Deux actifs : calcul explicite]
Supposons deux actifs avec
\[
\mu=
\begin{pmatrix}
6\%\\ 10\%
\end{pmatrix},
\qquad
\Sigma=
\begin{pmatrix}
0.04 & 0.01\\
0.01 & 0.09
\end{pmatrix}.
\]
On cherche le GMV :
\[
\Sigma^{-1}=
\frac{1}{0.04\cdot 0.09-0.01^2}
\begin{pmatrix}
0.09 & -0.01\\
-0.01 & 0.04
\end{pmatrix}
=
\frac{1}{0.0035}
\begin{pmatrix}
0.09 & -0.01\\
-0.01 & 0.04
\end{pmatrix}.
\]
Alors
\[
\Sigma^{-1}\1 = \frac{1}{0.0035}\begin{pmatrix}0.08\\0.03\end{pmatrix},
\qquad
A=\1^\top \Sigma^{-1}\1 = \frac{0.11}{0.0035}.
\]
Donc
\[
w_{\mathrm{GMV}}=\frac{1}{0.11}\begin{pmatrix}0.08\\0.03\end{pmatrix}\approx \begin{pmatrix}72.7\%\\27.3\%\end{pmatrix}.
\]
\end{example}

\begin{jurybox}
\begin{itemize}
\item Markowitz résume le risque par la variance : $\Var(R_p)=w^\top\Sigma w$.
\item L'optimisation min-variance pour un rendement cible $m$ a une solution fermée via Lagrange ; la frontière efficiente est $\sigma^2(m)=(A m^2-2Bm+C)/D$.
\item GMV : $w_{\mathrm{GMV}}=\Sigma^{-1}\1/(\1^\top\Sigma^{-1}\1)$.
\item En pratique, estimation risk et contraintes dominent : d'où shrinkage, factor models, risk parity.
\end{itemize}
\end{jurybox}

\chapter{Risk parity et risk budgeting : dérivations et intuition}
\label{chap:risk_parity}

\section{Pourquoi risk parity ?}
Markowitz dépend fortement de $\mu$ (rendements attendus), très bruités.
Risk parity propose une alternative :
\begin{itemize}
  \item on ignore (ou on réduit) le rôle de $\mu$ ;
  \item on construit un portefeuille où chaque actif contribue ``également'' au risque.
\end{itemize}
C'est souvent plus stable out-of-sample.

\section{Volatilité de portefeuille et contributions}
Soit $\sigma_p=\sqrt{w^\top \Sigma w}$ la volatilité du portefeuille.
La \emph{contribution marginale} au risque du poids $w_i$ est :
\[
\pd{\sigma_p}{w_i} = \frac{(\Sigma w)_i}{\sigma_p}.
\]
\begin{definition}[Risk contribution]
La contribution au risque de l'actif $i$ est
\[
\mathrm{RC}_i = w_i \pd{\sigma_p}{w_i} = w_i\frac{(\Sigma w)_i}{\sigma_p}.
\]
\end{definition}

\subsection{Somme des contributions}
On a :
\[
\sum_{i=1}^n \mathrm{RC}_i = \frac{1}{\sigma_p}\sum_i w_i(\Sigma w)_i = \frac{w^\top \Sigma w}{\sigma_p}=\sigma_p.
\]
Donc les contributions s'additionnent exactement à la volatilité du portefeuille.

\section{Risk parity (budgets égaux)}
Risk parity ``pur'' impose
\[
\mathrm{RC}_1=\dots=\mathrm{RC}_n=\frac{\sigma_p}{n}.
\]
Plus généralement, on fixe des budgets $b_i>0$ avec $\sum b_i=1$ et on impose :
\[
\mathrm{RC}_i = b_i \sigma_p.
\]
En remplaçant :
\[
w_i(\Sigma w)_i = b_i\, w^\top\Sigma w.
\]
C'est un système non linéaire en $w$.

\section{Cas particulier : actifs non corrélés}
Si $\Sigma$ est diagonale, $\Sigma=\mathrm{diag}(\sigma_1^2,\dots,\sigma_n^2)$,
alors $(\Sigma w)_i=\sigma_i^2 w_i$.
La condition RC égal s'écrit :
\[
w_i\frac{\sigma_i^2 w_i}{\sigma_p}=\frac{\sigma_p}{n}
\quad\Longrightarrow\quad
w_i^2\sigma_i^2 = \frac{\sigma_p^2}{n}.
\]
Donc $w_i\propto 1/\sigma_i$.
\textbf{Conclusion :} en l'absence de corrélations, risk parity = inverse-volatility.

\begin{example}[Deux actifs non corrélés]
Actif 1 : $\sigma_1=10\%$, actif 2 : $\sigma_2=20\%$.
Alors $w_1:w_2 = 1/0.1 : 1/0.2 = 2:1$.
Donc $w_1=2/3$, $w_2=1/3$.
Les deux contributions au risque sont égales (par construction).
\end{example}

\section{Résolution numérique : idées d'algorithmes}
\label{sec:rp_algo}

On peut résoudre risk parity via :
\begin{itemize}
  \item \textbf{optimisation} : minimiser une loss mesurant l'écart aux budgets, par exemple
  \[
  \min_w \sum_i \left(\mathrm{RC}_i - b_i\sigma_p\right)^2
  \quad \text{s.c.}\quad w_i\ge 0,\ \sum w_i=1.
  \]
  \item \textbf{log-weights} : poser $w_i=\exp(x_i)/\sum_j \exp(x_j)$ pour imposer $w_i>0$ et réduire le problème aux $x$.
  \item \textbf{méthodes itératives} : mise à jour multiplicative des poids selon l'écart des contributions.
\end{itemize}

\begin{remark}[Trade-off corrélations]
Quand les actifs sont corrélés, l'inverse-volatility n'est plus vrai.
Un actif très corrélé aux autres peut contribuer beaucoup au risque même avec un petit poids.
Risk parity ajuste implicitement les poids pour répartir le risque de corrélation.
\end{remark}

\section{Limites}
\begin{itemize}
  \item dépend toujours de $\Sigma$ (donc estimation risk sur la covariance) ;
  \item ignore $\mu$ : peut sous-investir des actifs à forte prime de risque ;
  \item peut concentrer dans des actifs très peu volatils mais très corrélés à un facteur latent.
\end{itemize}
D'où l'usage de ``risk budgeting'' par classes d'actifs ou facteurs.

\begin{jurybox}
\begin{itemize}
\item Contribution au risque : $\mathrm{RC}_i = w_i(\Sigma w)_i/\sigma_p$ ; la somme vaut $\sigma_p$.
\item Risk parity : imposer $\mathrm{RC}_i=b_i\sigma_p$ (budgets de risque).
\item Cas non corrélé : weights $\propto 1/\sigma_i$ (inverse-volatility).
\item En pratique : résolution numérique + dépendance à l'estimation de la covariance.
\end{itemize}
\end{jurybox}

\chapter{ACP (PCA), eigen-portfolios et modèles factoriels}
\label{chap:pca}

\section{Pourquoi PCA en finance ?}
Quand on a beaucoup d'actifs ($n$ grand), estimer $\Sigma$ est difficile.
PCA (Analyse en Composantes Principales) aide à :
\begin{itemize}
  \item comprendre les facteurs de risque dominants ;
  \item réduire la dimension (approximer $\Sigma$ par quelques facteurs) ;
  \item construire des portefeuilles exposés à des facteurs (eigen-portfolios).
\end{itemize}

\section{Décomposition spectrale de la covariance}
\begin{theorem}[Décomposition]
Si $\Sigma$ est symétrique définie positive, alors
\[
\Sigma = V\Lambda V^\top,
\]
où $V$ est orthogonale ($V^\top V=I$) et $\Lambda=\mathrm{diag}(\lambda_1,\dots,\lambda_n)$ avec $\lambda_1\ge\dots\ge\lambda_n>0$.
\end{theorem}

Les vecteurs colonnes $v_k$ de $V$ sont les directions principales (eigenvectors).
Les $\lambda_k$ mesurent la variance portée par chaque composante.

\subsection{Variance expliquée}
La part de variance expliquée par le facteur $k$ est
\[
\frac{\lambda_k}{\sum_{j=1}^n \lambda_j}.
\]
En finance actions, on observe souvent :
\begin{itemize}
  \item $\lambda_1$ très dominant (facteur ``marché'') ;
  \item quelques facteurs sectoriels/styles ;
  \item beaucoup de petits facteurs (bruit/idiosyncratique).
\end{itemize}

\section{Réduction de dimension : approximation à $m$ facteurs}
On approxime
\[
\Sigma \approx \sum_{k=1}^m \lambda_k v_k v_k^\top.
\]
Cela revient à un modèle factoriel :
\[
R \approx \mu + B F + \varepsilon,
\]
où $F\in\R^m$ sont les facteurs, $B$ les loadings, et $\varepsilon$ l'idiosyncratique.

\section{Eigen-portfolios}
Un eigen-portfolio associé au facteur $k$ est un portefeuille proportionnel à $v_k$.
Si on prend $w=v_k$ (normalisé), alors
\[
\Var(w^\top R)=v_k^\top \Sigma v_k = \lambda_k.
\]
Donc $v_1$ est le portefeuille de variance maximale sous contrainte $\norm{w}=1$.

\begin{example}[PCA sur 3 actifs (schéma)]
Supposons que la covariance a une première valeur propre très grande.
Alors le premier eigenvector est proche de $(1,1,1)/\sqrt{3}$ :
tous les actifs bougent ensemble (facteur commun).
Le second et le troisième représentent des spreads (long certains actifs, short d'autres).
\end{example}

\section{Shrinkage et régularisation via PCA}
Une technique pratique :
\begin{itemize}
  \item garder les $m$ premières composantes (structure) ;
  \item remplacer le reste par une variance moyenne (bruit).
\end{itemize}
On obtient une covariance ``plus stable'' et inversible, utile pour Markowitz.

\begin{remark}[Danger : instabilité des vecteurs propres]
Les eigenvectors estimés sont très sensibles au bruit,
surtout quand les valeurs propres sont proches.
Donc ``trader'' les eigen-portfolios de petits facteurs peut être extrêmement instable.
\end{remark}

\begin{jurybox}
\begin{itemize}
\item PCA diagonalise la covariance : $\Sigma=V\Lambda V^\top$ ; les $\lambda_k$ donnent la variance par facteur.
\item On peut approximer $\Sigma$ par $m$ facteurs dominants : réduction de dimension et stabilité.
\item Les eigen-portfolios ($v_k$) donnent des expositions directes aux facteurs de risque.
\item Attention : les petits facteurs (petites $\lambda$) sont souvent dominés par le bruit d'estimation.
\end{itemize}
\end{jurybox}

\chapter{Rebalancing, coûts de transaction et exécution : du modèle à la réalité}
\label{chap:rebalancing}

\section{Pourquoi le rebalancing est un problème non trivial ?}
Un portefeuille optimal théorique suppose qu'on peut :
\begin{itemize}
  \item ajuster instantanément les poids,
  \item sans coûts,
  \item sans impact sur les prix.
\end{itemize}
En réalité, chaque rebalancing crée :
\begin{itemize}
  \item du turnover (volume tradé),
  \item des coûts explicites (commissions) et implicites (spread, impact),
  \item des contraintes opérationnelles (liquidité, lots, contraintes réglementaires).
\end{itemize}
Donc on doit optimiser \emph{avec} les coûts.

\section{Mesure du turnover}
Entre deux dates, le turnover peut être approximé par
\[
\mathrm{TO} = \sum_{i=1}^n \abs{w_i^{\text{new}} - w_i^{\text{old}}}.
\]
Plus le turnover est grand, plus les coûts sont élevés.

\section{Modèles de coûts}
\subsection{Coûts proportionnels}
Simple :
\[
\text{coût} = c \sum_i \abs{\Delta q_i} \cdot \text{prix}_i,
\]
où $\Delta q_i$ est la variation de quantité.

\subsection{Impact quadratique (stylisé)}
Pour refléter que de gros ordres coûtent plus que proportionnellement :
\[
\text{coût} \approx \sum_i \left( c_1 \abs{\Delta q_i} + c_2 (\Delta q_i)^2\right).
\]
Cela incite à lisser l'exécution.

\section{Rebalancing périodique vs seuil}
\subsection{Périodique}
On rebalancera toutes les semaines/mois.
Avantage : simple.
Inconvénient : on trade même quand les poids ont peu dérivé.

\subsection{Seuil (band rebalancing)}
On rebalancera seulement si
\[
\abs{w_i^{\text{drift}} - w_i^{\text{target}}} > \epsilon_i.
\]
Avantage : réduit le turnover.
Inconvénient : plus complexe à analyser ; risque de dérive.

\begin{example}[Intuition]
Si le portefeuille cible est 60/40 actions/obligations et que le marché bouge peu,
un rebalancing hebdo crée des coûts inutiles.
Une règle à seuil attend que l'écart dépasse, par exemple, 2\% avant de trader.
\end{example}

\section{Optimisation ``cost-aware'' (idée)}
On peut modifier l'objectif Markowitz en ajoutant une pénalité de turnover :
\[
\min_w \quad w^\top \Sigma w - \lambda w^\top \mu + \eta \sum_i \abs{w_i-w_i^{\text{old}}},
\]
ou une pénalité quadratique $\eta\norm{w-w^{\text{old}}}^2$ (plus facile numériquement).
Cela produit des allocations plus ``lentes'' et plus réalistes.

\section{Lien avec le hedging d'options}
La logique est identique :
\begin{itemize}
  \item hedging = rebalancing d'une position $\Delta$ ;
  \item coûts de transaction = pénalité sur la variation de $\Delta$ ;
  \item compromis risque/coût = choix d'une fréquence ou d'un contrôle optimal.
\end{itemize}
C'est exactement ce qui motive les stratégies RL ``risk-aware'' (Partie III).

\begin{jurybox}
\begin{itemize}
\item Sans coûts, on rebalance ``tout le temps'' ; avec coûts, on cherche un compromis risque vs turnover.
\item Turnover $\sum \abs{\Delta w_i}$ est un proxy simple du coût.
\item Modèles de coûts : proportionnels (spread) et quadratiques (impact) ; ils changent la solution optimale.
\item Band rebalancing et pénalités de turnover sont des outils pratiques pour rendre un modèle allocable.
\end{itemize}
\end{jurybox}

\part{Exercices guidés (avec solutions) : consolider les notions}

\chapter{Feuille 1 : Brownien, Itô et Black--Scholes}
\label{chap:td1}

Ce chapitre propose une série d'exercices \emph{corrigés}.
L'objectif n'est pas de piéger, mais de forcer la maîtrise opérationnelle :
passer d'une intuition (``le bruit est en $\sqrt{t}$'') à un calcul concret.

\section{Exercice 1 --- incréments browniens et moments}
\begin{example}[Énoncé]
Soit $(W_t)_{t\ge 0}$ un Brownien standard.
\begin{enumerate}[label=\alph*)]
  \item Calculer $\E[W_t]$ et $\Var(W_t)$.
  \item Pour $0\le s<t$, calculer la loi de $W_t-W_s$ et $\Cov(W_s, W_t)$.
  \item Montrer que $\E[W_s W_t]=s$.
\end{enumerate}
\end{example}

\subsection*{Solution}
\begin{enumerate}[label=\alph*)]
  \item Par définition, $W_t\sim \Normal(0,t)$, donc $\E[W_t]=0$ et $\Var(W_t)=t$.
  \item Les incréments sont stationnaires et indépendants : $W_t-W_s\sim \Normal(0,t-s)$ et il est indépendant de $\F_s$.
  Pour la covariance, écrire $W_t=W_s+(W_t-W_s)$, d'où
  \[\Cov(W_s,W_t)=\Cov(W_s, W_s)+\Cov(W_s, W_t-W_s)=\Var(W_s)+0=s.\]
  \item Comme $\E[W_s]=0$ et $\E[W_t]=0$,
  \[\E[W_s W_t]=\Cov(W_s,W_t)=s.\]
\end{enumerate}

\begin{jurybox}
\begin{itemize}
  \item $W_t\sim\Normal(0,t)\Rightarrow$ échelle $\sqrt{t}$.
  \item $W_t=W_s+(W_t-W_s)$, avec incrément indépendant : outil de base pour covariances.
\end{itemize}
\end{jurybox}

\section{Exercice 2 --- variation quadratique : calcul guidé}
\begin{example}[Énoncé]
Soit une partition uniforme $0=t_0<t_1<\dots<t_n=T$ avec $t_k=k\Delta t$.
Définissons $Q_n=\sum_{k=0}^{n-1} (W_{t_{k+1}}-W_{t_k})^2$.
\begin{enumerate}[label=\alph*)]
  \item Calculer $\E[Q_n]$.
  \item Calculer $\Var(Q_n)$.
  \item Conclure que $Q_n\to T$ en probabilité.
\end{enumerate}
\end{example}

\subsection*{Solution}
Les incréments $\Delta W_k\triangleq W_{t_{k+1}}-W_{t_k}$ sont i.i.d. $\Normal(0,\Delta t)$.
\begin{enumerate}[label=\alph*)]
  \item $\E[(\Delta W_k)^2]=\Var(\Delta W_k)=\Delta t$, donc
  \[\E[Q_n]=\sum_{k=0}^{n-1} \Delta t = n\Delta t = T.\]
  \item Pour une gaussienne centrée $X\sim\Normal(0,\sigma^2)$, on a $\E[X^4]=3\sigma^4$.
  Donc $\Var(X^2)=\E[X^4]-\E[X^2]^2 = 3\sigma^4-\sigma^4=2\sigma^4$.
  Ici $\sigma^2=\Delta t$, donc $\Var((\Delta W_k)^2)=2(\Delta t)^2$.
  Par indépendance :
  \[\Var(Q_n)=\sum_{k=0}^{n-1}2(\Delta t)^2 = 2n(\Delta t)^2 = 2T\Delta t \xrightarrow[n\to\infty]{} 0.\]
  \item Comme $\E[Q_n]=T$ et $\Var(Q_n)\to 0$, on a $Q_n\to T$ en $L^2$, donc en probabilité.
\end{enumerate}

\section{Exercice 3 --- résoudre le GBM par Itô}
\begin{example}[Énoncé]
Sous $\Q$, supposons
\[\frac{\dd S_t}{S_t}=(r-q)\,\dd t + \sigma\,\dd W_t^{\Q}.
\]
\begin{enumerate}[label=\alph*)]
  \item Appliquer Itô à $\log S_t$ et en déduire une expression explicite de $S_T$.
  \item En déduire la loi de $\log S_T$.
  \item Calculer $\E^{\Q}[S_T]$ et vérifier qu'on retrouve le forward.
\end{enumerate}
\end{example}

\subsection*{Solution}
\begin{enumerate}[label=\alph*)]
  \item Pour $f(S)=\log S$, on a $f'(S)=1/S$ et $f''(S)=-1/S^2$.
  Par Itô :
  \[\dd\log S_t = \frac{1}{S_t}\dd S_t + \frac12\left(-\frac{1}{S_t^2}\right)(\dd S_t)^2.
  \]
  Or $(\dd S_t)^2=(\sigma S_t)^2\dd t=\sigma^2 S_t^2\dd t$.
  Donc
  \[\dd\log S_t = (r-q)\dd t + \sigma\dd W_t^{\Q} - \frac12\sigma^2\dd t
  = \left((r-q)-\tfrac12\sigma^2\right)\dd t + \sigma\dd W_t^{\Q}.
  \]
  En intégrant :
  \[\log S_T = \log S_0 + \left((r-q)-\tfrac12\sigma^2\right)T + \sigma W_T^{\Q}.
  \]
  Donc
  \[S_T = S_0\exp\Bigl(\left((r-q)-\tfrac12\sigma^2\right)T + \sigma W_T^{\Q}\Bigr).
  \]
  \item Comme $W_T^{\Q}\sim\Normal(0,T)$, on a
  \[\log S_T\sim \Normal\Bigl(\log S_0+\left((r-q)-\tfrac12\sigma^2\right)T,\ \sigma^2 T\Bigr).
  \]
  \item Par propriété de la lognormale, si $Y\sim\Normal(m,s^2)$ alors $\E[\ee^Y]=\ee^{m+\tfrac12 s^2}$.
  Ici, $m=\log S_0+((r-q)-\tfrac12\sigma^2)T$ et $s^2=\sigma^2 T$.
  Donc
  \[\E^{\Q}[S_T]=\E^{\Q}[\ee^{\log S_T}]=\exp\Bigl(m+\tfrac12 s^2\Bigr)=S_0\ee^{(r-q)T}=F_{0,T}.
  \]
\end{enumerate}

\section{Exercice 4 --- retrouver la formule de Black--Scholes}
\begin{example}[Énoncé]
On considère un call européen $C_0$ de maturité $T$ et strike $K$.
En utilisant la loi de $\log S_T$ sous $\Q$, montrer que
\[C_0=S_0\ee^{-qT}\Phi(d_1)-K\ee^{-rT}\Phi(d_2)
\quad\text{avec}\quad
d_1=\frac{\log(S_0/K)+(r-q+\tfrac12\sigma^2)T}{\sigma\sqrt{T}},\ d_2=d_1-\sigma\sqrt{T}.
\]
\end{example}

\subsection*{Solution (détaillée)}
Le prix risque-neutre est
\[C_0=\E^{\Q}[\ee^{-rT}(S_T-K)^+].\]
On sépare :
\[C_0=\ee^{-rT}\E^{\Q}[S_T\1_{\{S_T>K\}}]-\ee^{-rT}K\Q(S_T>K).\]
Posons $X\triangleq\log S_T$ ; alors $X\sim\Normal(m,s^2)$ avec
\[m=\log S_0+\left((r-q)-\tfrac12\sigma^2\right)T,\qquad s=\sigma\sqrt{T}.
\]
La condition $S_T>K$ est $X>\log K$.
\begin{itemize}
  \item Pour la probabilité :
  \[\Q(S_T>K)=\Q(X>\log K)=\Phi\left(\frac{m-\log K}{s}\right)=\Phi(d_2),\]
  où l'on vérifie que
  \[d_2=\frac{\log(S_0/K)+(r-q-\tfrac12\sigma^2)T}{\sigma\sqrt{T}}.
  \]
  \item Pour l'espérance tronquée $\E[S_T\1_{\{S_T>K\}}]=\E[\ee^X\1_{\{X>\log K\}}]$,
  on utilise un calcul standard de gaussienne complétée au carré.
  Écrivons la densité de $X$ :
  \[f(x)=\frac{1}{s\sqrt{2\pi}}\exp\left(-\frac{(x-m)^2}{2s^2}\right).
  \]
  Alors
  \[\E[\ee^X\1_{\{X>\log K\}}]=\int_{\log K}^{\infty}\ee^x f(x)\,\dd x.
  \]
  On combine les exponentielles et on complète le carré :
  \[
  \ee^x\exp\left(-\frac{(x-m)^2}{2s^2}\right)
  =\exp\left(-\frac{(x-(m+s^2))^2}{2s^2}\right)\exp\left(m+\tfrac12 s^2\right).
  \]
  Donc
  \[\E[\ee^X\1_{\{X>\log K\}}]=\exp\left(m+\tfrac12 s^2\right)\Phi\left(\frac{m+s^2-\log K}{s}\right).
  \]
  Or $\exp(m+\tfrac12 s^2)=\E[S_T]=S_0\ee^{(r-q)T}$ et
  \[\frac{m+s^2-\log K}{s}=\frac{\log(S_0/K)+(r-q+\tfrac12\sigma^2)T}{\sigma\sqrt{T}}=d_1.
  \]
  Finalement
  \[\E[S_T\1_{\{S_T>K\}}]=S_0\ee^{(r-q)T}\Phi(d_1).
  \]
\end{itemize}
En remplaçant dans la formule du prix :
\[C_0=\ee^{-rT}\Bigl(S_0\ee^{(r-q)T}\Phi(d_1)-K\Phi(d_2)\Bigr)=S_0\ee^{-qT}\Phi(d_1)-K\ee^{-rT}\Phi(d_2).
\]

\begin{jurybox}
\begin{itemize}
  \item La formule BS se déduit d'une espérance risque-neutre + loi lognormale.
  \item Le ``truc'' technique est la complétion du carré pour intégrer $\ee^x$ contre une densité gaussienne.
\end{itemize}
\end{jurybox}

\chapter{Feuille 2 : volatilité implicite, surface et calibration (Heston/SABR)}
\label{chap:td2}

\section{Exercice 1 --- densité risque-neutre et convexité}
\begin{example}[Énoncé]
Fixons une maturité $T$.
On rappelle la relation de Breeden--Litzenberger :
\[\pdd{C}{K}(K,T)=\ee^{-rT} f_{S_T}(K).
\]
\begin{enumerate}[label=\alph*)]
  \item Expliquer pourquoi la convexité de $C(K,T)$ en $K$ est une condition de non-arbitrage.
  \item Proposer une approximation discrète de $\pdd{C}{K}$ à partir de prix de calls observés sur une grille de strikes.
  \item Comment vérifier qualitativement, à partir de la densité, la présence d'un skew ``put'' ?
\end{enumerate}
\end{example}

\subsection*{Solution}
\begin{enumerate}[label=\alph*)]
  \item Un call est une espérance de $(S_T-K)^+$, qui est convexe en $K$.
  L'espérance conserve la convexité : si $C$ n'est pas convexe, on peut construire un butterfly statique (combinaison linéaire de calls) avec payoff positif mais prix négatif $\Rightarrow$ arbitrage.
  \item Sur strikes équidistants $K_{i-1},K_i,K_{i+1}$, une seconde différence finie :
  \[\pdd{C}{K}(K_i) \approx \frac{C(K_{i+1})-2C(K_i)+C(K_{i-1})}{(\Delta K)^2}.
  \]
  \item Un skew put sur indices correspond à une queue gauche plus lourde en $\Q$ : la densité a plus de masse pour des $S_T$ bas.
  Qualitativement, cela se traduit par des calls très OTM (grands $K$) relativement bon marché et des puts OTM (petits $K$) chers ; via la relation ci-dessus, la densité est plus asymétrique vers la gauche.
\end{enumerate}

\section{Exercice 2 --- inversion de la volatilité implicite (Newton sécurisé)}
\begin{example}[Énoncé]
On observe un prix de call $C^{\mathrm{mkt}}$.
On veut l'IV $\sigma$ telle que $C^{\mathrm{BS}}(\sigma)=C^{\mathrm{mkt}}$.
\begin{enumerate}[label=\alph*)]
  \item Écrire l'itération de Newton.
  \item Expliquer pourquoi Newton peut diverger pour des options très OTM.
  \item Proposer un schéma "robuste" combinant bisection et Newton.
\end{enumerate}
\end{example}

\subsection*{Solution}
\begin{enumerate}[label=\alph*)]
  \item $\sigma_{n+1}=\sigma_n-\frac{C^{\mathrm{BS}}(\sigma_n)-C^{\mathrm{mkt}}}{\vega(\sigma_n)}$.
  \item Quand l'option est très OTM ou très proche de maturité, $\vega$ est faible.
  La division par un petit nombre amplifie l'erreur, et l'approximation linéaire de Newton est mauvaise $\Rightarrow$ saut trop grand.
  \item Encadrer $\sigma$ dans un intervalle $[\sigma_{\min},\sigma_{\max}]$ (ex. $[10^{-4},5]$).
  À chaque itération, faire un pas de Newton ; s'il sort de l'intervalle ou n'améliore pas l'erreur, remplacer par une bisection.
  Ce schéma converge toujours (bisection), tout en restant rapide quand Newton est bien conditionné.
\end{enumerate}

\section{Exercice 3 --- condition de Feller et interprétation}
\begin{example}[Énoncé]
Dans Heston : $\dd v_t=\kappa(\theta-v_t)\dd t+\xi\sqrt{v_t}\dd W_t$.
\begin{enumerate}[label=\alph*)]
  \item Rappeler la condition de Feller.
  \item Pour $(\kappa,\theta,\xi)=(2,0.04,0.6)$, la condition est-elle satisfaite ?
  \item Que se passe-t-il numériquement si $v_t$ touche 0 ?
\end{enumerate}
\end{example}

\subsection*{Solution}
\begin{enumerate}[label=\alph*)]
  \item $2\kappa\theta\ge \xi^2$.
  \item $2\kappa\theta=2\cdot 2\cdot 0.04=0.16$ et $\xi^2=0.36$ $\Rightarrow$ non.
  \item Si $v_t$ approche 0, la diffusion $\xi\sqrt{v_t}$ s'éteint mais le terme de drift peut réinjecter.
  En simulation discrète, on peut obtenir des valeurs négatives si l'on n'est pas prudent (Euler naïf) ; d'où l'usage de schémas adaptés (full truncation, QE, réflexion) et/ou reparamétrisations.
\end{enumerate}

\section{Exercice 4 --- SABR : calcul ATM}
\begin{example}[Énoncé]
Dans la formule de Hagan, l'IV ATM (au premier ordre) s'écrit
\[\sigma_{\mathrm{ATM}}\approx \frac{\alpha}{F^{1-\beta}}\left(1 + \left(\frac{(1-\beta)^2}{24}\frac{\alpha^2}{F^{2-2\beta}}+\frac{\rho\beta\nu\alpha}{4F^{1-\beta}}+\frac{2-3\rho^2}{24}\nu^2\right)T\right).
\]
Pour $F=100$, $\beta=0.5$, $\alpha=0.2$, $\nu=0.6$, $\rho=-0.3$, $T=1$ :
\begin{enumerate}[label=\alph*)]
  \item Calculer le terme principal $\alpha/F^{1-\beta}$.
  \item Évaluer la correction en $T$ (ordre de grandeur).
\end{enumerate}
\end{example}

\subsection*{Solution}
\begin{enumerate}[label=\alph*)]
  \item $F^{1-\beta}=100^{0.5}=10$, donc terme principal $0.2/10=0.02$ (soit 2\%).
  \item On calcule séparément les trois contributions dans la parenthèse :
  \[
  \frac{(1-\beta)^2}{24}\frac{\alpha^2}{F^{2-2\beta}}=\frac{0.25}{24}\frac{0.04}{100^{1}}\approx \frac{0.25}{24}\cdot 0.0004 \approx 4.2\times 10^{-6},
  \]
  \[
  \frac{\rho\beta\nu\alpha}{4F^{1-\beta}}=\frac{-0.3\cdot 0.5\cdot 0.6\cdot 0.2}{4\cdot 10}\approx \frac{-0.018}{40}\approx -4.5\times 10^{-4},
  \]
  \[
  \frac{2-3\rho^2}{24}\nu^2=\frac{2-3\cdot 0.09}{24}\cdot 0.36=\frac{1.73}{24}\cdot 0.36\approx 0.02595.
  \]
  La correction est dominée par le terme en $\nu^2$ : environ $+2.6\%$ en relatif.
  Donc $\sigma_{\mathrm{ATM}}\approx 0.02\times(1+0.026)\approx 0.0205$ (2.05\%).
\end{enumerate}

\begin{jurybox}
\begin{itemize}
  \item L'IV de marché encode une densité risque-neutre ; les contraintes d'arbitrage sont des contraintes de convexité/monotonie.
  \item Newton pour l'IV est rapide mais doit être sécurisé (vega faible).
  \item Dans Heston, Feller est une contrainte de robustesse ; dans SABR, l'expansion de Hagan permet des calculs quasi fermés.
\end{itemize}
\end{jurybox}

\chapter{Feuille 3 : courbe des taux, optimisation de portefeuille et risque}
\label{chap:td3}

\section{Exercice 1 --- bootstrapping minimal (zéro-coupons)}
\begin{example}[Énoncé]
On observe des taux zéro-coupon (compounded en continu) :
\[R(0,1)=3\%,\qquad R(0,2)=3.5\%,\qquad R(0,3)=4\%.
\]
\begin{enumerate}[label=\alph*)]
  \item Calculer les discount factors $P(0,T)$ pour $T=1,2,3$.
  \item Calculer le forward (continu) $f(1,2)$ tel que $P(0,2)=P(0,1)\ee^{-f(1,2)(2-1)}$.
\end{enumerate}
\end{example}

\subsection*{Solution}
\begin{enumerate}[label=\alph*)]
  \item En continu, $P(0,T)=\ee^{-R(0,T)T}$.
  Donc $P(0,1)=\ee^{-0.03}\approx 0.9704$,
  $P(0,2)=\ee^{-0.035\cdot 2}=\ee^{-0.07}\approx 0.9324$,
  $P(0,3)=\ee^{-0.04\cdot 3}=\ee^{-0.12}\approx 0.8869$.
  \item On impose $P(0,2)=P(0,1)\ee^{-f(1,2)}$.
  Donc $f(1,2)=\log\left(\frac{P(0,1)}{P(0,2)}\right)=\log\left(\frac{0.9704}{0.9324}\right)\approx 0.0399$ (3.99\%).
\end{enumerate}

\section{Exercice 2 --- Markowitz (2 actifs) : solution fermée}
\begin{example}[Énoncé]
Deux actifs avec rendements espérés $\mu=(8\%,4\%)$ et covariance
\[\Sigma=\begin{pmatrix}0.04 & 0.01\\ 0.01 & 0.01\end{pmatrix}.
\]
On cherche le portefeuille de variance minimale sous contrainte $w_1+w_2=1$.
\begin{enumerate}[label=\alph*)]
  \item Résoudre par Lagrangien.
  \item Calculer la variance du portefeuille optimal.
\end{enumerate}
\end{example}

\subsection*{Solution}
On minimise $w^\top\Sigma w$ sous $\1^\top w=1$.
Lagrangien : $\mathcal{L}(w,\lambda)=w^\top\Sigma w-\lambda(\1^\top w-1)$.
Condition du premier ordre : $2\Sigma w-\lambda\1=0$, donc
\[w=\frac{\lambda}{2}\Sigma^{-1}\1.
\]
Imposer $\1^\top w=1$ donne
\[\frac{\lambda}{2}\1^\top\Sigma^{-1}\1=1\quad\Rightarrow\quad \lambda=\frac{2}{\1^\top\Sigma^{-1}\1}.
\]
Donc
\[w^{\mathrm{MV}}=\frac{\Sigma^{-1}\1}{\1^\top\Sigma^{-1}\1}.
\]
Calculons $\Sigma^{-1}$ : déterminant $0.04\cdot 0.01-0.01^2=0.0004-0.0001=0.0003$.
\[\Sigma^{-1}=\frac{1}{0.0003}\begin{pmatrix}0.01 & -0.01\\ -0.01 & 0.04\end{pmatrix}
=\begin{pmatrix}33.33 & -33.33\\ -33.33 & 133.33\end{pmatrix}.
\]
Alors $\Sigma^{-1}\1=(0,100)^{\top}$ et $\1^\top\Sigma^{-1}\1=100$.
Donc $w=(0,1)$ : tout sur l'actif 2 (beaucoup moins risqué).
Variance : $w^\top\Sigma w=0.01$.

\section{Exercice 3 --- risk parity à deux actifs}
\begin{example}[Énoncé]
Deux actifs non corrélés de volatilités $\sigma_1=20\%$ et $\sigma_2=10\%$.
En risk parity (sans levier), on veut des contributions au risque égales.
\begin{enumerate}[label=\alph*)]
  \item Trouver les poids $w_1,w_2$.
  \item Donner l'intuition économique.
\end{enumerate}
\end{example}

\subsection*{Solution}
Si la corrélation est nulle, la volatilité du portefeuille est
\[\sigma_p=\sqrt{(w_1\sigma_1)^2+(w_2\sigma_2)^2}.
\]
La contribution marginale au risque de l'actif $i$ vaut (ici)
\[\mathrm{RC}_i\propto w_i\sigma_i^2.
\]
L'égalité des contributions impose $w_1\sigma_1^2=w_2\sigma_2^2$.
Donc
\[\frac{w_1}{w_2}=\frac{\sigma_2^2}{\sigma_1^2}=\frac{0.1^2}{0.2^2}=\frac{1}{4}.
\]
Avec $w_1+w_2=1$, on obtient $w_1=0.2$ et $w_2=0.8$.

\textbf{Intuition :} risk parity met plus de poids sur l'actif moins risqué pour équilibrer les contributions.
Ce n'est pas "min-variance" : la contrainte d'égalité des contributions est une préférence de diversification du risque.

\begin{jurybox}
\begin{itemize}
  \item Bootstrapping : $P(0,T)=\ee^{-R(0,T)T}$ (continu) ; les forwards se déduisent par ratios de discount.
  \item Markowitz : la solution min-variance a une forme fermée via $\Sigma^{-1}\1$.
  \item Risk parity : égaliser des contributions revient souvent à surpondérer les actifs moins volatils.
\end{itemize}
\end{jurybox}

\chapter{Feuille 4 : méthodes numériques (Monte Carlo, greeks, PDE) --- exercices}
\label{chap:td4}

Cette feuille vise à relier les méthodes numériques des annexes à des calculs concrets.
L'idée centrale : \emph{toute} méthode numérique est un compromis entre biais (erreur systématique) et variance (bruit),
et ce compromis se pilote par la discrétisation (pas de temps/grille) et le budget de calcul (nombre de trajectoires, taille de grille).

\section{Exercice 1 --- Monte Carlo : erreur en $1/\sqrt{M}$}
\begin{example}[Énoncé]
On veut estimer $V_0=\E[\ee^{-rT}g(S_T)]$ par Monte Carlo.
On simule $M$ trajectoires i.i.d. et on calcule
\[\widehat{V}_0=\frac{1}{M}\sum_{m=1}^M \ee^{-rT} g\bigl(S_T^{(m)}\bigr).
\]
\begin{enumerate}[label=\alph*)]
  \item Montrer que $\widehat{V}_0$ est un estimateur sans biais de $V_0$.
  \item Exprimer $\Var(\widehat{V}_0)$ en fonction de $\Var(\ee^{-rT}g(S_T))$.
  \item Déduire un intervalle de confiance asymptotique à 95\%.
\end{enumerate}
\end{example}

\subsection*{Solution}
\begin{enumerate}[label=\alph*)]
  \item Par linéarité de l'espérance et i.i.d. :
  \[\E[\widehat{V}_0]=\frac{1}{M}\sum_{m=1}^M \E\bigl[\ee^{-rT} g(S_T^{(m)})\bigr]=\E\bigl[\ee^{-rT}g(S_T)\bigr]=V_0.
  \]
  \item Si $Y=\ee^{-rT}g(S_T)$, alors $\widehat{V}_0=\frac{1}{M}\sum_{m=1}^M Y^{(m)}$.
  Donc $\Var(\widehat{V}_0)=\Var(Y)/M$.
  \item Par TCL : $\sqrt{M}(\widehat{V}_0-V_0)\Rightarrow \Normal(0,\Var(Y))$.
  Un IC 95\% est
  \[\widehat{V}_0 \pm 1.96\,\frac{\widehat{\sigma}_Y}{\sqrt{M}},\]
  où $\widehat{\sigma}_Y^2$ est la variance empirique des $Y^{(m)}$.
\end{enumerate}

\begin{jurybox}
\begin{itemize}
  \item Monte Carlo a une convergence lente mais robuste : erreur typique $\propto 1/\sqrt{M}$.
  \item Doubler la précision (diviser l'erreur par 2) demande environ \(4\times\) plus de scénarios.
\end{itemize}
\end{jurybox}

\section{Exercice 2 --- réduction de variance : variable de contrôle}
\begin{example}[Énoncé]
On price un call européen dans Black--Scholes par Monte Carlo. On connaît aussi le prix exact $C_{\BS}$.
On propose une variable de contrôle : $Z=(S_T-K)^+$ simulé, et un contrôle $Z_0=C_{\BS}\ee^{rT}$ (valeur non actualisée).
On définit
\[\widehat{V}_{\mathrm{cv}} = \ee^{-rT}\left(\overline{Z} - \beta(\overline{Z}-Z_0)\right)
=\ee^{-rT}\left((1-\beta)\overline{Z}+\beta Z_0\right).
\]
\begin{enumerate}[label=\alph*)]
  \item Montrer que l'estimateur est sans biais pour tout $\beta$.
  \item Trouver $\beta^*$ qui minimise la variance (en fonction de $\Var(Z)$).
  \item Interpréter économiquement.
\end{enumerate}
\end{example}

\subsection*{Solution}
\begin{enumerate}[label=\alph*)]
  \item $\E[\overline{Z}]=\E[Z]=C_{\BS}\ee^{rT}=Z_0$.
  Donc $\E[\widehat{V}_{\mathrm{cv}}]=\ee^{-rT}((1-\beta)Z_0+\beta Z_0)=\ee^{-rT}Z_0=C_{\BS}$.
  \item Ici, la variable de contrôle est exactement le payoff du call, donc $\overline{Z}-Z_0$ a moyenne nulle.
  On a
  \[\Var(\widehat{V}_{\mathrm{cv}})=\ee^{-2rT}(1-\beta)^2\,\Var(\overline{Z})
  =\ee^{-2rT}(1-\beta)^2\,\frac{\Var(Z)}{M}.
  \]
  Le minimum est atteint pour $\beta^*=1$ : variance nulle.
  \item C'est volontairement ``trop beau'' : si on connaît $C_{\BS}$, on n'a pas besoin de MC.
  Le message : une bonne variable de contrôle est un produit \emph{proche} du payoff dont on connaît le prix.
  En pratique : contrôle via le prix BS avec une volatilité proche, via des greeks, ou via des approximations rapides.
\end{enumerate}

\section{Exercice 3 --- greeks par bump \& revalue : choisir le bon epsilon}
\begin{example}[Énoncé]
On estime le delta par différences finies :
\[\widehat{\Delta} = \frac{\widehat{V}(S_0+\varepsilon)-\widehat{V}(S_0-\varepsilon)}{2\varepsilon}.
\]
\begin{enumerate}[label=\alph*)]
  \item Expliquer pourquoi $\varepsilon$ trop petit \emph{augmente} le bruit.
  \item Expliquer pourquoi $\varepsilon$ trop grand \emph{biais} l'estimation.
  \item Donner une règle de choix d'ordre de grandeur.
\end{enumerate}
\end{example}

\subsection*{Solution}
\begin{enumerate}[label=\alph*)]
  \item Les deux estimations Monte Carlo ont une variance $\propto 1/M$.
  La différence au numérateur a donc un bruit typique de taille $\propto 1/\sqrt{M}$.
  En divisant par $2\varepsilon$, la variance de $\widehat{\Delta}$ se comporte comme $\propto 1/(M\varepsilon^2)$.
  Donc $\varepsilon\downarrow 0$ fait \emph{exploser} la variance.
  \item L'approximation par différences finies suppose une linéarité locale :
  \(V(S_0+\varepsilon)=V(S_0)+\varepsilon V'(S_0)+O(\varepsilon^2)\).
  Le biais de la différence centrée est en $O(\varepsilon^2)$.
  \item On choisit $\varepsilon$ pour équilibrer variance ($\propto 1/(M\varepsilon^2)$) et biais ($\propto \varepsilon^2$).
  Cela suggère $\varepsilon\propto M^{-1/4}$ (règle de pouce).
  En pratique : commencer par $\varepsilon\approx 10^{-3}S_0$ et tester la stabilité.
\end{enumerate}

\section{Exercice 4 --- PDE : schéma implicite et stabilité}
\begin{example}[Énoncé]
Considérons la PDE de Black--Scholes en $S$ et $t$.
On discrétise $t$ par pas $\Delta t$ et $S$ par un pas $\Delta S$.
\begin{enumerate}[label=\alph*)]
  \item Expliquer qualitativement pourquoi un schéma \emph{explicite} peut être instable si $\Delta t$ est trop grand.
  \item Expliquer pourquoi un schéma \emph{implicite} est plus stable mais demande de résoudre un système linéaire.
  \item Donner l'intuition de Crank--Nicolson (moyenne explicite/implicite).
\end{enumerate}
\end{example}

\subsection*{Solution}
\begin{enumerate}[label=\alph*)]
  \item La PDE contient un terme de diffusion $\tfrac12\sigma^2 S^2\pdd{V}{S}$.
  Comme pour l'équation de la chaleur, un schéma explicite impose une condition CFL du type
  \(\Delta t\lesssim c\,(\Delta S)^2\) (ici avec un facteur dépendant de $S$).
  Sinon, l'erreur numérique s'amplifie à chaque pas de temps.
  \item Le schéma implicite ``met'' les inconnues au temps futur dans le terme de diffusion.
  Il est généralement inconditionnellement stable, mais il impose de résoudre un système tridiagonal à chaque pas (coût mais contrôlable).
  \item Crank--Nicolson prend la moyenne : moitié explicite, moitié implicite.
  Il est d'ordre 2 en temps (plus précis) mais peut osciller si les conditions aux bords sont mal gérées ; on le stabilise parfois par \emph{Rannacher} (quelques pas implicites puis CN).
\end{enumerate}

\begin{jurybox}
\begin{itemize}
  \item MC : variance $\propto 1/M$ ; pour gagner un chiffre significatif, il faut beaucoup de scénarios.
  \item Réduction de variance : antithétiques, contrôles, stratification, quasi-MC (souvent plus rentable que d'augmenter $M$).
  \item Greeks : différences finies = compromis biais/variance ; méthode pathwise/LR quand possible.
  \item PDE : explicite = condition CFL ; implicite = stable mais système linéaire ; CN = précision/stabilité intermédiaire.
\end{itemize}
\end{jurybox}

\section{Exercice 5 --- option américaine : algorithme de Longstaff--Schwartz (LSM)}
\begin{example}[Énoncé]
On veut pricer un put américain de maturité $T$ et strike $K$ sur un sous-jacent suivant Black--Scholes.
On simule $M$ trajectoires $S_{t_0},S_{t_1},\dots,S_{t_N}$, où $t_N=T$.
\begin{enumerate}[label=\alph*)]
  \item Expliquer la difficulté spécifique d'une option américaine (comparée à l'européenne) en termes d'\emph{exercice optimal}.
  \item Décrire l'algorithme LSM : backward induction et régression de la valeur de continuation.
  \item Donner une base de fonctions (features) simple pour approximer l'espérance conditionnelle (ex. polynômes de Laguerre ou polynômes en $S$).
\end{enumerate}
\end{example}

\subsection*{Solution (pas à pas)}
\begin{enumerate}[label=\alph*)]
  \item Pour un américain, à chaque date $t_k$ on choisit entre :
  \(\text{exercer maintenant : } g(S_{t_k})\) ou \(\text{continuer : } \E[\ee^{-r\Delta t}V_{k+1}\mid\F_{t_k}]\).
  La frontière d'exercice (ensemble des $S$ où l'exercice est optimal) est inconnue \emph{a priori}.
  \item LSM approxime la valeur de continuation par une régression sur les trajectoires simulées.
  Schéma :
  \begin{enumerate}[label=\roman*)]
    \item Initialiser à maturité : $V_N^{(m)}=g(S_{t_N}^{(m)})$.
    \item Pour $k=N-1,\dots,1$ :
      \begin{itemize}
        \item Considérer les trajectoires \emph{dans la monnaie} (pour un put : $S_{t_k}^{(m)}<K$).
        \item Calculer les cashflows actualisés au pas suivant : $Y^{(m)}=\ee^{-r\Delta t}V_{k+1}^{(m)}$.
        \item Régression : approximer $\E[Y\mid S_{t_k}]\approx \sum_{j} a_j\,\phi_j(S_{t_k})$.
        \item Décision : exercer si $g(S_{t_k}^{(m)}) \ge \widehat{\text{cont}}(S_{t_k}^{(m)})$, sinon continuer.
        \item Mettre à jour $V_k^{(m)}$ en conséquence.
      \end{itemize}
    \item Le prix est la moyenne des cashflows au temps 0 : $V_0\approx \frac{1}{M}\sum_m \ee^{-rt_{\tau^{(m)}}}g(S_{\tau^{(m)}}^{(m)})$, où $\tau^{(m)}$ est le temps d'exercice simulé.
  \end{enumerate}
  \item Bases classiques :
  \(\phi_0(S)=1\), \(\phi_1(S)=S\), \(\phi_2(S)=S^2\) (simple mais peut être instable),
  ou bases de Laguerre (souvent plus stables) :
  \[L_0(x)=1,\quad L_1(x)=1-x,\quad L_2(x)=1-2x+\tfrac12x^2,\quad x=S/K.
  \]
\end{enumerate}

\begin{jurybox}
\begin{itemize}
  \item Américain = optimisation de stopping : exercice vs continuation à chaque date.
  \item LSM : on remonte dans le temps en régressant la continuation sur les trajectoires (approximation de $\E[\cdot\mid S]$).
  \item Qualité dépend de (i) la base de régression, (ii) le nombre de trajectoires, (iii) la discrétisation temporelle.
\end{itemize}
\end{jurybox}

\section{Exercice 6 --- FFT (Carr--Madan) : choisir la grille et éviter l'aliasing}
\begin{example}[Énoncé]
Dans Carr--Madan, on price un call via une transformée de Fourier du prix amorti.
On utilise une grille en fréquence $u_j=j\,\Delta u$ et une grille en log-strike $k_m=k_{\min}+m\,\Delta k$ avec la relation $\Delta u\,\Delta k=2\pi/N$.
\begin{enumerate}[label=\alph*)]
  \item Expliquer l'intuition derrière le paramètre d'amortissement $\alpha>0$ (pourquoi on multiplie par $\ee^{\alpha k}$).
  \item Expliquer qualitativement le phénomène d'\emph{aliasing} quand la grille en $k$ est trop courte.
  \item Donner deux règles pratiques pour choisir $(N,\Delta u,k_{\min})$.
\end{enumerate}
\end{example}

\subsection*{Solution (intuitions pratiques)}
\begin{enumerate}[label=\alph*)]
  \item Le payoff d'un call n'est pas intégrable sur $k=\log K$ : la transformée diverge.
  On amortit en considérant un prix modifié $\widetilde{C}(k)=\ee^{\alpha k}C(k)$.
  Pour $\alpha$ assez grand, $\widetilde{C}$ devient intégrable et la Fourier est bien définie.
  \item L'FFT suppose une périodisation implicite : on reconstruit une fonction périodique de période $N\Delta k$.
  Si la vraie fonction n'est pas quasi nulle aux bords, les ``copies'' se recollent et créent des erreurs (aliasing), typiquement des oscillations et des prix négatifs dans les ailes.
  \item Règles de pouce :
  \begin{itemize}
    \item prendre $N$ une puissance de 2 (efficacité FFT) et suffisamment grand (ex. $N\ge 2^{10}$) ;
    \item choisir $\Delta u$ petit pour bien intégrer les oscillations (meilleure précision) ;
    \item choisir $k_{\min}$ et la largeur $N\Delta k$ pour couvrir une large plage de strikes (ex. de $K\approx 0.05F$ à $K\approx 20F$) ;
    \item tester la stabilité en faisant varier $\alpha$ et la largeur de grille : un modèle sain donne des prix quasi inchangés.
  \end{itemize}
\end{enumerate}

\section{Exercice 7 --- quasi-Monte Carlo : quand et pourquoi cela aide}
\begin{example}[Énoncé]
On remplace les tirages pseudo-aléatoires uniformes par une suite à faible discrépance (Sobol, Halton) puis on transforme en gaussiennes (inverse de la CDF $\Phi$ ou Box--Muller).
\begin{enumerate}[label=\alph*)]
  \item Expliquer pourquoi, pour une intégrale $I=\int_{[0,1]^d} f(u)\,\dd u$, l'erreur quasi-MC peut être \emph{meilleure} que $1/\sqrt{M}$ quand $f$ est régulière.
  \item En pricing, pourquoi parle-t-on d'une \emph{dimension effective} souvent plus faible que $d=N$ (nombre de pas de temps) ?
  \item Donner une intuition de la technique \emph{Brownian bridge} pour réduire la dimension effective.
\end{enumerate}
\end{example}

\subsection*{Solution (intuition + points pratiques)}
\begin{enumerate}[label=\alph*)]
  \item Les suites à faible discrépance couvrent l'hypercube plus uniformément que des tirages i.i.d.
  Pour des fonctions assez régulières, l'erreur d'intégration est liée à la discrépance via l'inégalité de Koksma--Hlawka (idée : "bonne couverture" $\times$ "variation" de $f$).
  Cela peut donner, dans les cas favorables, une convergence proche de $O\bigl((\log M)^d/M\bigr)$, donc plus rapide que $1/\sqrt{M}$.
  \item Dans un MC discretisé, on croit intégrer en dimension $d=N$ (un uniforme/gaussienne par pas), mais souvent le payoff dépend surtout de quelques "facteurs" :
  niveau du spot à maturité, moyenne, ou quelques combinaisons.
  Les contributions des pas tardifs peuvent être dominantes ; les pas très fins apportent peu.
  C'est ce qu'on appelle la \emph{dimension effective}.
  \item Brownian bridge réordonne la construction de la trajectoire : on simule d'abord $W_T$, puis les points intermédiaires conditionnellement.
  On "met" la variance la plus importante sur les premières dimensions de la suite Sobol, ce qui réduit la dimension effective et améliore l'efficacité quasi-MC.
\end{enumerate}

\begin{figure}[h!]
\centering
\begin{tikzpicture}
\begin{axis}[
  width=0.75\textwidth,
  height=0.32\textwidth,
  xlabel={$M$ (nombre de scénarios)},
  ylabel={erreur typique (échelle relative)},
  xmode=log,
  ymode=log,
  legend pos=south west,
  grid=both
]
\addplot+[mark=none, thick] coordinates {(1e2,1e-1) (1e3,3.16e-2) (1e4,1e-2) (1e5,3.16e-3)};
\addlegendentry{$1/\sqrt{M}$ (MC)}
\addplot+[mark=none, thick, dashed] coordinates {(1e2,5e-2) (1e3,1.5e-2) (1e4,4.5e-3) (1e5,1.4e-3)};
\addlegendentry{profil plus rapide (quasi-MC)}
\end{axis}
\end{tikzpicture}
\caption{Illustration qualitative : MC décroît en $M^{-1/2}$ ; quasi-MC peut faire mieux sur des payoffs réguliers et une faible dimension effective.}
\end{figure}

\section{Exercice 8 --- greeks Monte Carlo : méthode pathwise vs likelihood ratio}
\begin{example}[Énoncé]
On veut estimer le delta d'un payoff européen $g(S_T)$ en Black--Scholes.
\begin{enumerate}[label=\alph*)]
  \item (Pathwise) En supposant $g$ dérivable, montrer formellement que
  \[\Delta = \pd{V_0}{S_0}=\E\left[\ee^{-rT} g'(S_T)\,\pd{S_T}{S_0}\right].\]
  \item Dans Black--Scholes, exprimer $\pd{S_T}{S_0}$.
  \item (Likelihood ratio) Expliquer le principe : dériver la \emph{densité} plutôt que le payoff, utile quand $g$ n'est pas dérivable (digitals).
\end{enumerate}
\end{example}

\subsection*{Solution}
\begin{enumerate}[label=\alph*)]
  \item Si on peut permuter dérivation et espérance (hypothèses de domination), alors
  \[V_0=\E[\ee^{-rT}g(S_T(S_0,\omega))]\quad\Rightarrow\quad \pd{V_0}{S_0}=\E\left[\ee^{-rT}g'(S_T)\,\pd{S_T}{S_0}\right].\]
  C'est la méthode \emph{pathwise} : on dérive trajectoire par trajectoire.
  \item Sous BS : $S_T=S_0\exp\bigl((r-q-\tfrac12\sigma^2)T+\sigma W_T\bigr)$, donc
  \[\pd{S_T}{S_0}=\exp\bigl((r-q-\tfrac12\sigma^2)T+\sigma W_T\bigr)=\frac{S_T}{S_0}.
  \]
  Donc une estimation MC pathwise est
  \[\widehat{\Delta}_{\mathrm{pw}}=\frac{1}{M}\sum_{m=1}^M \ee^{-rT}g'(S_T^{(m)})\,\frac{S_T^{(m)}}{S_0}.
  \]
  \item La méthode \emph{likelihood ratio} écrit $V_0=\int g(x)f(x;S_0)\,\dd x$ et dérive $f$ :
  \[\pd{V_0}{S_0}=\int g(x)\,\pd{f}{S_0}(x;S_0)\,\dd x = \E\left[\ee^{-rT}g(S_T)\,\pd{}{S_0}\log f(S_T;S_0)\right].\]
  Elle ne demande pas $g'$ et marche pour des payoffs non lisses, mais peut avoir plus de variance.
\end{enumerate}

\chapter{Feuille 5 : risk management (VaR, ES/CVaR, backtesting)}
\label{chap:td5}

Cette feuille complète le cours en rendant \emph{calculables} des notions de risque utilisées en pratique :
quantile (VaR), queue moyenne (ES/CVaR), et validation par backtesting.
Les objectifs pédagogiques :
\begin{itemize}
  \item savoir \textbf{définir} correctement VaR et ES (ce n'est pas la même chose) ;
  \item comprendre les \textbf{hypothèses} implicites (stationnarité, normalité...) ;
  \item savoir \textbf{vérifier} une méthode (exceedances et tests simples).
\end{itemize}

\section{Exercice 1 --- VaR et ES historiques sur un échantillon}
\begin{example}[Énoncé]
On observe 20 rendements journaliers (en \%) :
\[
\{-1.5,\ -0.7,\ 0.2,\ 0.1,\ 0.4,\ -2.1,\ 0.3,\ 1.0,\ -0.3,\ 0.6,\ -0.9,\ 0.2,\ 0.8,\ -1.1,\ 0.5,\ -0.4,\ 0.1,\ 0.2,\ -0.6,\ 0.7\}.
\]
On considère la perte $L=-R$ (donc une perte est positive).
\begin{enumerate}[label=\alph*)]
  \item Calculer la VaR historique à 95\% (sur cet échantillon discret).
  \item Calculer l'ES (Expected Shortfall) historique à 95\%.
  \item Expliquer pourquoi l'ES est plus sensible aux queues que la VaR.
\end{enumerate}
\end{example}

\subsection*{Solution}
\begin{enumerate}[label=\alph*)]
  \item Trier les pertes (c.-à-d. les $-R$). Les plus grandes pertes correspondent aux rendements les plus négatifs : $-2.1$, $-1.5$, $-1.1$, ...
  Avec $n=20$, le quantile 95\% correspond à environ le rang $\lceil 0.95\,n\rceil=19$ pour les pertes (convention discrète).
  Ici, une façon simple : la VaR 95\% est la \textbf{deuxième} plus grande perte (car 5\% de 20 = 1 observation dans la queue).
  Donc $\VaR_{0.95}\approx 1.5\%$ ou $2.1\%$ selon la convention d'interpolation ; l'important est d'énoncer la règle choisie.
  \item L'ES 95\% est la moyenne des pertes au-delà de la VaR.
  Sur 20 observations, la queue 5\% contient environ 1 observation : l'ES devient alors (presque) la pire perte, soit $\approx 2.1\%$.
  Si l'on prend 2 observations de queue pour stabiliser (méthode pratique), l'ES est la moyenne de $2.1\%$ et $1.5\%$, soit $1.8\%$.
  \item La VaR ne regarde que le \emph{seuil} (un quantile) : elle ignore la gravité au-delà.
  L'ES moyenne la queue : elle ``voit'' si les scénarios extrêmes deviennent plus sévères.
\end{enumerate}

\section{Exercice 2 --- VaR paramétrique sous normalité}
\begin{example}[Énoncé]
On suppose un P\&L journalier $X\sim\Normal(\mu,\sigma^2)$.
\begin{enumerate}[label=\alph*)]
  \item Exprimer la VaR à niveau $\alpha$ d'une perte $L=-X$.
  \item Pour $\mu=0$ et $\sigma=1\%$, donner la VaR 99\%.
  \item Quelles limites de cette hypothèse pour des rendements actions ?
\end{enumerate}
\end{example}

\subsection*{Solution}
\begin{enumerate}[label=\alph*)]
  \item $L=-X$ est normal de moyenne $-\mu$ et variance $\sigma^2$.
  La VaR à niveau $\alpha$ est le quantile :
  \[\VaR_{\alpha}=F_L^{-1}(\alpha)=-\mu+\sigma\,z_{\alpha},\]
  où $z_{\alpha}$ est le quantile de la loi normale standard.
  \item $z_{0.99}\approx 2.33$, donc $\VaR_{0.99}\approx 2.33\%$.
  \item Les rendements ont souvent des queues plus épaisses et de l'asymétrie ; la normalité sous-estime les pertes extrêmes.
  De plus, $\sigma$ varie dans le temps (volatilité conditionnelle), ce que la VaR normale stationnaire ne capture pas.
\end{enumerate}

\section{Exercice 3 --- backtesting : test de Kupiec (couverture inconditionnelle)}
\begin{example}[Énoncé]
On produit chaque jour une VaR 99\% sur un portefeuille.
Sur $N=250$ jours, on observe $x=6$ dépassements (jours où la perte réalise $>\VaR$).
\begin{enumerate}[label=\alph*)]
  \item Sous le modèle correct, quelle est la loi de $x$ ?
  \item Écrire le taux empirique $\hat{p}=x/N$ et le comparer à $p=1-0.99=0.01$.
  \item Donner l'intuition du test : que signifie ``trop de dépassements'' ?
\end{enumerate}
\end{example}

\subsection*{Solution}
\begin{enumerate}[label=\alph*)]
  \item Si chaque jour la probabilité de dépassement vaut $p=0.01$ et est indépendante, alors $x\sim\mathrm{Bin}(N,p)$.
  \item $\hat{p}=6/250=2.4\%$, supérieur à 1\% : cela suggère une VaR trop optimiste.
  \item Trop de dépassements signifie que la VaR est sous-calibrée (elle sous-estime les queues), ou que les hypothèses (indépendance, stationnarité) sont fausses.
\end{enumerate}

\begin{jurybox}
\begin{itemize}
  \item VaR = quantile : seuil de perte ; ES/CVaR = moyenne de la queue : sévérité au-delà du seuil.
  \item Historique : simple, peu d'hypothèses, mais instable en petite taille et aveugle aux régimes futurs.
  \item Paramétrique normale : rapide, mais risque de sous-estimer les queues et d'ignorer la vol variable.
  \item Backtesting : compter les dépassements et comprendre ce qu'ils impliquent sur la calibration du risque.
\end{itemize}
\end{jurybox}

\appendix
\part{Annexes}

\chapter{Méthodes de pricing : arbre binomial, Monte Carlo, PDE}
\label{app:pricing_methods}

Cette annexe synthétise des méthodes de pricing complémentaires.
L'objectif n'est pas de remplacer des cours de numérique, mais de donner une vue claire des principes,
des formules de base, et des pièges.

\section{Arbre binomial (Cox--Ross--Rubinstein)}
\label{sec:crr}

\subsection{Discrétisation du sous-jacent}
On discrétise $[0,T]$ en $N$ pas de taille $\Delta t = T/N$.
À chaque pas, le prix monte ou baisse :
\[
S_{k+1} =
\begin{cases}
u S_k & \text{avec probabilité } p,\\
d S_k & \text{avec probabilité } 1-p,
\end{cases}
\]
avec typiquement
\[
u=\ee^{\sigma\sqrt{\Delta t}},\qquad d=\ee^{-\sigma\sqrt{\Delta t}}=\frac{1}{u}.
\]

\subsection{Probabilité risque-neutre}
L'absence d'arbitrage impose que l'espérance risque-neutre du sous-jacent croisse à $(r-q)$ :
\[
\E^\Q[S_{k+1}\mid S_k] = \ee^{(r-q)\Delta t}S_k.
\]
Donc
\[
p u + (1-p)d = \ee^{(r-q)\Delta t}
\quad\Longrightarrow\quad
p = \frac{\ee^{(r-q)\Delta t}-d}{u-d}.
\]
Condition de no-arbitrage : $d < \ee^{(r-q)\Delta t} < u$ (sinon $p\notin(0,1)$).

\subsection{Pricing européen}
On calcule les payoffs au temps $T$ sur les noeuds de l'arbre,
puis on remonte par récurrence :
\[
V_k(S)=\ee^{-r\Delta t}\Bigl(p V_{k+1}(uS) + (1-p)V_{k+1}(dS)\Bigr).
\]
À $k=0$, on obtient $V_0$.

\subsection{Options américaines}
Pour une option américaine, on compare à chaque noeud :
\[
V_k(S)=\max\Bigl(g(S),\ \ee^{-r\Delta t}(p V_{k+1}(uS)+(1-p)V_{k+1}(dS))\Bigr).
\]
C'est une condition d'exercice anticipé.

\begin{example}[Pourquoi l'américaine est plus difficile ?]
Pour un call sur action sans dividende ($q=0$), il n'est jamais optimal d'exercer tôt :
l'américaine = l'européenne.
Pour un put, exercer tôt peut être optimal : l'option américaine vaut plus.
Le binomial capture cet effet via le $\max$ noeud par noeud.
\end{example}

\section{Monte Carlo (pricing comme espérance)}
\label{sec:mc}

\subsection{Principe}
Sous $\Q$ :
\[
V_0 = \E^\Q\left[\ee^{-rT} g(S_T)\right].
\]
Monte Carlo approxime cette espérance par une moyenne empirique :
\[
\widehat{V}_0 = \frac{1}{M}\sum_{m=1}^M \ee^{-rT} g(S_T^{(m)}),
\]
où $S_T^{(m)}$ sont des scénarios simulés.

\subsection{Erreur statistique}
Par le TCL :
\[
\sqrt{M}\left(\widehat{V}_0 - V_0\right)\Rightarrow \Normal(0,\sigma_g^2),
\]
donc l'erreur décroît comme $1/\sqrt{M}$.
C'est lent : doubler la précision demande $\times 4$ scénarios.

\subsection{Réduction de variance (outils classiques)}
\begin{itemize}
  \item \textbf{antithétiques} : utiliser $Z$ et $-Z$.
  \item \textbf{control variate} : utiliser une variable $H$ de prix connu, corrélée au payoff.
  \item \textbf{stratification} : mieux répartir les tirages.
\end{itemize}

\begin{example}[Control variate avec Black--Scholes]
Si on price un payoff compliqué mais corrélé à un call vanilla,
on peut utiliser le call BS comme contrôle : cela réduit fortement la variance.
\end{example}

\subsection{American options : régression (LSMC)}
Monte Carlo direct ne gère pas l'exercice anticipé.
Longstaff--Schwartz approxime la valeur de continuation par régression :
on simule des trajectoires, puis on estime $\E[V_{t+1}|S_t]$ via une base de fonctions (polynômes).
C'est un outil standard.

\section{PDE et différences finies}
\label{sec:fd}

Pour une diffusion, le prix $V(t,S)$ satisfait une PDE (ex. Black--Scholes).
On peut résoudre numériquement via différences finies :
\begin{itemize}
  \item schéma explicite (simple mais condition de stabilité) ;
  \item implicite (stable mais résout un système linéaire) ;
  \item Crank--Nicolson (ordre 2 en temps, compromis).
\end{itemize}

\subsection{Schéma explicite (idée)}
On discrétise $t$ et $S$ et on approxime les dérivées par des différences.
Le schéma explicite avance $V$ dans le temps via une formule directe.
Il est conditionnellement stable (condition CFL).

\subsection{American options : condition de complémentarité}
L'américaine impose
\[
V(t,S)\ge g(S),\qquad
\pd{V}{t}+\mathcal{L}V-rV \le 0,\qquad
\bigl(V-g\bigr)\Bigl(\pd{V}{t}+\mathcal{L}V-rV\Bigr)=0.
\]
Numériquement, cela devient un problème de type obstacle (LCP), résolu par méthodes projetées.

\begin{jurybox}
\begin{itemize}
\item Binomial CRR : discretise $S$ en up/down ; la probabilité risque-neutre $p$ vient de l'absence d'arbitrage.
\item Options américaines : backward induction avec $V=\max(\text{exercice},\text{continuation})$.
\item Monte Carlo : prix = moyenne des payoffs actualisés ; erreur $\sim 1/\sqrt{M}$, d'où l'importance de la réduction de variance.
\item PDE/FD : alternative déterministe ; indispensable pour américaines via obstacle.
\end{itemize}
\end{jurybox}

\chapter{Payoffs : options vanilles et stratégies (graphes)}
\label{app:payoffs}

Cette annexe sert de ``mémoire visuelle''.
On représente le payoff à maturité en fonction de $S_T$.
Les axes sont volontairement simplifiés (pas de prime payée) : on veut voir la structure.

\section{Call et put}
\begin{figure}[h!]
\centering
\begin{tikzpicture}
\begin{axis}[
  width=0.82\textwidth,
  height=0.45\textwidth,
  grid=both,
  xlabel={$S_T$},
  ylabel={Payoff},
  xmin=0, xmax=200,
  ymin=-10, ymax=120
]
\addplot+[
  domain=0:200,
  samples=400
]
{max(x-100,0)};
\addlegendentry{Call (K=100)}
\end{axis}
\end{tikzpicture}
\caption{Payoff d'un call européen.}
\end{figure}

\begin{figure}[h!]
\centering
\begin{tikzpicture}
\begin{axis}[
  width=0.82\textwidth,
  height=0.45\textwidth,
  grid=both,
  xlabel={$S_T$},
  ylabel={Payoff},
  xmin=0, xmax=200,
  ymin=-10, ymax=120
]
\addplot+[
  domain=0:200,
  samples=400
]
{max(100-x,0)};
\addlegendentry{Put (K=100)}
\end{axis}
\end{tikzpicture}
\caption{Payoff d'un put européen.}
\end{figure}

\section{Straddle et strangle}
\begin{figure}[h!]
\centering
\begin{tikzpicture}
\begin{axis}[
  width=0.82\textwidth,
  height=0.45\textwidth,
  grid=both,
  xlabel={$S_T$},
  ylabel={Payoff},
  xmin=0, xmax=200,
  ymin=-10, ymax=120
]
\addplot+[
  domain=0:200,
  samples=400
]
{max(x-100,0)+max(100-x,0)};
\addlegendentry{Straddle (K=100)}
\end{axis}
\end{tikzpicture}
\caption{Straddle : long call + long put, même strike.}
\end{figure}

\begin{figure}[h!]
\centering
\begin{tikzpicture}
\begin{axis}[
  width=0.82\textwidth,
  height=0.45\textwidth,
  grid=both,
  xlabel={$S_T$},
  ylabel={Payoff},
  xmin=0, xmax=200,
  ymin=-10, ymax=120
]
\addplot+[
  domain=0:200,
  samples=400
]
{max(x-110,0)+max(90-x,0)};
\addlegendentry{Strangle (Kp=90,Kc=110)}
\end{axis}
\end{tikzpicture}
\caption{Strangle : long put OTM + long call OTM.}
\end{figure}

\section{Spreads et structures à risque borné}
\begin{figure}[h!]
\centering
\begin{tikzpicture}
\begin{axis}[
  width=0.82\textwidth,
  height=0.45\textwidth,
  grid=both,
  xlabel={$S_T$},
  ylabel={Payoff},
  xmin=0, xmax=200,
  ymin=-10, ymax=80
]
\addplot+[
  domain=0:200,
  samples=400
]
{max(x-90,0)-max(x-110,0)};
\addlegendentry{Bull call spread (90/110)}
\end{axis}
\end{tikzpicture}
\caption{Bull call spread : long call strike bas, short call strike haut (gain borné).}
\end{figure}

\begin{figure}[h!]
\centering
\begin{tikzpicture}
\begin{axis}[
  width=0.82\textwidth,
  height=0.45\textwidth,
  grid=both,
  xlabel={$S_T$},
  ylabel={Payoff},
  xmin=0, xmax=200,
  ymin=-10, ymax=60
]
\addplot+[
  domain=0:200,
  samples=400
]
{max(x-90,0)-2*max(x-100,0)+max(x-110,0)};
\addlegendentry{Butterfly (90/100/110)}
\end{axis}
\end{tikzpicture}
\caption{Butterfly (calls) : payoff positif autour du strike central, borné.}
\end{figure}

\section{Table mémo}
\begin{longtable}{@{}p{3.4cm}p{4.2cm}p{6.8cm}@{}}
\toprule
Stratégie & Composition & Intuition \\
\midrule
Call & $(S_T-K)^+$ & Exposition haussière convex. \\
Put & $(K-S_T)^+$ & Protection baisse. \\
Straddle & Call + Put (même $K$) & Pari sur forte volatilité (mouvement dans un sens ou l'autre). \\
Strangle & Put OTM + Call OTM & Volatilité, mais plus ``large'' (coût souvent plus faible). \\
Call spread & Long call bas, short call haut & Exposition haussière à gain borné. \\
Butterfly & Long/short/long calls & Cible une zone de prix (pari sur faible vol autour du centre). \\
\bottomrule
\end{longtable}

\begin{jurybox}
\begin{itemize}
\item Visualiser les payoffs est indispensable : c'est la manière la plus rapide de vérifier une intuition.
\item Straddle/strangle sont des paris directionnellement neutres mais volatils.
\item Les spreads et butterflies bornent le risque (et le gain) : ce sont des outils de structuration.
\end{itemize}
\end{jurybox}

\chapter{Glossaire \& carte des notions}
\label{app:glossary}

\section{Glossaire (définitions ultra-courtes)}
\begin{longtable}{@{}p{3.2cm}p{11cm}@{}}
\toprule
Terme & Définition (1 ligne) \\
\midrule
Arbitrage & Stratégie à profit certain sans risque et sans investissement net. \\
Mesure risque-neutre $\Q$ & Probabilité sous laquelle les prix actualisés sont des martingales. \\
Brownien & Bruit continu gaussien à incréments indépendants. \\
Itô & Calcul différentiel adapté aux semi-martingales ; règle $(\dd W)^2=\dd t$. \\
Greeks & Dérivées du prix d'option (sensibilités). \\
Volatilité implicite & $\sigma$ BS qui reproduit un prix de marché. \\
Smile/Skew & Variation de l'IV selon le strike ; asymétrie. \\
Heston & Modèle à variance stochastique CIR corrélée au spot. \\
SABR & Modèle stochastique sur forward avec élasticité $\beta$. \\
Jump diffusion & Diffusion + sauts (Poisson composé). \\
Carr--Madan & Pricing via transformée de Fourier amortie et FFT. \\
Hedging & Construction d'un portefeuille pour réduire le risque d'un payoff. \\
PnL & Profit and Loss : résultat financier d'une stratégie. \\
MDP & Modèle de décision séquentielle (états, actions, rewards). \\
SARSA / Q-learning & Algorithmes RL basés sur Bellman/TD. \\
DQN & Q-learning où $Q$ est approximé par un réseau de neurones. \\
Discount factor & $P(0,T)$ : valeur actuelle d'1 euro à $T$. \\
Forward & Taux ou prix impliqué par les discount factors (absence d'arbitrage). \\
Nelson--Siegel & Paramétrisation level/slope/curvature d'une courbe de taux. \\
Markowitz & Optimisation moyenne--variance via $\mu$ et $\Sigma$. \\
Risk parity & Allocation égalisant les contributions au risque. \\
PCA & Diagonalisation de $\Sigma$ pour extraire des facteurs principaux. \\
\bottomrule
\end{longtable}

\section{Carte rapide : où chaque brique est utilisée}
\begin{longtable}{@{}p{4.2cm}p{9.8cm}@{}}
\toprule
Brique & Où elle intervient \\
\midrule
Brownien / Itô & Base de Black--Scholes, Heston, diffusion de taux ; simulation MC. \\
Mesure risque-neutre & Formule d'espérance des prix, pricing sans $\mu$, calibration cohérente. \\
Greeks & Couverture (delta/gamma/vega), sensitivités de calibration et risk management. \\
Vol implicite / surface & Calibration : comparaison modèle vs marché ; arbitrage en strike/maturité. \\
Heston / sauts / SABR & Modèles pour reproduire smile/skew ; pricing via Fourier ; calibration. \\
Rough vol & Smile court terme et rugosité de la volatilité ; simulation plus difficile. \\
RL hedging & Optimisation du PnL discret sous coûts/risque ; approche MDP. \\
Courbe des taux & Discounting et forwards : base de tous les prix en taux/FX. \\
Markowitz / RP / PCA & Décision allocation : construire des portefeuilles robustes et comprendre les risques. \\
\bottomrule
\end{longtable}

\begin{jurybox}
\begin{itemize}
\item Le glossaire doit être connu : c'est le vocabulaire minimal pour parler finance quantitative proprement.
\item La carte des notions montre la cohérence : chaque brique sert ailleurs (pas de chapitre ``isolé'').
\end{itemize}
\end{jurybox}

\chapter*{Bibliographie indicative}
\addcontentsline{toc}{chapter}{Bibliographie indicative}

\begin{thebibliography}{99}

\bibitem{shreve}
S.~E. Shreve,
\emph{Stochastic Calculus for Finance I \& II},
Springer.

\bibitem{bjork}
T.~Bj{\"o}rk,
\emph{Arbitrage Theory in Continuous Time},
Oxford University Press.

\bibitem{gatheral}
J.~Gatheral,
\emph{The Volatility Surface: A Practitioner's Guide},
Wiley.

\bibitem{wilmott}
P.~Wilmott,
\emph{Paul Wilmott Introduces Quantitative Finance},
Wiley.

\bibitem{glasserman}
P.~Glasserman,
\emph{Monte Carlo Methods in Financial Engineering},
Springer.

\bibitem{heston}
S.~Heston,
``A Closed-Form Solution for Options with Stochastic Volatility with Applications to Bond and Currency Options'',
\emph{Review of Financial Studies}, 1993.

\bibitem{merton}
R.~Merton,
``Option Pricing When Underlying Stock Returns Are Discontinuous'',
\emph{Journal of Financial Economics}, 1976.

\bibitem{sabr_hagan}
P.~Hagan, D.~Kumar, A.~Lesniewski, D.~Woodward,
``Managing Smile Risk'',
Wilmott Magazine, 2002.

\bibitem{rebonato}
R.~Rebonato,
\emph{Interest-Rate Option Models},
Wiley.

\bibitem{carr_madan}
P.~Carr, D.~Madan,
``Option Valuation Using the Fast Fourier Transform'',
\emph{Journal of Computational Finance}, 1999.

\bibitem{bernoulli_rough}
J.~Gatheral, T.~Jaisson, M.~Rosenbaum,
``Volatility is rough'',
\emph{Quantitative Finance}, 2018.

\bibitem{sutton_barto}
R.~Sutton, A.~Barto,
\emph{Reinforcement Learning: An Introduction},
MIT Press.

\end{thebibliography}

\end{document}
